\documentclass[12pt,letterpaper]{article}

% ============================================================
% PACKAGES
% ============================================================

% PACKAGES
% ============================================================
\usepackage[margin=1in]{geometry}
\usepackage{amsmath,amssymb,amsfonts}
\usepackage{listings}
\usepackage{xcolor}
\usepackage{hyperref}
\usepackage{booktabs}
\usepackage{longtable}
\usepackage{array}
\usepackage{graphicx}
\usepackage{float}
\usepackage{fancyhdr}
\usepackage{titlesec}
\usepackage{enumitem}
\usepackage{mdframed}
\usepackage{caption}
\usepackage{subcaption}
\usepackage{tikz}
\usetikzlibrary{arrows.meta, positioning, shapes.geometric, fit, backgrounds}

% ============================================================
% COLORS
% ============================================================
\definecolor{fortranblue}{RGB}{0,70,140}
\definecolor{fortrangreen}{RGB}{0,120,0}
\definecolor{fortranred}{RGB}{160,0,0}
\definecolor{fortranorange}{RGB}{200,100,0}
\definecolor{codebg}{RGB}{248,248,248}
\definecolor{framecolor}{RGB}{180,180,180}
\definecolor{commentcolor}{RGB}{80,120,80}
\definecolor{keywordcolor}{RGB}{0,0,180}
\definecolor{stringcolor}{RGB}{160,0,0}
\definecolor{highlightbg}{RGB}{255,255,200}
\definecolor{sectioncolor}{RGB}{0,70,140}
\definecolor{boxbg}{RGB}{240,248,255}
\definecolor{boxframe}{RGB}{100,150,200}
\definecolor{warnbg}{RGB}{255,248,230}
\definecolor{warnframe}{RGB}{200,150,50}

% ============================================================
% LISTINGS STYLE
% ============================================================
\lstdefinelanguage{Fortran90GGL90}{
  language=[90]Fortran,
  morekeywords={SUBROUTINE,END,IF,THEN,ELSE,ENDIF,DO,ENDDO,CALL,
                RETURN,IMPLICIT,NONE,INTEGER,REAL,LOGICAL,CHARACTER,
                PARAMETER,COMMON,INCLUDE,IFDEF,IFNDEF,ENDIF,DEFINE,
                UNDEF,MIN,MAX,SQRT,ABS,SIGN,EXP,LOG,INT,FLOAT},
  keywordstyle=\color{keywordcolor}\bfseries,
  commentstyle=\color{commentcolor}\itshape,
  stringstyle=\color{stringcolor},
  basicstyle=\ttfamily\footnotesize,
  backgroundcolor=\color{codebg},
  frame=single,
  framerule=0.8pt,
  rulecolor=\color{framecolor},
  numbers=left,
  numberstyle=\tiny\color{gray},
  numbersep=8pt,
  breaklines=true,
  breakatwhitespace=false,
  tabsize=2,
  showspaces=false,
  showstringspaces=false,
  captionpos=b,
  xleftmargin=15pt,
  xrightmargin=5pt,
  aboveskip=8pt,
  belowskip=8pt,
  sensitive=false,
  morecomment=[l]{!},
  morecomment=[l]{c},
  morecomment=[l]{C},
  morecomment=[s]{/*}{*/},
  morestring=[b]",
  morestring=[b]',
  escapeinside={(*@}{@*)}
}

\lstset{
  language=Fortran90GGL90,
  style={}
}

% ============================================================
% CUSTOM ENVIRONMENTS
% ============================================================
\newmdenv[
  backgroundcolor=boxbg,
  linecolor=boxframe,
  linewidth=1.5pt,
  roundcorner=4pt,
  innerleftmargin=10pt,
  innerrightmargin=10pt,
  innertopmargin=8pt,
  innerbottommargin=8pt
]{infobox}

\newmdenv[
  backgroundcolor=warnbg,
  linecolor=warnframe,
  linewidth=1.5pt,
  roundcorner=4pt,
  innerleftmargin=10pt,
  innerrightmargin=10pt,
  innertopmargin=8pt,
  innerbottommargin=8pt
]{warnbox}

\newmdenv[
  backgroundcolor=codebg,
  linecolor=framecolor,
  linewidth=1pt,
  roundcorner=2pt,
  innerleftmargin=8pt,
  innerrightmargin=8pt,
  innertopmargin=6pt,
  innerbottommargin=6pt
]{codebox}

% ============================================================
% HEADER / FOOTER
% ============================================================
\pagestyle{fancy}
\fancyhf{}
\fancyhead[L]{\textcolor{sectioncolor}{\textbf{MITgcm GGL90 Package}}}
\fancyhead[R]{\textcolor{gray}{\textit{Technical Reference}}}
\fancyfoot[C]{\thepage}
\renewcommand{\headrulewidth}{1pt}
\renewcommand{\headrule}{\color{sectioncolor}\hrule width\headwidth height\headrulewidth}

% ============================================================
% SECTION FORMATTING
% ============================================================
\titleformat{\section}
  {\color{sectioncolor}\Large\bfseries}
  {\thesection}{1em}{}[\color{sectioncolor}\titlerule]

\titleformat{\subsection}
  {\color{fortranblue}\large\bfseries}
  {\thesubsection}{1em}{}

\titleformat{\subsubsection}
  {\color{fortranblue}\normalsize\bfseries}
  {\thesubsubsection}{1em}{}

% ============================================================
% HYPERREF SETUP
% ============================================================
\hypersetup{
  colorlinks=true,
  linkcolor=fortranblue,
  citecolor=fortrangreen,
  urlcolor=fortranred,
  pdftitle={MITgcm GGL90 Package Technical Reference},
  pdfauthor={Auto-generated from source analysis},
  pdfsubject={Ocean vertical mixing parameterization}
}

% ============================================================
% MATH OPERATORS AND SHORTCUTS
% ============================================================
\DeclareMathOperator{\sgn}{sgn}
\newcommand{\TKE}{\mathrm{TKE}}
\newcommand{\KM}{\kappa_m}
\newcommand{\KH}{\kappa_h}
\newcommand{\KE}{\kappa_E}
\newcommand{\Lmix}{L_{\mathrm{mix}}}
\newcommand{\ustar}{u_*}
\newcommand{\Nsq}{N^2}
\newcommand{\Ssq}{S^2}
\newcommand{\Ri}{\mathrm{Ri}}
\newcommand{\Pr}{\mathrm{Pr}}

% ============================================================
\begin{document}
% ============================================================

% ============================================================
% TITLE PAGE
% ============================================================
\begin{titlepage}
\begin{center}
\vspace*{2cm}

{\Huge\color{sectioncolor}\bfseries MITgcm GGL90 Package}\\[0.5cm]
{\Large\color{fortranblue} Turbulent Kinetic Energy-Based Vertical Ocean Mixing}\\[0.5cm]
{\large\color{gray} Technical Reference: Architecture, Physics, and Implementation}

\vspace{1.5cm}
\rule{\linewidth}{2pt}
\vspace{0.5cm}

\begin{infobox}
\textbf{Package Location:} \texttt{MITgcm/pkg/ggl90/}\\[4pt]
\textbf{Primary Reference:} Gaspar, P., Y.\ Gregoris, and J.-M.\ Lefevre (1990).
\textit{A simple eddy kinetic energy model for simulations of the oceanic vertical mixing}.
Journal of Geophysical Research, 95(C9), 16,179--16,193.\\[4pt]
\textbf{Implementation Reference:} Blanke, B., and P.\ Delecluse (1993).
\textit{Variability of the Tropical Atlantic Ocean Simulated by a General Circulation Model
with Two Different Mixed-Layer Physics}. Journal of Physical Oceanography, 23, 1363--1388.
\end{infobox}

\vspace{1cm}
\rule{\linewidth}{1pt}
\vspace{0.5cm}

{\large\textbf{Compiled from source files:}}\\[6pt]
\begin{tabular}{ll}
\texttt{GGL90.h}                      & State variable declarations \\
\texttt{GGL90\_OPTIONS.h}              & Compile-time options \\
\texttt{ggl90\_calc.F}                 & Main driver routine (1,177 lines) \\
\texttt{ggl90\_mixinglength.F}         & Mixing length computation (421 lines) \\
\texttt{ggl90\_calc\_diff.F}           & Diffusivity interface (74 lines) \\
\texttt{ggl90\_calc\_visc.F}           & Viscosity interface (64 lines) \\
\texttt{ggl90\_readparms.F}            & Parameter file reader (451 lines) \\
\texttt{ggl90\_check.F}                & Consistency checks (166 lines) \\
\texttt{ggl90\_init\_fixed.F}          & Time-invariant initialization (67 lines) \\
\texttt{ggl90\_init\_varia.F}          & Time-varying initialization (156 lines) \\
\texttt{ggl90\_diagnostics\_init.F}    & Diagnostics registration (202 lines) \\
\texttt{ggl90\_output.F}               & Diagnostics and I/O (98 lines) \\
\texttt{ggl90\_read\_pickup.F}         & Restart file reader (69 lines) \\
\texttt{ggl90\_write\_pickup.F}        & Restart file writer (60 lines) \\
\texttt{ggl90\_exchanges.F}            & Halo exchanges (48 lines) \\
\texttt{ggl90\_idemix.F}               & IDEMIX internal wave model (598 lines) \\
\texttt{ggl90\_add\_stokesdrift.F}     & Langmuir circulation (79 lines) \\
\end{tabular}

\vspace{0.5cm}
{\small\color{gray} Total package size: $\sim$3,958 lines of Fortran code}

\vfill
{\color{gray}\today}
\end{center}
\end{titlepage}

% ============================================================
\tableofcontents
\newpage
% ============================================================

% ============================================================
\section{Introduction and Scientific Background}
% ============================================================

The GGL90 package implements a one-equation turbulence closure scheme based on
a prognostic equation for turbulent kinetic energy (TKE). Originally proposed by
Gaspar, Gregoris, and Lefevre (1990, hereafter GGL90), this parameterization
provides a physically based representation of vertical mixing throughout the
ocean water column. The MITgcm implementation closely follows the Blanke and
Delecluse (1993) formulation, which was developed for the OPA ocean model and
has been extensively validated in global ocean simulations.

Unlike diagnostic mixing schemes (such as KPP), GGL90 carries TKE as a
prognostic variable, explicitly simulating its production, dissipation, and
transport. This approach captures the time-dependent evolution of turbulence,
making it particularly well-suited for scenarios with rapidly varying forcing
or when turbulence memory effects are important.

\subsection{Physical Motivation}

Vertical mixing in the ocean arises from turbulent motions that span a wide
range of spatial and temporal scales. While direct numerical simulation of
these turbulent flows is computationally infeasible in global ocean models,
turbulence closure schemes provide a practical framework for parameterizing
the effects of unresolved turbulent motions on the resolved-scale flow.

GGL90 adopts a \textbf{one-equation turbulence closure} approach, wherein:

\begin{enumerate}[leftmargin=*, label=\textbf{\arabic*.}]
  \item \textbf{Turbulent kinetic energy (TKE)} is treated as a prognostic
        variable satisfying a budget equation that balances production,
        buoyancy work, dissipation, and diffusion.

  \item \textbf{Mixing length ($\Lmix$)} is diagnosed from the vertical
        structure of TKE and stratification, subject to geometric constraints
        (distance to boundaries).

  \item \textbf{Eddy viscosity ($\KM$) and diffusivity ($\KH$)} are computed
        from TKE and mixing length using dimensional analysis:
        \[
        \KM \sim \Lmix \sqrt{\TKE}
        \]

  \item \textbf{Stability functions} relate the eddy diffusivity to eddy
        viscosity through an empirical Prandtl number that depends on local
        stratification.
\end{enumerate}

This framework is grounded in the classical turbulence theory of Kolmogorov
(1942) and extends it to stratified, sheared oceanic flows where buoyancy
effects play a critical role in modulating turbulence.

\subsubsection{Key Physical Processes}

The GGL90 scheme captures several fundamental mechanisms of oceanic turbulence:

\begin{itemize}[leftmargin=*]
  \item \textbf{Shear production} -- Vertical shear in the horizontal velocity
        field ($S^2 = |\partial \mathbf{u}/\partial z|^2$) generates TKE
        through the extraction of energy from the mean flow. This mechanism
        dominates in regions with strong currents or wind-driven surface layers.

  \item \textbf{Buoyancy work} -- Stratification ($N^2 = -g/\rho_0 \,
        \partial \rho / \partial z$) either suppresses turbulence (stable
        stratification, $N^2 > 0$) or enhances it (convective instability,
        $N^2 < 0$). The buoyancy term in the TKE budget can act as either
        a sink or source depending on the sign of $N^2$ and the direction
        of turbulent buoyancy flux.

  \item \textbf{Dissipation} -- Turbulent kinetic energy is irreversibly
        converted to heat at the smallest (Kolmogorov) scales. GGL90
        parameterizes this dissipation rate as $\varepsilon \sim
        \TKE^{3/2} / \Lmix$, consistent with dimensional analysis.

  \item \textbf{Vertical transport} -- TKE can be transported vertically
        by turbulent diffusion and pressure work. GGL90 includes a diffusion
        term for TKE with diffusivity $\KE \sim \KM$.

  \item \textbf{Surface and bottom forcing} -- Wind stress at the ocean
        surface injects TKE into the boundary layer. Bottom drag can also
        generate TKE near the seafloor. These boundary sources are critical
        for maintaining realistic mixing levels.
\end{itemize}

\subsection{GGL90 vs.\ KPP: A Comparison}

Both GGL90 and KPP are widely used in ocean general circulation models, but
they differ fundamentally in their approach to parameterizing vertical mixing.
Table~\ref{tab:ggl90_vs_kpp} summarizes the key distinctions.

\begin{table}[h!]
\centering
\caption{Comparison of GGL90 and KPP vertical mixing schemes.}
\label{tab:ggl90_vs_kpp}
\begin{tabular}{@{}lll@{}}
\toprule
\textbf{Aspect} & \textbf{GGL90} & \textbf{KPP} \\
\midrule
\textbf{Type} & Prognostic (TKE-based) & Diagnostic (profile-based) \\
\textbf{Primary variable} & TKE & Boundary layer depth ($h_{\mathrm{bl}}$) \\
\textbf{Time dependence} & Explicit TKE evolution & Instantaneous diagnosis \\
\textbf{Mixing length} & Geometric + TKE-based & Geometric (distance to $h_{\mathrm{bl}}$) \\
\textbf{Boundary layer} & Emergent from TKE & Explicitly diagnosed \\
\textbf{Memory effects} & Yes (TKE persists) & No (local diagnosis) \\
\textbf{Computational cost} & Higher (extra PDE) & Lower (diagnostic) \\
\textbf{Calibration} & Few parameters & Many profile functions \\
\textbf{Deep convection} & Natural (TKE grows) & Requires special treatment \\
\textbf{Adjoint complexity} & Moderate (Appendix~\ref{app:adjoint_considerations}) & Higher (many branches) \\
\bottomrule
\end{tabular}
\end{table}

\subsubsection{When to Use GGL90}

GGL90 is particularly advantageous in the following scenarios:

\begin{itemize}[leftmargin=*]
  \item \textbf{Rapidly varying forcing} -- Diurnal or synoptic-scale wind
        variations, where TKE memory affects the response time of mixing.

  \item \textbf{Deep convection} -- High-latitude regions where convective
        plumes penetrate to great depths. GGL90 naturally allows TKE to grow
        and propagate downward without requiring special convective adjustment.

  \item \textbf{Adjoint-based optimization} -- GGL90's simpler parameter space
        facilitates gradient-based data assimilation (see
        Appendix~\ref{app:adjoint_considerations} for detailed discussion).

  \item \textbf{Coupled models} -- When the ocean model is coupled to sea ice
        or atmosphere models with rapid timescale interactions, the prognostic
        nature of TKE can improve synchronization.
\end{itemize}

\subsubsection{When to Use KPP}

KPP remains preferred in many applications due to:

\begin{itemize}[leftmargin=*]
  \item \textbf{Lower computational cost} -- No additional PDE to solve.

  \item \textbf{Well-established calibration} -- Decades of tuning in numerous
        ocean models (e.g., CESM, ACCESS, GFDL models).

  \item \textbf{Explicit non-local flux} -- KPP includes a counter-gradient
        term for boundary layer mixing, which GGL90 does not.

  \item \textbf{Smoother mixed layer base} -- KPP's explicit $h_{\mathrm{bl}}$
        can produce sharper transitions than GGL90's gradual TKE decay.
\end{itemize}

In practice, both schemes can produce realistic simulations when properly
tuned. The choice often depends on the specific application, available
computational resources, and user familiarity.

\subsection{Key Equations Overview}

This subsection provides a high-level summary of the governing equations in
GGL90. Detailed derivations and code implementations are presented in
subsequent sections.

\subsubsection{TKE Evolution Equation}

The prognostic equation for TKE (denoted $e = \TKE$ in the code) is:
\begin{equation}
\frac{\partial \TKE}{\partial t} = P + B - \varepsilon + D
\label{eq:tke_budget}
\end{equation}
where:
\begin{align*}
P &= \KM \, S^2 \quad &&\text{(shear production)} \\
B &= -\KH \, N^2 \quad &&\text{(buoyancy term)} \\
\varepsilon &= c_\varepsilon \frac{\TKE^{3/2}}{\Lmix} \quad &&\text{(dissipation)} \\
D &= \frac{\partial}{\partial z} \left( \KE \frac{\partial \TKE}{\partial z} \right) \quad &&\text{(vertical diffusion)}
\end{align*}

The dissipation constant $c_\varepsilon$ (denoted \texttt{GGL90ceps} in the
code) is typically set to $c_\varepsilon \approx 0.7$, following empirical
constraints from laboratory and atmospheric boundary layer studies.

\subsubsection{Eddy Viscosity and Diffusivity}

Eddy coefficients are diagnosed from TKE and mixing length:
\begin{align}
\KM &= c_k \, \Lmix \sqrt{\TKE} \label{eq:kappa_m} \\
\KH &= \frac{\KM}{\alpha} \label{eq:kappa_h}
\end{align}
where $c_k \approx 0.1$ (constant \texttt{GGL90ck}) and $\alpha$ (constant
\texttt{GGL90alpha}) is the turbulent Prandtl number. In the original Gaspar
et al.\ (1990) formulation, $\alpha = 1$. However, global ocean applications
often use larger values (see Appendix~\ref{app:eccov4_config} for ECCOv4's
$\alpha = 30$ choice) to reduce diffusivity relative to viscosity, thereby
maintaining sharp thermoclines.

\subsubsection{Mixing Length}

The mixing length is diagnosed from TKE and stratification following Gaspar
et al.\ (1990), equation (9):
\begin{equation}
\Lmix^{\mathrm{(Gaspar)}} = \frac{\sqrt{2} \, \sqrt{\TKE}}{\sqrt{N^2}}
\label{eq:l_gaspar}
\end{equation}
This expression is then limited by geometric constraints (distance to surface,
distance to bottom) using one of three methods (\texttt{mxlMaxFlag} = 0, 1, or 2).
The Blanke and Delecluse (1993) Method 2 (two-way sweep) is most commonly used
(see Appendix~\ref{app:eccov4_config}).

\subsubsection{Boundary Conditions}

\textbf{Surface ($z = 0$):}
\begin{equation}
\KE \frac{\partial \TKE}{\partial z}\bigg|_{z=0} = m_2 \, u_*^3
\label{eq:bc_surface}
\end{equation}
where $u_* = \sqrt{\tau / \rho_0}$ is the friction velocity and $m_2 \approx
3.75$ (constant \texttt{GGL90m2}) determines the fraction of wind energy flux
entering the TKE budget.

\textbf{Bottom ($z = -H$):}
\begin{equation}
\TKE\big|_{z=-H} = \max(\TKE_{\mathrm{bottom}}, \TKE_{\mathrm{min}})
\label{eq:bc_bottom}
\end{equation}
where $\TKE_{\mathrm{bottom}}$ (\texttt{GGL90TKEbottom}) represents background
TKE generated by bottom drag or internal wave breaking.

These equations form the core of the GGL90 parameterization and are solved at
every model timestep. The following sections describe their implementation in
MITgcm in detail, with extensive code references and examples.

% ============================================================
\section{Package Architecture and Call Flow}
% ============================================================

The GGL90 package follows MITgcm's modular architecture, with clearly separated
initialization, computation, and I/O phases. This section describes the overall
package structure, the sequence of subroutine calls, and how GGL90 integrates
with the main model timestepping loop.

\subsection{Package Initialization Sequence}

GGL90 initialization occurs in two phases: \textbf{fixed initialization} (once
at model startup) and \textbf{variable initialization} (which may occur multiple
times, e.g., when restarting from a pickup file).

\subsubsection{Phase 1: Fixed Initialization (\texttt{GGL90\_INIT\_FIXED})}

Called from \texttt{PACKAGES\_INIT\_FIXED} in the main model initialization.
This routine (\texttt{ggl90\_init\_fixed.F:7}) performs time-invariant setup:

\begin{lstlisting}[caption={\texttt{ggl90\_init\_fixed.F}, lines 7--44 -- Fixed initialization.}]
SUBROUTINE GGL90_INIT_FIXED( myThid )
  ! Purpose: Initialize GGL90 variables that are kept fixed during the run

#ifdef ALLOW_DIAGNOSTICS
  IF ( useDiagnostics ) THEN
    CALL GGL90_DIAGNOSTICS_INIT( myThid )  ! Register diagnostic fields
  ENDIF
#endif

#ifdef ALLOW_GGL90_SMOOTH
  ! Initialize corner-point mask for horizontal smoothing
  DO bj=myByLo(myThid),myByHi(myThid)
    DO bi=myBxLo(myThid),myBxHi(myThid)
      DO j=1-OLy,sNy+OLy
        DO i=1-OLx,sNx+OLx
          mskCor(i,j,bi,bj) = 1.0
        ENDDO
      ENDDO
      ! Special treatment for cubed-sphere corners
      IF ( useCubedSphereExchange ) THEN
        CALL FILL_CS_CORNER_TR_RL( ... )
      ENDIF
    ENDDO
  ENDDO
#endif

END SUBROUTINE
\end{lstlisting}

\textbf{Key actions:}
\begin{itemize}[leftmargin=*]
  \item Register diagnostic output fields (TKE, viscosity, diffusivity, etc.)
        via \texttt{GGL90\_DIAGNOSTICS\_INIT}.
  \item If horizontal smoothing is enabled (\texttt{ALLOW\_GGL90\_SMOOTH}),
        initialize corner-point masks for proper averaging at cubed-sphere
        edges.
\end{itemize}

\subsubsection{Phase 2: Variable Initialization (\texttt{GGL90\_INIT\_VARIA})}

Called from \texttt{PACKAGES\_INIT\_VARIABLES} after parameters are read but
before timestepping begins. This routine (\texttt{ggl90\_init\_varia.F:7})
initializes the prognostic TKE field:

\begin{lstlisting}[caption={\texttt{ggl90\_init\_varia.F}, lines 42--80 (excerpts) -- TKE initialization.}]
SUBROUTINE GGL90_INIT_VARIA( myThid )
  ! Initialize TKE field either from file or with default values

  IF ( GGL90TKEFile .NE. ' ' ) THEN
    ! Read TKE from binary file
    CALL READ_FLD_XYZ_RL( GGL90TKEFile, ' ', GGL90TKE, 0, myThid )
  ELSE
    ! Initialize with minimum TKE
    DO bj = myByLo(myThid), myByHi(myThid)
      DO bi = myBxLo(myThid), myBxHi(myThid)
        DO k = 1, Nr
          DO j = 1-OLy, sNy+OLy
            DO i = 1-OLx, sNx+OLx
              GGL90TKE(i,j,k,bi,bj) = GGL90TKEmin
            ENDDO
          ENDDO
        ENDDO
      ENDDO
    ENDDO
  ENDIF

  ! Initialize IDEMIX fields if enabled
#ifdef ALLOW_GGL90_IDEMIX
  IF ( useIDEMIX ) THEN
    CALL GGL90_IDEMIX_INIT_VARIA( myThid )
  ENDIF
#endif

END SUBROUTINE
\end{lstlisting}

\textbf{Key actions:}
\begin{itemize}[leftmargin=*]
  \item If \texttt{GGL90TKEFile} is specified in \texttt{data.ggl90}, read
        initial TKE from file. Otherwise, initialize to \texttt{GGL90TKEmin}.
  \item If IDEMIX is enabled, initialize internal wave energy fields.
  \item Exchange halo regions for parallel runs.
\end{itemize}

\subsection{Per-Timestep Execution Flow}

At each model timestep, GGL90 is called from the main timestepping loop
(\texttt{model/src/do\_oceanic\_phys.F}) as part of the vertical mixing
calculations. The top-level entry point is \texttt{GGL90\_CALC}.

Figure~\ref{fig:call_flow} illustrates the complete call sequence.

\begin{figure}[h!]
\centering
\begin{tikzpicture}[
  node distance=0.8cm and 1.5cm,
  box/.style={rectangle, draw=boxframe, fill=boxbg, thick,
              minimum width=3.5cm, minimum height=0.7cm, align=center},
  subbox/.style={rectangle, draw=fortranblue, fill=boxbg!30, thick,
                 minimum width=3cm, minimum height=0.6cm, align=center, font=\small},
  arrow/.style={-{Stealth[length=3mm]}, thick, draw=fortranblue},
]

% Main driver
\node[box, fill=fortranred!20] (calc) {\textbf{GGL90\_CALC}};

% First level subroutines
\node[subbox, below=of calc, xshift=-4cm] (idemix) {GGL90\_IDEMIX\\{\tiny(optional)}};
\node[subbox, below=of calc] (mixlen) {GGL90\_MIXINGLENGTH};
\node[subbox, below=of calc, xshift=4cm] (langmuir) {GGL90\_ADD\_STOKESDRIFT\\{\tiny(optional)}};

% Second level
\node[subbox, below=of mixlen, xshift=-2cm] (visc) {GGL90\_CALC\_VISC};
\node[subbox, below=of mixlen, xshift=2cm] (diff) {GGL90\_CALC\_DIFF};

% Third level
\node[subbox, below=of visc, xshift=1cm] (exch) {GGL90\_EXCHANGES};

% Arrows
\draw[arrow] (calc) -- (idemix);
\draw[arrow] (calc) -- (mixlen);
\draw[arrow] (calc) -- (langmuir);
\draw[arrow] (mixlen) -- (visc);
\draw[arrow] (mixlen) -- (diff);
\draw[arrow] (diff) -- (exch);

% Annotations
\node[right=0.3cm of calc, font=\tiny, text width=2.5cm, align=left]
     {Main driver:\\compute TKE,\\$\KM$, $\KH$};
\node[below=0.1cm of mixlen, font=\tiny, text width=2.5cm, align=center]
     {Compute $\Lmix$\\(3 methods)};
\node[below=0.1cm of visc, font=\tiny, text width=2cm, align=center]
     {Transfer to\\U/V points};
\node[below=0.1cm of diff, font=\tiny, text width=2cm, align=center]
     {Transfer to\\T/S points};

\end{tikzpicture}
\caption{GGL90 call flow diagram showing the hierarchy of subroutines invoked
         during each timestep. Optional features (IDEMIX, Langmuir) are marked.}
\label{fig:call_flow}
\end{figure}

\subsubsection{Main Driver: \texttt{GGL90\_CALC}}

The primary computational routine (\texttt{ggl90\_calc.F:13}) orchestrates all
GGL90 calculations. Its execution proceeds as follows:

\begin{enumerate}[leftmargin=*, label=\textbf{Step \arabic*.}]
  \item \textbf{Setup} (lines 190--260)
    \begin{itemize}[label=$\circ$]
      \item Determine coordinate system (Z-coordinates vs.\ P-coordinates)
      \item Initialize local arrays for mixing coefficients, TKE, etc.
      \item Compute interface thickness factors (\texttt{hFacI})
    \end{itemize}

  \item \textbf{Optional: IDEMIX} (lines 259--266)
    \begin{itemize}[label=$\circ$]
      \item If \texttt{useIDEMIX = .TRUE.}, call \texttt{GGL90\_IDEMIX} to
            compute internal wave energy dissipation (\texttt{IDEMIX\_gTKE})
      \item This provides an additional source term for the TKE budget
    \end{itemize}

  \item \textbf{Compute stratification and shear} (lines 300--550)
    \begin{itemize}[label=$\circ$]
      \item Buoyancy frequency squared: $N^2 = -g/\rho_0 \, \partial\rho/\partial z$
            (from input \texttt{sigmaR})
      \item Vertical shear squared: $S^2 = (\partial u/\partial z)^2 + (\partial v/\partial z)^2$
      \item Surface friction velocity: $u_*^2 = \tau/\rho_0$
    \end{itemize}

  \item \textbf{Optional: Langmuir} (lines 450--550)
    \begin{itemize}[label=$\circ$]
      \item If \texttt{useLANGMUIR = .TRUE.}, augment shear with Stokes drift
            contribution
    \end{itemize}

  \item \textbf{Compute mixing length} (lines 550--600)
    \begin{itemize}[label=$\circ$]
      \item Call \texttt{GGL90\_MIXINGLENGTH} to compute $\Lmix$ using one of
            three methods (selected by \texttt{mxlMaxFlag})
    \end{itemize}

  \item \textbf{Compute eddy coefficients} (lines 600--750)
    \begin{itemize}[label=$\circ$]
      \item $\KM = c_k \, \Lmix \sqrt{\max(\TKE, \TKE_{\mathrm{min}})}$
      \item $\KH = \KM / \alpha$
      \item Apply background floors and maximum caps
    \end{itemize}

  \item \textbf{TKE budget terms} (lines 750--900)
    \begin{itemize}[label=$\circ$]
      \item Production: $P = \KM \, S^2$
      \item Buoyancy: $B = -\KH \, N^2$
      \item Dissipation: $\varepsilon = c_\varepsilon \, \TKE^{3/2} / \Lmix$
    \end{itemize}

  \item \textbf{Build tridiagonal system} (lines 900--1000)
    \begin{itemize}[label=$\circ$]
      \item Set up implicit time-stepping matrix for TKE equation
      \item Surface BC: wind-driven TKE flux
      \item Bottom BC: minimum TKE or drag-generated TKE
    \end{itemize}

  \item \textbf{Solve for TKE} (lines 1000--1050)
    \begin{itemize}[label=$\circ$]
      \item Solve tridiagonal system using Thomas algorithm
      \item Update \texttt{GGL90TKE} array
    \end{itemize}

  \item \textbf{Optional: Horizontal smoothing} (lines 1050--1100)
    \begin{itemize}[label=$\circ$]
      \item If \texttt{ALLOW\_GGL90\_SMOOTH} is defined, apply horizontal
            averaging (OPA method) to reduce grid-scale noise
    \end{itemize}

  \item \textbf{Transfer to model arrays} (lines 1100--1150)
    \begin{itemize}[label=$\circ$]
      \item \texttt{GGL90\_CALC\_VISC}: Transfer $\KM$ to U/V-point arrays
            (\texttt{GGL90viscArU}, \texttt{GGL90viscArV})
      \item \texttt{GGL90\_CALC\_DIFF}: Transfer $\KH$ to T/S-point array
            (\texttt{GGL90diffKr})
    \end{itemize}

  \item \textbf{Diagnostics} (lines 1150--1177)
    \begin{itemize}[label=$\circ$]
      \item Fill diagnostic arrays if requested (e.g., TKE snapshots, mixing
            length, viscosity profiles)
    \end{itemize}
\end{enumerate}

This 12-step sequence is executed once per timestep for each horizontal tile
(in parallel domain decomposition).

\subsection{File Organization}

The GGL90 package consists of 23 source files totaling approximately 3,958
lines of Fortran code. Table~\ref{tab:file_summary} summarizes the purpose
of each file.

\begin{table}[h!]
\centering
\caption{Summary of GGL90 source files with line counts and descriptions.}
\label{tab:file_summary}
\small
\begin{tabular}{@{}llr@{}}
\toprule
\textbf{File} & \textbf{Description} & \textbf{Lines} \\
\midrule
\multicolumn{3}{l}{\textit{Core Computation}} \\
\texttt{ggl90\_calc.F} & Main driver, TKE budget, coefficient calc. & 1,177 \\
\texttt{ggl90\_mixinglength.F} & Mixing length (3 methods) & 421 \\
\texttt{ggl90\_calc\_diff.F} & Transfer diffusivity to T/S points & 74 \\
\texttt{ggl90\_calc\_visc.F} & Transfer viscosity to U/V points & 64 \\
\midrule
\multicolumn{3}{l}{\textit{Configuration}} \\
\texttt{GGL90.h} & Common blocks, variable declarations & 178 \\
\texttt{GGL90\_OPTIONS.h} & Compile-time CPP flags & 42 \\
\texttt{ggl90\_readparms.F} & Read parameters from data.ggl90 & 451 \\
\texttt{ggl90\_check.F} & Validate parameter consistency & 166 \\
\midrule
\multicolumn{3}{l}{\textit{Initialization}} \\
\texttt{ggl90\_init\_fixed.F} & Time-invariant initialization & 67 \\
\texttt{ggl90\_init\_varia.F} & Time-varying initialization (TKE) & 156 \\
\midrule
\multicolumn{3}{l}{\textit{I/O and Diagnostics}} \\
\texttt{ggl90\_diagnostics\_init.F} & Register diagnostic fields & 202 \\
\texttt{ggl90\_output.F} & Output state to files & 98 \\
\texttt{ggl90\_read\_pickup.F} & Read restart (pickup) files & 69 \\
\texttt{ggl90\_write\_pickup.F} & Write restart (pickup) files & 60 \\
\midrule
\multicolumn{3}{l}{\textit{Optional Features}} \\
\texttt{ggl90\_idemix.F} & IDEMIX internal wave model & 598 \\
\texttt{ggl90\_add\_stokesdrift.F} & Langmuir circulation & 79 \\
\midrule
\multicolumn{3}{l}{\textit{Utilities}} \\
\texttt{ggl90\_exchanges.F} & Halo exchanges for parallel runs & 48 \\
\texttt{ggl90\_ad\_*.h} & Adjoint model directives (4 files) & 8 \\
\bottomrule
\multicolumn{2}{l}{\textbf{Total}} & \textbf{3,958} \\
\end{tabular}
\end{table}

\subsubsection{File Interdependencies}

\begin{itemize}[leftmargin=*]
  \item \textbf{Header files} (\texttt{GGL90.h}, \texttt{GGL90\_OPTIONS.h})
        are included by all computational routines to access parameters and
        state variables.

  \item \textbf{Parameter reading} (\texttt{ggl90\_readparms.F}) must occur
        before \texttt{ggl90\_init\_varia.F} because initialization depends
        on parameter values (e.g., \texttt{GGL90TKEmin}).

  \item \textbf{Diagnostics initialization} (\texttt{ggl90\_diagnostics\_init.F})
        must occur before the first timestep to register output fields with
        MITgcm's diagnostics package.

  \item \textbf{Main driver} (\texttt{ggl90\_calc.F}) calls helper routines
        (\texttt{ggl90\_mixinglength.F}, \texttt{ggl90\_calc\_diff.F},
        \texttt{ggl90\_calc\_visc.F}) but does not depend on I/O routines.
\end{itemize}

\subsubsection{Optional Feature Compilation}

Several features are controlled by CPP flags defined in
\texttt{GGL90\_OPTIONS.h} (see Section~\ref{sec:cpp_options} for details):

\begin{itemize}[leftmargin=*]
  \item \texttt{ALLOW\_GGL90\_IDEMIX}: Enables \texttt{ggl90\_idemix.F}
        (internal wave forcing).
  \item \texttt{ALLOW\_GGL90\_LANGMUIR}: Enables
        \texttt{ggl90\_add\_stokesdrift.F} (Langmuir turbulence).
  \item \texttt{ALLOW\_GGL90\_SMOOTH}: Enables horizontal smoothing in
        \texttt{ggl90\_calc.F}.
  \item \texttt{ALLOW\_GGL90\_HORIZDIFF}: Enables horizontal diffusion of TKE.
\end{itemize}

If these flags are undefined, the corresponding code blocks are excluded at
compile time, reducing memory footprint and computational cost.

% ============================================================
\section{Parameter Initialization}
% ============================================================

GGL90 parameters are read from the namelist file \texttt{data.ggl90} during
model initialization. The reading is handled by \texttt{ggl90\_readparms.F}
(lines 6--451), which sets default values and then overwrites them with
user-specified values if present. This section documents all GGL90 parameters
with their default values, physical meanings, and typical ranges.

Parameters are organized into three namelist groups:
\begin{itemize}[leftmargin=*]
  \item \texttt{GGL90\_PARM01}: Core GGL90 parameters (lines 48--57)
  \item \texttt{GGL90\_PARM02}: IDEMIX parameters (optional, lines 63--68)
  \item \texttt{GGL90\_PARM03}: Langmuir parameters (optional, lines 75--76)
\end{itemize}

\subsection{Core Physics Parameters}

Table~\ref{tab:core_params} lists the fundamental parameters governing TKE
dynamics, eddy coefficients, and mixing length computation.

\begin{table}[h!]
\centering
\caption{Core physics parameters for GGL90 (from \texttt{ggl90\_readparms.F:108--120}).}
\label{tab:core_params}
\small
\begin{tabular}{@{}lllp{6cm}@{}}
\toprule
\textbf{Parameter} & \textbf{Default} & \textbf{Units} & \textbf{Description} \\
\midrule
\texttt{GGL90ck} & 0.1 & -- & Constant $c_k$ in eddy viscosity formula
                                $\KM = c_k \Lmix \sqrt{\TKE}$ (eq.\ 10 of Gaspar
                                et al.\ 1990). Dimensionless. \\[6pt]

\texttt{GGL90ceps} & 0.7 & -- & Dissipation constant $c_\varepsilon$ in
                                $\varepsilon = c_\varepsilon \TKE^{3/2}/\Lmix$.
                                Based on Kolmogorov (1942) theory. \\[6pt]

\texttt{GGL90alpha} & 1.0 & -- & Turbulent Prandtl number relating diffusivity
                                 to viscosity: $\KH = \KM / \alpha$. Gaspar
                                 et al.\ use $\alpha=1$. See
                                 Appendix~\ref{app:eccov4_config} for alternative
                                 values in global applications. \\[6pt]

\texttt{GGL90m2} & 3.75 & -- & Surface TKE flux constant $m_2$ in boundary
                               condition $\KE \partial\TKE/\partial z|_{z=0} =
                               m_2 u_*^3$. Value from Blanke \& Delecluse (1993). \\[6pt]

\texttt{GGL90TKEmin} & $10^{-11}$ & m$^2$/s$^2$ & Minimum TKE enforced
                                    throughout the water column. Represents
                                    background mixing from internal waves and
                                    other unresolved processes. \\[6pt]

\texttt{GGL90TKEsurfMin} & $10^{-4}$ & m$^2$/s$^2$ & Minimum TKE at surface.
                                        Higher than interior \texttt{GGL90TKEmin}
                                        to ensure adequate surface boundary layer
                                        mixing. Value from Blanke \& Delecluse. \\[6pt]

\texttt{GGL90TKEbottom} & \texttt{UNSET\_RL} & m$^2$/s$^2$ & Bottom TKE.
                                                If unset, defaults to
                                                \texttt{GGL90TKEmin}. Can be
                                                increased to represent bottom
                                                drag or rough topography. \\[6pt]

\texttt{GGL90mixingLengthMin} & $10^{-8}$ & m & Minimum mixing length. Prevents
                                               division by zero and provides
                                               regularization. Physically
                                               represents molecular scales. \\[6pt]

\texttt{GGL90viscMax} & $10^{2}$ & m$^2$/s & Maximum eddy viscosity. Caps
                                             unrealistic values in poorly
                                             resolved regions or during
                                             convective events. \\[6pt]

\texttt{GGL90diffMax} & $10^{2}$ & m$^2$/s & Maximum eddy diffusivity. Similar
                                             purpose as \texttt{GGL90viscMax}. \\

\bottomrule
\end{tabular}
\end{table}

\subsubsection{Physical Interpretation}

\begin{itemize}[leftmargin=*]
  \item \textbf{$c_k$ and $c_\varepsilon$} -- These constants arise from
        dimensional analysis and are calibrated against laboratory experiments,
        large-eddy simulations, and observations. Their default values (0.1 and
        0.7) are well-established in the turbulence closure literature.

  \item \textbf{$\alpha$ (Prandtl number)} -- In neutral boundary layers,
        momentum and heat are transported similarly ($\alpha \approx 1$).
        However, in the stratified ocean interior, $\alpha > 1$ is often
        needed to prevent excessive tracer diffusion. The choice of $\alpha$
        is application-specific and often tuned to match observed mixed layer
        depths and thermocline sharpness.

  \item \textbf{$m_2$ (surface flux fraction)} -- Not all wind energy input
        to the ocean surface is converted to TKE. Some goes into surface
        gravity waves and Langmuir circulation. $m_2 \approx 3.75$ implies
        that a few percent of wind power enters the TKE budget.

  \item \textbf{Minimum TKE values} -- Even in quiescent, stratified regions,
        some level of mixing persists due to internal wave breaking, double
        diffusion, and other sub-grid processes. \texttt{GGL90TKEmin} ensures
        that $\KM$ and $\KH$ never vanish completely, which would be unphysical
        and numerically problematic.
\end{itemize}

\subsection{Mixing Length Control Parameters}

The computation of mixing length is controlled by \texttt{mxlMaxFlag} (lines
121--123) and related flags. Table~\ref{tab:mixlen_params} summarizes these.

\begin{table}[h!]
\centering
\caption{Mixing length control parameters.}
\label{tab:mixlen_params}
\begin{tabular}{@{}llp{8cm}@{}}
\toprule
\textbf{Parameter} & \textbf{Default} & \textbf{Description} \\
\midrule
\texttt{mxlMaxFlag} & 0 & Method for limiting mixing length:\\
                          \quad 0 = Distance to boundaries (simplest)\\
                          \quad 1 = One-way downward sweep\\
                          \quad 2 = Two-way sweep (Blanke \& Delecluse 1993)\\
                          Method 2 is recommended (see
                          Appendix~\ref{app:eccov4_config}). \\[6pt]

\texttt{adMxlMaxFlag} & \texttt{UNSET\_I} & Alternative \texttt{mxlMaxFlag}
                                             for adjoint runs. If unset, uses
                                             \texttt{mxlMaxFlag}. See
                                             Appendix~\ref{app:adjoint_considerations}
                                             for adjoint-specific settings. \\[6pt]

\texttt{mxlSurfFlag} & \texttt{.FALSE.} & If \texttt{.TRUE.}, forces mixing
                                           length between surface and first
                                           level to equal the layer thickness.
                                           Ensures adequate surface boundary
                                           layer mixing. \\

\bottomrule
\end{tabular}
\end{table}

\textbf{Mixing Length Method Comparison:}

\begin{itemize}[leftmargin=*]
  \item \textbf{Method 0} -- Simplest: $\Lmix = \min(L_{\mathrm{Gaspar}},
        z_{\mathrm{surf}}, z_{\mathrm{bot}})$ where $z_{\mathrm{surf}}$ and
        $z_{\mathrm{bot}}$ are distances to surface and bottom. Fast but can
        produce discontinuities.

  \item \textbf{Method 1} -- Downward sweep smooths $\Lmix$ in the vertical:
        starting from the surface, ensures $\Lmix(k) \leq \Lmix(k-1) +
        \Delta z$. Prevents abrupt jumps but is asymmetric.

  \item \textbf{Method 2} -- Two-way sweep (downward then upward) produces the
        smoothest profiles and is most physical. Used in OPA and many global
        applications (Appendix~\ref{app:eccov4_config}). Adds minimal
        computational cost. \textbf{Recommended for production runs.}
\end{itemize}

\subsection{Numerical and I/O Parameters}

Table~\ref{tab:numerical_params} lists parameters controlling numerics, I/O,
and optional features.

\begin{table}[h!]
\centering
\caption{Numerical and I/O parameters (from \texttt{ggl90\_readparms.F}).}
\label{tab:numerical_params}
\begin{tabular}{@{}lllp{5.5cm}@{}}
\toprule
\textbf{Parameter} & \textbf{Default} & \textbf{Units/Type} & \textbf{Description} \\
\midrule
\texttt{GGL90diffTKEh} & 0.0 & m$^2$/s & Horizontal diffusivity for TKE.
                                         Only active if
                                         \texttt{ALLOW\_GGL90\_HORIZDIFF}
                                         is defined. Usually kept at zero. \\[6pt]

\texttt{GGL90\_dirichlet} & \texttt{.TRUE.} & logical & Use Dirichlet
                                                         boundary conditions
                                                         for TKE at boundaries.
                                                         If \texttt{.FALSE.},
                                                         uses Neumann conditions. \\[6pt]

\texttt{calcMeanVertShear} & \texttt{.FALSE.} & logical & Calculate mean
                                                           vertical shear at
                                                           cell centers instead
                                                           of shear of mean
                                                           velocity. Affects
                                                           $S^2$ computation. \\[6pt]

\texttt{GGL90TKEFile} & \texttt{' '} & string & Binary file containing
                                                 initial TKE field. If blank,
                                                 initializes to
                                                 \texttt{GGL90TKEmin}. \\[6pt]

\texttt{GGL90dumpFreq} & \texttt{dumpFreq} & s & Output frequency for GGL90
                                                  state files (matches main
                                                  model by default). \\[6pt]

\texttt{GGL90mixingMaps} & \texttt{.FALSE.} & logical & Output mixing
                                                         diagnostics to
                                                         standard out. \\[6pt]

\texttt{GGL90writeState} & \texttt{.FALSE.} & logical & Write GGL90 state
                                                         to separate output
                                                         files. \\[6pt]

\texttt{useIDEMIX} & \texttt{.FALSE.} & logical & Enable IDEMIX internal
                                                   wave model (see Section~\ref{sec:idemix}). \\[6pt]

\texttt{useLANGMUIR} & \texttt{.FALSE.} & logical & Enable Langmuir
                                                     circulation
                                                     parameterization. \\

\bottomrule
\end{tabular}
\end{table}

\subsection{Example Parameter File}

A typical \texttt{data.ggl90} file for a global ocean simulation might look
like:

\begin{lstlisting}[caption={Example \texttt{data.ggl90} parameter file.},language=bash,basicstyle=\ttfamily\small]
&GGL90_PARM01
# Core physics parameters
  GGL90ck              = 0.1,
  GGL90ceps            = 0.7,
  GGL90alpha           = 30.0,      # ECCOv4 value (default is 1.0)
  GGL90m2              = 3.75,

# TKE thresholds
  GGL90TKEmin          = 1.0e-7,    # ECCOv4 value (default 1.0e-11)
  GGL90TKEsurfMin      = 1.0e-4,
  GGL90TKEbottom       = 1.0e-6,    # ECCOv4 value

# Mixing length
  GGL90mixingLengthMin = 1.0e-8,
  mxlMaxFlag           = 2,          # Two-way sweep (recommended)
  mxlSurfFlag          = .TRUE.,    # Force surface mixing

# Caps
  GGL90viscMax         = 1.0e2,
  GGL90diffMax         = 1.0e2,

# I/O
  GGL90dumpFreq        = 86400.0,   # Daily output
/
\end{lstlisting}

\subsection{Application-Specific Configurations}

GGL90 has been successfully deployed in a variety of ocean modeling applications,
from idealized process studies to global data-assimilating state estimates. While
the default parameter values (based on Gaspar et al.\ 1990 and Blanke \& Delecluse
1993) provide a reasonable starting point, certain applications benefit from
targeted parameter adjustments.

For example, global ocean state estimation systems using adjoint-based
optimization often employ non-default settings to balance thermocline sharpness,
deep ocean mixing representation, and gradient quality. Appendix~\ref{app:eccov4_config}
documents the complete ECCOv4 Release 4 configuration as a well-validated
example for $\sim$1$^\circ$ global applications.

\subsection{Parameter Validation}

After reading \texttt{data.ggl90}, the routine \texttt{ggl90\_check.F} (lines
7--166) validates parameter consistency. It checks for:

\begin{itemize}[leftmargin=*]
  \item \textbf{Positive values}: $c_k, c_\varepsilon, \alpha, m_2, \TKE_{\mathrm{min}}$
        must all be $> 0$.
  \item \textbf{Valid flags}: \texttt{mxlMaxFlag} $\in \{0, 1, 2\}$.
  \item \textbf{Implicit timestepping}: \texttt{implicitDiffusion} and
        \texttt{implicitViscosity} must be \texttt{.TRUE.} in \texttt{data}
        (required for stability of TKE equation).
  \item \textbf{IDEMIX requirements}: If \texttt{useIDEMIX = .TRUE.}, checks
        that overlap regions (OLx, OLy) are $\geq 3$ for horizontal operators.
\end{itemize}

If any check fails, the model stops with a descriptive error message. This
prevents common configuration mistakes.

% ============================================================
\section{Main Driver: \texttt{GGL90\_CALC}}
\label{sec:ggl90_calc}
% ============================================================

The subroutine \texttt{GGL90\_CALC} (\texttt{ggl90\_calc.F:13--1177}) is the
computational heart of the GGL90 package. It is called once per timestep for
each tile in the domain decomposition. This section provides a detailed
walkthrough of its execution, with code excerpts and line number references.

\subsection{Subroutine Signature and Input Parameters}

\begin{lstlisting}[caption={\texttt{ggl90\_calc.F}, lines 13--63 -- Subroutine interface.}]
SUBROUTINE GGL90_CALC(
     I     bi, bj, sigmaR, myTime, myIter, myThid )

  ! Input parameters:
  !   bi, bj   :: Current tile indices
  !   sigmaR   :: Vertical gradient of iso-neutral density (3D)
  !   myTime   :: Current time in simulation
  !   myIter   :: Current time-step number
  !   myThid   :: Thread ID number

  ! Global variables updated:
  !   GGL90TKE(i,j,k,bi,bj)     :: Turbulent kinetic energy (m^2/s^2)
  !   GGL90viscArU(i,j,k,bi,bj) :: Eddy viscosity at U-points (m^2/s)
  !   GGL90viscArV(i,j,k,bi,bj) :: Eddy viscosity at V-points (m^2/s)
  !   GGL90diffKr(i,j,k,bi,bj)  :: Eddy diffusivity at T/S points (m^2/s)
\end{lstlisting}

\textbf{Key input:} The array \texttt{sigmaR} contains the vertical gradient
of potential density referenced to the local pressure. This is computed by the
main model's equation of state and passed to GGL90. From \texttt{sigmaR}, the
buoyancy frequency squared ($N^2$) is derived as:
\[
N^2 = -\frac{g}{\rho_0} \frac{\partial \rho}{\partial z}
    = -\frac{g}{\rho_0} \sigma_R
\]
where $g$ is gravitational acceleration and $\rho_0$ is reference density.

\subsection{Coordinate System Handling}

MITgcm supports both Z-coordinates (depth, $z < 0$) and P-coordinates
(pressure, $p > 0$). GGL90 must handle both conventions. Lines 198--209:

\begin{lstlisting}[caption={\texttt{ggl90\_calc.F}, lines 198--209 -- Coordinate system setup.}]
IF ( usingPCoords ) THEN
  kSrf = Nr      ! Surface at highest k index
  kTop = Nr
ELSE
  kSrf = 1       ! Surface at k=1 (Z-coordinates)
  kTop = 2
ENDIF
deltaTloc = dTtracerLev(kSrf)  ! Timestep for tracer level

coordFac = 1.0
IF ( usingPCoords ) coordFac = gravity * rhoConst
recip_coordFac = 1.0 / coordFac
\end{lstlisting}

\textbf{Coordinate scaling factor:} In P-coordinates, vertical derivatives
involve pressure gradients. The factor \texttt{coordFac} = $g \rho_0$ converts
pressure units to depth-equivalent units. This ensures that mixing length and
TKE budget terms have consistent dimensions regardless of coordinate choice.

\subsection{Initialization of Local Arrays}

Lines 268--300 initialize working arrays to safe default values:

\begin{lstlisting}[caption={\texttt{ggl90\_calc.F}, lines 268--300 (excerpts) -- Array initialization.}]
DO k=1,Nr
  DO j=1-OLy,sNy+OLy
    DO i=1-OLx,sNx+OLx
      rMixingLength(i,j,k)     = 0.0
      GGL90visctmp(i,j,k)      = 0.0
      KappaE(i,j,k)            = 0.0
      TKEPrandtlNumber(i,j,k)  = 1.0
      GGL90mixingLength(i,j,k) = GGL90mixingLengthMin
      Nsquare(i,j,k)           = 0.0
      SQRTTKE(i,j,k)           = 0.0
    ENDDO
  ENDDO
ENDDO
DO j=1-OLy,sNy+OLy
  DO i=1-OLx,sNx+OLx
    KappaM(i,j)        = 0.0
    uStarSquare(i,j)   = 0.0
    verticalShear(i,j) = 0.0
  ENDDO
ENDDO
\end{lstlisting}

Initializing arrays is essential for robustness, especially in regions where
topography creates "dry" cells (masked points).

\subsection{Interface Thickness Factors}

Lines 242--257 compute thickness factors at cell interfaces (\texttt{hFacI}),
which are needed for vertical operators:

\begin{lstlisting}[caption={\texttt{ggl90\_calc.F}, lines 242--257 -- Interface thickness calculation.}]
DO k=1,Nr
  km1 = MAX(k-1,1)
  DO j=1-OLy,sNy+OLy
    DO i=1-OLx,sNx+OLx
      hFac = MIN( 0.5, _hFacC(i,j,km1,bi,bj) )
           + MIN( 0.5, _hFacC(i,j,k,  bi,bj) )
      recip_hFacI(i,j,k) = 0.0
      IF ( hFac .NE. 0.0 )
     &    recip_hFacI(i,j,k) = 1.0 / hFac
#ifdef ALLOW_GGL90_IDEMIX
      hFacI(i,j,k) = hFac
#endif
    ENDDO
  ENDDO
ENDDO
\end{lstlisting}

\textbf{Explanation:} The interface at $k$ lies between cell centers $k-1$ and
$k$. The thickness factor \texttt{hFacI(k)} is the average of the two adjacent
cell fractions, capped at 0.5 each. This ensures proper handling of partial
cells near topography.

\subsection{Step-by-Step Execution Summary}

The main computational loop (lines 300--1177) proceeds as follows. Each step
is detailed in subsequent sections of this document.

\begin{enumerate}[leftmargin=*, label=\textbf{Step \arabic*.}]
  \item \textbf{Compute buoyancy frequency squared ($N^2$)} -- From
        \texttt{sigmaR}. See Section~\ref{sec:stratification}.

  \item \textbf{Compute vertical shear squared ($S^2$)} -- From velocity
        fields U and V. See Section~\ref{sec:stratification}.

  \item \textbf{Compute surface friction velocity ($u_*^2$)} -- From wind
        stress. See Section~\ref{sec:stratification}.

  \item \textbf{Compute mixing length ($\Lmix$)} -- Call
        \texttt{GGL90\_MIXINGLENGTH} with method selected by
        \texttt{mxlMaxFlag}. See Section~\ref{sec:mixing_length}.

  \item \textbf{Compute eddy coefficients ($\KM$, $\KH$)} -- From TKE and
        mixing length. See Section~\ref{sec:eddy_coeffs}.

  \item \textbf{Compute TKE budget terms} -- Production, buoyancy,
        dissipation. See Section~\ref{sec:tke_budget}.

  \item \textbf{Build implicit matrix} -- Set up tridiagonal system for TKE
        equation. See Section~\ref{sec:tke_budget}.

  \item \textbf{Apply boundary conditions} -- Surface flux and bottom TKE.
        See Section~\ref{sec:boundary_conditions}.

  \item \textbf{Solve tridiagonal system} -- Advance TKE to next timestep.
        See Section~\ref{sec:tke_budget}.

  \item \textbf{Transfer to model arrays} -- Call \texttt{GGL90\_CALC\_VISC}
        and \texttt{GGL90\_CALC\_DIFF}. See Section~\ref{sec:interface}.
\end{enumerate}

The remainder of this document examines each step in detail with code snippets
and physical interpretation.

% ============================================================
\section{Stratification and Shear Computation}
\label{sec:stratification}
% ============================================================

The TKE budget depends critically on two fundamental quantities: the buoyancy
frequency squared ($N^2$, representing stratification) and the vertical shear
squared ($S^2$, representing velocity gradients). These appear in the
production and buoyancy terms of the TKE equation. This section describes how
GGL90 computes these quantities from the model state.

\subsection{Buoyancy Frequency Squared ($N^2$)}

The buoyancy frequency squared (also called the Brunt-V\"ais\"al\"a frequency
squared) quantifies the strength of stable stratification:
\[
N^2 = -\frac{g}{\rho_0} \frac{\partial \rho}{\partial z}
\]
where $g$ is gravitational acceleration, $\rho_0$ is a reference density, and
$\partial \rho / \partial z$ is the vertical density gradient. Positive $N^2$
indicates stable stratification (density decreases upward), while negative
$N^2$ indicates convective instability.

\subsubsection{Implementation}

In MITgcm, the equation of state computes \texttt{sigmaR}, the vertical
gradient of potential density referenced to local pressure. GGL90 converts
this to $N^2$ (lines 346--348):

\begin{lstlisting}[caption={\texttt{ggl90\_calc.F}, lines 346--348 -- Buoyancy frequency squared.}]
! Loop over k=2,Nr (interfaces)
Nsquare(i,j,k) = gravity * gravitySign * recip_rhoConst
     &         * sigmaR(i,j,k) * coordFac
\end{lstlisting}

\textbf{Notes:}
\begin{itemize}[leftmargin=*]
  \item \texttt{gravitySign} = $+1$ for normal gravity, $-1$ if gravity is
        reversed (e.g., for testing).
  \item \texttt{recip\_rhoConst} = $1/\rho_0$ where $\rho_0 \approx 1035$~kg/m$^3$
        is the reference density.
  \item \texttt{coordFac} accounts for coordinate system (1.0 for Z-coords,
        $g\rho_0$ for P-coords).
  \item $N^2$ is defined at interfaces (between levels $k-1$ and $k$), matching
        the staggering of TKE.
\end{itemize}

\subsubsection{Physical Interpretation}

\begin{itemize}[leftmargin=*]
  \item \textbf{Stable stratification} ($N^2 > 0$): Buoyancy acts as a
        restoring force, suppressing vertical motion. Turbulence is damped,
        and mixing is weak.

  \item \textbf{Neutral stratification} ($N^2 \approx 0$): Density is
        vertically uniform. Turbulence can develop freely if shear is present.

  \item \textbf{Convective instability} ($N^2 < 0$): Denser water overlies
        lighter water. Buoyancy drives overturning motions, generating TKE and
        enhancing mixing. This occurs during winter cooling, brine rejection
        from sea ice, or evaporative heat loss.
\end{itemize}

In the TKE budget, $N^2$ enters through the buoyancy term:
\[
B = -\kappa_h N^2
\]
When $N^2 < 0$, $B > 0$ (source of TKE). When $N^2 > 0$ and mixing transports
light water downward, $B < 0$ (sink of TKE).

\subsection{Vertical Shear Squared ($S^2$)}

Vertical shear of the horizontal velocity field is a primary source of
turbulence in the ocean. The shear squared is defined as:
\[
S^2 = \left( \frac{\partial u}{\partial z} \right)^2
    + \left( \frac{\partial v}{\partial z} \right)^2
\]
where $u$ and $v$ are zonal and meridional velocity components.

\subsubsection{Implementation: Two Methods}

GGL90 offers two methods for computing $S^2$, controlled by the parameter
\texttt{calcMeanVertShear}.

\textbf{Method 1: Shear of mean velocity (default, \texttt{calcMeanVertShear = .FALSE.})}

Lines 484--498:

\begin{lstlisting}[caption={\texttt{ggl90\_calc.F}, lines 484--498 -- Shear of mean velocity.}]
! Average velocity to cell centers, then compute shear
tempU = ( ( uVel(i,j,km1,bi,bj) + uVel(i+1,j,km1,bi,bj) )
     &   -( uVel(i,j,k  ,bi,bj) + uVel(i+1,j,k  ,bi,bj) )
     &  ) * 0.5 * recip_drC(k) * coordFac

tempV = ( ( vVel(i,j,km1,bi,bj) + vVel(i,j+1,km1,bi,bj) )
     &   -( vVel(i,j,k  ,bi,bj) + vVel(i,j+1,k  ,bi,bj) )
     &  ) * 0.5 * recip_drC(k) * coordFac

verticalShear(i,j) = tempU*tempU + tempV*tempV
\end{lstlisting}

This method:
\begin{enumerate}[label=\textbf{\arabic*.}]
  \item Averages $u$ from adjacent U-points to the T-point.
  \item Averages $v$ from adjacent V-points to the T-point.
  \item Computes $\partial u/\partial z$ and $\partial v/\partial z$ from the
        averaged velocities.
  \item Squares and sums to get $S^2$.
\end{enumerate}

\textbf{Method 2: Mean of vertical shear (\texttt{calcMeanVertShear = .TRUE.})}

Lines 468--482:

\begin{lstlisting}[caption={\texttt{ggl90\_calc.F}, lines 468--482 -- Mean of shear components.}]
! Compute shear at each U/V point, then average
tempU  = ( uVel( i ,j,km1,bi,bj) - uVel( i ,j,k,bi,bj) )
tempUp = ( uVel(i+1,j,km1,bi,bj) - uVel(i+1,j,k,bi,bj) )
tempV  = ( vVel(i, j ,km1,bi,bj) - vVel(i, j ,k,bi,bj) )
tempVp = ( vVel(i,j+1,km1,bi,bj) - vVel(i,j+1,k,bi,bj) )

verticalShear(i,j) = (
     &       ( tempU*tempU + tempUp*tempUp )
     &     + ( tempV*tempV + tempVp*tempVp )
     &               ) * 0.5 * recip_drC(k) * recip_drC(k)
     &                       * coordFac * coordFac
\end{lstlisting}

This method:
\begin{enumerate}[label=\textbf{\arabic*.}]
  \item Computes $(\partial u/\partial z)^2$ at both U-points adjacent to the
        T-point.
  \item Computes $(\partial v/\partial z)^2$ at both V-points adjacent to the
        T-point.
  \item Averages the four squared shears.
\end{enumerate}

\textbf{Comparison:}

Method 1 (default) computes shear of averaged velocity. Method 2 computes
average of squared shear. Method 2 is more conservative in that it captures
localized shear spikes better, but can produce noisier fields. The default
(Method 1) is smoother and is recommended for most applications.

\subsubsection{Physical Interpretation}

Shear extracts kinetic energy from the mean flow and converts it to turbulent
kinetic energy. The production term in the TKE budget is:
\[
P = \kappa_m S^2
\]
Strong shear (e.g., in western boundary currents, equatorial jets, or beneath
strong winds) generates vigorous turbulence. Conversely, regions with weak
currents have low $S^2$ and rely on other processes (convection, internal
waves) for mixing.

\subsection{Surface Stress ($u_*^2$)}

The surface friction velocity $u_*$ characterizes the rate at which wind
transfers momentum (and energy) to the ocean:
\[
u_* = \sqrt{\frac{\tau}{\rho_0}}
\]
where $\tau$ is the wind stress magnitude.

\subsubsection{Implementation}

Surface stress is computed in the surface layer (interface at \texttt{kTop})
by averaging U and V stresses (lines not shown in prior excerpts but typically
around lines 700--750):

\begin{lstlisting}[caption={Conceptual code for surface stress (typically lines 700--750).}]
! At surface interface (k = kTop)
IF ( usingZCoords ) THEN
  uStarSquare(i,j) = SQRT(
       ( 0.5*(surfaceForcingU(i,j,bi,bj) + surfaceForcingU(i+1,j,bi,bj)) )**2
     + ( 0.5*(surfaceForcingV(i,j,bi,bj) + surfaceForcingV(i,j+1,bi,bj)) )**2
     )
ENDIF
! surfaceForcingU/V are in units of N/m^2 (stress)
! Convert to u_star^2 by dividing by rho_0
\end{lstlisting}

\textbf{Use in TKE budget:}

$u_*^2$ sets the surface boundary condition for TKE. The wind-driven TKE flux
at the surface is:
\[
\kappa_E \frac{\partial \TKE}{\partial z}\bigg|_{z=0} = m_2 u_*^3
\]
This boundary condition injects TKE into the surface layer, which then
diffuses and is transported downward, energizing the mixed layer.

% ============================================================
\section{Mixing Length Computation: \texttt{GGL90\_MIXINGLENGTH}}
% ============================================================

The mixing length $\ell$ is one of the most critical elements of the GGL90 scheme, determining the spatial scale over which turbulent eddies can transport momentum and scalars. Unlike some other turbulence closures that prescribe the mixing length purely from geometric considerations, GGL90 computes it from both turbulent energetics (TKE and buoyancy frequency) and geometric constraints imposed by boundaries and the water column structure.

The mixing length computation proceeds in two distinct stages:

\begin{enumerate}
\item \textbf{Initial estimate from Gaspar et al.\ (1990):} Compute an energetically-based mixing length $\ell^{(Gaspar)}$ that balances TKE with stratification (Section~\ref{sec:mxl_gaspar}).
\item \textbf{Geometric constraints:} Apply boundary-aware constraints using one of three methods (Sections~\ref{sec:mxl_method0}--\ref{sec:mxl_method2}), controlled by runtime parameter \verb|mxlMaxFlag|.
\end{enumerate}

The final mixing length is then used to compute eddy viscosity and diffusivity (Section~\ref{sec:eddy_coeffs}), and also appears in the TKE dissipation term (Section~\ref{sec:tke_dissipation}).

\paragraph{Historical context.} The original Gaspar et al.\ (1990) paper proposed the energetic estimate, but left open the question of how to constrain the mixing length near boundaries. Blanke \& Delecluse (1993) introduced a sophisticated two-way sweep algorithm (Method 2) that prevents unphysical penetration of mixing across stable layers, and this has become the de facto standard in many ocean models (see Appendix~\ref{app:eccov4_config}).

\paragraph{Code organization.} The mixing length computation is split between two subroutines:
\begin{itemize}
\item \verb|GGL90_CALC| (lines 351--357): Computes the initial Gaspar estimate $\ell^{(Gaspar)}$.
\item \verb|GGL90_MIXINGLENGTH| (lines 7--421): Applies geometric constraints via one of three methods.
\end{itemize}

This section documents both stages in detail, with particular emphasis on Method 2, the most physically sophisticated approach.

% ------------------------------------------------------------
\subsection{Initial Estimate from Gaspar et al.\ (1990)}
\label{sec:mxl_gaspar}
% ------------------------------------------------------------

Before any geometric constraints are applied, GGL90 computes an initial mixing length estimate based purely on the local turbulent kinetic energy and stratification. This is the energetic mixing length from Gaspar et al.\ (1990), which represents the vertical scale over which a turbulent eddy with kinetic energy $e$ can work against the ambient stratification.

\subsubsection{Physical motivation}

Consider a turbulent eddy with characteristic velocity scale $\sqrt{e}$ attempting to overturn a stratified water column with buoyancy frequency $N = \sqrt{N^2}$. The eddy can penetrate vertically until its kinetic energy is exhausted by work against buoyancy forces. Dimensional analysis gives:

\begin{equation}
\ell^{(Gaspar)} \sim \frac{\sqrt{e}}{N}
\end{equation}

Gaspar et al.\ (1990) found that a factor of $\sqrt{2}$ improves agreement with observations and large-eddy simulations, yielding their equation (2.35):

\begin{equation}
\ell^{(Gaspar)} = \sqrt{2} \frac{\sqrt{e}}{\sqrt{\max(N^2, \varepsilon)}}
\label{eq:mxl_gaspar}
\end{equation}

where $\varepsilon$ is a small positive constant (typically $10^{-12}$ to $10^{-14}$ m$^{-2}$ s$^{-2}$) to prevent division by zero in weakly stratified or unstratified regions.

\paragraph{Physical interpretation.}
\begin{itemize}
\item In strongly stratified regions ($N^2 \gg \varepsilon$): $\ell^{(Gaspar)} \propto e/N^2$, so mixing is suppressed by stratification.
\item In weakly stratified or convective regions ($N^2 \approx 0$ or $N^2 < 0$): The denominator is floored at $\sqrt{\varepsilon}$, allowing $\ell^{(Gaspar)}$ to grow in proportion to $\sqrt{e}$.
\item At high TKE (e.g., near-surface wind-driven turbulence): $\ell^{(Gaspar)}$ increases, allowing deeper mixing.
\end{itemize}

\subsubsection{Implementation in \texttt{GGL90\_CALC}}

The initial Gaspar mixing length is computed in the main driver routine \verb|GGL90_CALC|, immediately after the buoyancy frequency $N^2$ has been diagnosed from the density field:

\begin{lstlisting}[language=Fortran,caption={(ggl90\_calc.F:L346--357) Initial mixing length estimate}]
C     buoyancy frequency
         Nsquare(i,j,k) = gravity*gravitySign*recip_rhoConst
     &                  * sigmaR(i,j,k) * coordFac
C     vertical shear term (dU/dz)^2+(dV/dz)^2 is computed later
C     to save some memory
C     Initialise mixing length (eq. 2.35)
         GGL90mixingLength(i,j,k) = SQRTTWO *
     &        SQRTTKE(i,j,k)/SQRT( MAX(Nsquare(i,j,k),GGL90eps) )
     &        * mskLoc
        ENDDO
       ENDDO
      ENDDO
\end{lstlisting}

\paragraph{Key code details:}
\begin{itemize}
\item \verb|SQRTTWO| is the constant $\sqrt{2} \approx 1.41421356$ (defined in \verb|EEPARAMS.h|).
\item \verb|SQRTTKE(i,j,k)| is $\sqrt{e}$, precomputed earlier in the routine (lines 339--343) to avoid repeated square root operations.
\item \verb|GGL90eps| is the floor value $\varepsilon$ for $N^2$, set by runtime parameter (default $10^{-12}$ m$^{-2}$ s$^{-2}$).
\item \verb|mskLoc| is a local cell mask (\verb|maskC(i,j,k)*maskC(i,j,k-1)|) that zeroes out the mixing length in dry cells or at dry interfaces.
\item \verb|coordFac| is a coordinate scaling factor (1 in z-coordinates, $g \rho_0$ in pressure coordinates) to convert between geometric and pressure vertical coordinates.
\end{itemize}

\subsubsection{Coordinate transformation}

Note that \verb|Nsquare| is multiplied by \verb|coordFac| to ensure dimensional consistency:
\begin{itemize}
\item In z-coordinates: \verb|coordFac = 1|, so \verb|Nsquare| has units of s$^{-2}$.
\item In pressure coordinates: \verb|coordFac = g*rhoConst|, so \verb|Nsquare| is scaled appropriately for pressure-based vertical coordinates.
\end{itemize}

This ensures that the mixing length $\ell^{(Gaspar)}$ always has units of meters, regardless of the vertical coordinate system.

\subsubsection{Typical values}

For representative ocean conditions:
\begin{itemize}
\item \textbf{Surface mixed layer} (winter, deep convection): $e \sim 10^{-3}$ m$^2$ s$^{-2}$, $N^2 \sim 10^{-6}$ s$^{-2}$ $\Rightarrow$ $\ell^{(Gaspar)} \sim 45$ m.
\item \textbf{Thermocline} (strong stratification): $e \sim 10^{-5}$ m$^2$ s$^{-2}$, $N^2 \sim 10^{-4}$ s$^{-2}$ $\Rightarrow$ $\ell^{(Gaspar)} \sim 0.45$ m.
\item \textbf{Deep ocean} (weak stratification): $e \sim 10^{-6}$ m$^2$ s$^{-2}$, $N^2 \sim 10^{-7}$ s$^{-2}$ $\Rightarrow$ $\ell^{(Gaspar)} \sim 4.5$ m.
\end{itemize}

These values are purely local and do not yet account for distance to boundaries, which is why the subsequent geometric constraint step is essential.

% ------------------------------------------------------------
\subsection{Method 0: Simple Distance to Boundaries}
\label{sec:mxl_method0}
% ------------------------------------------------------------

Method 0 (\verb|mxlMaxFlag = 0|) is the simplest geometric constraint: cap the Gaspar mixing length at the total water column thickness, without any consideration of boundary proximity or vertical propagation of mixing constraints.

\subsubsection{Algorithm}

For each interior cell $k = 2, \ldots, N_r$:
\begin{equation}
\ell(i,j,k) = \min\left( \ell^{(Gaspar)}(i,j,k), \; H_{\text{col}}(i,j) \right)
\label{eq:mxl_method0}
\end{equation}

where $H_{\text{col}}$ is the total water column thickness:
\begin{equation}
H_{\text{col}}(i,j) = R_{\text{surf}}(i,j) - R_{\text{bottom}}(i,j)
\end{equation}

\paragraph{Physical interpretation.} This constraint simply prevents the mixing length from exceeding the total depth of the ocean at this horizontal location. It is a minimal sanity check, but does not capture the physically important constraint that mixing cannot extend beyond nearby boundaries (surface, bottom, or stable layers).

\subsubsection{Implementation}

The code for Method 0 is straightforward:

\begin{lstlisting}[language=Fortran,caption={(ggl90\_mixinglength.F:L166--179) Method 0 constraint}]
      IF ( locMxlMaxFlag .EQ. 0 ) THEN

       DO k=2,Nr
        DO j=jMin,jMax
         DO i=iMin,iMax
C     Use thickness of water column (inverse of recip_Rcol) as the
C     maximum length.
          MaxLength =
     &         ( Ro_surf(i,j,bi,bj) - R_low(i,j,bi,bj) )*recip_coordFac
          GGL90mixingLength(i,j,k) = MIN(GGL90mixingLength(i,j,k),
     &                                   MaxLength)
         ENDDO
        ENDDO
       ENDDO
\end{lstlisting}

\paragraph{Key variables:}
\begin{itemize}
\item \verb|Ro_surf(i,j,bi,bj)|: Position of the surface (z-coords) or top pressure (p-coords).
\item \verb|R_low(i,j,bi,bj)|: Position of the bottom (z-coords) or bottom pressure (p-coords).
\item \verb|recip_coordFac|: Inverse of coordinate scaling factor (1 in z-coords, $1/(g\rho_0)$ in p-coords).
\end{itemize}

\subsubsection{When to use Method 0}

Method 0 is rarely used in practice because it is too permissive: it allows unrealistically large mixing lengths in deep ocean columns, where turbulence should be confined to a thin layer near boundaries or shear zones. However, it can be useful for:
\begin{itemize}
\item \textbf{Debugging}: If you suspect that the more sophisticated methods (1 or 2) are causing numerical issues, Method 0 provides a minimal baseline.
\item \textbf{Shallow water applications}: In very shallow domains (e.g., $H < 50$ m), the entire water column may be actively mixed, and Method 0 suffices.
\item \textbf{Theoretical studies}: Method 0 isolates the energetic mixing length behavior without geometric constraints.
\end{itemize}

For realistic ocean simulations, Method 1 or (preferably) Method 2 should be used instead.

% ------------------------------------------------------------
\subsection{Method 1: One-Way Downward Sweep}
\label{sec:mxl_method1}
% ------------------------------------------------------------

Method 1 (\verb|mxlMaxFlag = 1|) constrains the mixing length at each cell by the minimum distance to either the surface above or the bottom below. This is a significant improvement over Method 0, as it prevents mixing from penetrating beyond the nearest boundary.

\subsubsection{Algorithm}

For each interior cell $k = 2, \ldots, N_r$:
\begin{equation}
\ell(i,j,k) = \min\left( \ell^{(Gaspar)}(i,j,k), \; d_{\text{surf}}(i,j,k), \; d_{\text{bottom}}(i,j,k) \right)
\label{eq:mxl_method1}
\end{equation}

where:
\begin{equation}
\begin{aligned}
d_{\text{surf}}(i,j,k)   &= R_{\text{surf}}(i,j) - r_F(k) \\
d_{\text{bottom}}(i,j,k) &= r_F(k) - R_{\text{bottom}}(i,j)
\end{aligned}
\end{equation}

and $r_F(k)$ is the vertical position of the lower face of cell $k$ (i.e., the interface between cells $k$ and $k+1$).

\paragraph{Physical interpretation.} At each grid cell, the mixing length is limited by whichever boundary (top or bottom) is closer. For example:
\begin{itemize}
\item Near the surface: $d_{\text{surf}} < d_{\text{bottom}}$, so $\ell \leq d_{\text{surf}}$.
\item Near the bottom: $d_{\text{bottom}} < d_{\text{surf}}$, so $\ell \leq d_{\text{bottom}}$.
\item Mid-depth: The constraint depends on which boundary is closer.
\end{itemize}

This prevents the mixing length from extending into the "shadow" of a nearby boundary.

\subsubsection{Implementation}

\begin{lstlisting}[language=Fortran,caption={(ggl90\_mixinglength.F:L181--193) Method 1 constraint}]
      ELSEIF ( locMxlMaxFlag .EQ. 1 ) THEN

       DO k=2,Nr
        DO j=jMin,jMax
         DO i=iMin,iMax
          MaxLength=MIN(Ro_surf(i,j,bi,bj)-rF(k),rF(k)-R_low(i,j,bi,bj))
     &         * recip_coordFac
c         MaxLength=MAX(MaxLength,20. _d 0)
          GGL90mixingLength(i,j,k) = MIN(GGL90mixingLength(i,j,k),
     &                                   MaxLength)
         ENDDO
        ENDDO
       ENDDO
\end{lstlisting}

\paragraph{Key details:}
\begin{itemize}
\item \verb|rF(k)|: Vertical position of the lower face of cell $k$ (interface $k/k+1$).
\item The commented-out line \verb|MaxLength=MAX(MaxLength,20. _d 0)| would impose a lower bound of 20 m on the mixing length; this is disabled by default but could be useful in coarse-resolution models to prevent overly small mixing lengths.
\end{itemize}

\subsubsection{Limitations}

Method 1 is a substantial improvement over Method 0, but it has a critical limitation: \textbf{it does not account for vertical propagation of boundary constraints}. Consider a stratified water column with a shallow mixed layer of depth 50 m and a total column depth of 1000 m:

\begin{itemize}
\item At 30 m depth (within the mixed layer): $d_{\text{surf}} = 30$ m, $d_{\text{bottom}} = 970$ m $\Rightarrow$ $\ell \leq 30$ m.
\item At 70 m depth (just below the mixed layer base): $d_{\text{surf}} = 70$ m, $d_{\text{bottom}} = 930$ m $\Rightarrow$ $\ell \leq 70$ m.
\end{itemize}

The problem: even though there is a strong pycnocline at 50 m depth that should block mixing, Method 1 allows $\ell$ to jump from 30 m to 70 m as we descend through the pycnocline. This violates the physical principle that mixing cannot propagate vertically through stable layers.

Method 2 (next subsection) addresses this limitation via a two-way sweep algorithm.

\subsubsection{When to use Method 1}

Method 1 is a reasonable choice for:
\begin{itemize}
\item \textbf{Weakly stratified domains}: If the water column is relatively homogeneous (e.g., polar regions in winter), the lack of vertical propagation is less problematic.
\item \textbf{Coarse vertical resolution}: In models with $\Delta z > 50$ m, the distinction between Methods 1 and 2 may be small.
\item \textbf{Computational cost}: Method 1 is slightly cheaper than Method 2 (single pass vs.\ double pass), though the difference is negligible in most applications.
\end{itemize}

For high-resolution simulations with realistic stratification, Method 2 is strongly preferred.

% ------------------------------------------------------------
\subsection{Method 2: Two-Way Sweep (Blanke \& Delecluse 1993)}
\label{sec:mxl_method2}
% ------------------------------------------------------------

Method 2 (\verb|mxlMaxFlag = 2|) is the most sophisticated and physically realistic approach, implementing the two-way sweep algorithm of Blanke \& Delecluse (1993). This method ensures that the mixing length is constrained not only by distance to the nearest boundary, but also by the accumulated constraint from all cells above and below, preventing unphysical propagation of mixing across stable layers.

\paragraph{Recommended method.} Method 2 is the standard choice for most modern ocean GCMs (e.g., Appendix~\ref{app:eccov4_config}), as it produces the most realistic mixed layer depths and stratification profiles.

\subsubsection{Physical motivation}

The key insight of Blanke \& Delecluse (1993) is that the mixing length at a given depth should be constrained by \emph{all} the mixing lengths between that depth and the boundaries, not just the direct geometric distance to the boundaries. Specifically:

\begin{itemize}
\item If there is a thin layer with $\ell_{\text{small}}$ above the current cell, then the current cell's mixing length should not exceed $\ell_{\text{small}} + \Delta z$, where $\Delta z$ is the thickness of one grid cell.
\item Similarly, if there is a thin layer below with $\ell_{\text{small}}$, the current cell should inherit this constraint.
\end{itemize}

This implements the physical constraint that mixing cannot "jump over" a stable layer with small mixing length.

\subsubsection{Algorithm overview}

Method 2 proceeds in three stages:

\begin{enumerate}
\item \textbf{Downward sweep}: Starting from the surface, propagate mixing length constraints downward:
\begin{equation}
\ell_{\text{down}}(k) = \min\left( \ell^{(Gaspar)}(k), \; \ell_{\text{down}}(k-1) + \Delta z_{k-1} \right)
\label{eq:mxl_downsweep}
\end{equation}

\item \textbf{Upward sweep}: Starting from the bottom, propagate mixing length constraints upward:
\begin{equation}
\ell_{\text{up}}(k) = \min\left( \ell_{\text{up}}(k+1) + \Delta z_k, \; \ell_{\text{down}}(k) \right)
\label{eq:mxl_upsweep}
\end{equation}

\item \textbf{Final constraint}: Take the minimum of the downward-swept and upward-swept values:
\begin{equation}
\ell(k) = \min\left( \ell_{\text{up}}(k), \; \ell_{\text{down}}(k) \right)
\label{eq:mxl_final}
\end{equation}
\end{enumerate}

The downward sweep ensures that surface constraints propagate down, the upward sweep ensures that bottom constraints propagate up, and the final minimum ensures that the most restrictive constraint (from either direction) is applied.

\subsubsection{Implementation in MITgcm}

The code for Method 2 is more complex than Methods 0 or 1, and it must handle two different cases depending on whether the model uses z-coordinates or pressure coordinates. The vertical sweep direction differs between the two coordinate systems:

\begin{itemize}
\item \textbf{Z-coordinates}: Surface is at $k=1$, bottom is at $k=N_r$. Downward sweep goes from $k=2$ to $N_r$; upward sweep goes from $k=N_r-1$ to $2$.
\item \textbf{Pressure coordinates}: Surface is at $k=N_r$, bottom is at $k=1$. Downward sweep goes from $k=N_r-1$ to $2$; upward sweep goes from $k=2$ to $N_r$.
\end{itemize}

We document the z-coordinate case here (lines 240--299), which is the most common. The p-coordinate case (lines 197--238) follows the same logic with reversed loop bounds.

\paragraph{Downward sweep (z-coordinates):}

\begin{lstlisting}[language=Fortran,caption={(ggl90\_mixinglength.F:L245--258) Downward sweep in z-coordinates}]
C     Downward sweep
        DO k=2,Nr
#ifdef ALLOW_AUTODIFF_TAMC
         kkey = k + (tkey-1)*Nr
CADJ STORE mxLength_Dn(:,:,k-1)
CADJ &     = comlev1_bibj_k, key = kkey, kind=isbyte
#endif /* ALLOW_AUTODIFF_TAMC */
         DO j=jMin,jMax
          DO i=iMin,iMax
           mxLength_Dn(i,j,k) = MIN(GGL90mixingLength(i,j,k),
     &          mxLength_Dn(i,j,k-1)+drF(k-1)*recip_coordFac)
          ENDDO
         ENDDO
        ENDDO
\end{lstlisting}

\paragraph{Key details:}
\begin{itemize}
\item \verb|mxLength_Dn(i,j,k)| is the downward-swept mixing length, stored in a temporary array.
\item The initialization \verb|mxLength_Dn(i,j,1) = GGL90mixingLengthMin| (line 131) provides the surface boundary condition.
\item \verb|drF(k-1)| is the thickness of cell $k-1$, so \verb|mxLength_Dn(i,j,k-1)+drF(k-1)| is the mixing length from the cell above plus the increment needed to reach cell $k$.
\end{itemize}

\paragraph{Upward sweep (z-coordinates):}

\begin{lstlisting}[language=Fortran,caption={(ggl90\_mixinglength.F:L260--283) Upward sweep in z-coordinates}]
C     Upward sweep, extra treatment of k=Nr for z-coordinates
C     because level Nr+1 is not available
        DO j=jMin,jMax
         DO i=iMin,iMax
          GGL90mixingLength(i,j,Nr) = MIN(GGL90mixingLength(i,j,Nr),
     &         GGL90mixingLengthMin+drF(Nr)*recip_coordFac)
         ENDDO
        ENDDO
        DO k=Nr-1,2,-1
#ifdef ALLOW_AUTODIFF_TAMC
         kkey = k + (tkey-1)*Nr
C     It is important that the two k-levels of these fields are stored
C     in one statement because otherwise taf will only store one, which
C     is wrong (i.e. was wrong in previous versions).
CADJ STORE GGL90mixingLength(:,:,k+1), GGL90mixingLength(:,:,k)
CADJ &     = comlev1_bibj_k, key = kkey, kind=isbyte
#endif /* ALLOW_AUTODIFF_TAMC */
         DO j=jMin,jMax
          DO i=iMin,iMax
           GGL90mixingLength(i,j,k) = MIN(GGL90mixingLength(i,j,k),
     &          GGL90mixingLength(i,j,k+1)+drF(k)*recip_coordFac)
          ENDDO
         ENDDO
        ENDDO
\end{lstlisting}

\paragraph{Key details:}
\begin{itemize}
\item The upward sweep overwrites \verb|GGL90mixingLength| directly (rather than using a temporary array), imposing the constraint from below.
\item Special treatment at $k=N_r$ (bottom cell): Since there is no cell below, the constraint is simply \verb|GGL90mixingLengthMin + drF(Nr)|, which allows at most a one-cell-thick mixing layer at the bottom.
\item The upward sweep runs in reverse order ($k = N_r-1$ down to $2$).
\end{itemize}

\paragraph{Final constraint:}

\begin{lstlisting}[language=Fortran,caption={(ggl90\_mixinglength.F:L291--299) Final minimum of downward and upward sweeps}]
C     Impose minimum from downward sweep
       DO k=2,Nr
        DO j=jMin,jMax
         DO i=iMin,iMax
          GGL90mixingLength(i,j,k) = MIN(GGL90mixingLength(i,j,k),
     &                                  mxLength_Dn(i,j,k))
         ENDDO
        ENDDO
       ENDDO
\end{lstlisting}

After the upward sweep, \verb|GGL90mixingLength| contains the upward-swept value $\ell_{\text{up}}$. This final loop takes the minimum with the downward-swept value $\ell_{\text{down}}$ stored in \verb|mxLength_Dn|, yielding the final constrained mixing length.

\subsubsection{Worked example: Idealized stratification}

Consider an idealized water column with a shallow mixed layer and a sharp pycnocline:

\begin{center}
\begin{tabular}{cccc}
\toprule
Depth (m) & $N^2$ (s$^{-2}$) & $\ell^{(Gaspar)}$ (m) & $\Delta z$ (m) \\
\midrule
10 (k=2)  & $10^{-6}$ & 45 & 20 \\
30 (k=3)  & $10^{-6}$ & 45 & 20 \\
50 (k=4)  & $10^{-4}$ & 4.5 & 20 \\
70 (k=5)  & $10^{-5}$ & 14 & 20 \\
90 (k=6)  & $10^{-5}$ & 14 & 20 \\
\bottomrule
\end{tabular}
\end{center}

\paragraph{Downward sweep:}
\begin{align*}
\ell_{\text{down}}(2) &= \min(45, 0 + 0) = 0 \quad \text{(initialized at surface)} \\
\ell_{\text{down}}(3) &= \min(45, 0 + 20) = 20 \\
\ell_{\text{down}}(4) &= \min(4.5, 20 + 20) = 4.5 \quad \text{(strong stratification limits mixing)} \\
\ell_{\text{down}}(5) &= \min(14, 4.5 + 20) = 14 \\
\ell_{\text{down}}(6) &= \min(14, 14 + 20) = 14
\end{align*}

\paragraph{Upward sweep:}
Starting from the bottom (assuming a bottom boundary condition $\ell_{\text{up}}(N_r) = \ell_{\text{min}} + \Delta z$):
\begin{align*}
\ell_{\text{up}}(6) &= \ell_{\text{down}}(6) = 14 \\
\ell_{\text{up}}(5) &= \min(14 + 20, 14) = 14 \\
\ell_{\text{up}}(4) &= \min(14 + 20, 4.5) = 4.5 \\
\ell_{\text{up}}(3) &= \min(4.5 + 20, 20) = 20 \\
\ell_{\text{up}}(2) &= \min(20 + 20, 0) = 0
\end{align*}

\paragraph{Final mixing length:}
\begin{align*}
\ell(2) &= \min(\ell_{\text{up}}(2), \ell_{\text{down}}(2)) = \min(0, 0) = 0 \\
\ell(3) &= \min(20, 20) = 20 \\
\ell(4) &= \min(4.5, 4.5) = 4.5 \\
\ell(5) &= \min(14, 14) = 14 \\
\ell(6) &= \min(14, 14) = 14
\end{align*}

Notice how the pycnocline at $k=4$ (with $\ell^{(Gaspar)} = 4.5$ m) blocks the downward propagation of the large mixed-layer mixing lengths from $k=2,3$. The upward sweep ensures that this constraint also affects cells above the pycnocline.

\subsubsection{Comparison of Methods 0, 1, and 2}

Table~\ref{tab:mxl_methods} summarizes the three methods and their typical use cases.

\begin{table}[ht]
\centering
\caption{Comparison of GGL90 mixing length methods}
\label{tab:mxl_methods}
\begin{tabular}{llp{5cm}p{4cm}}
\toprule
\textbf{Method} & \textbf{Flag} & \textbf{Algorithm} & \textbf{Best for} \\
\midrule
0 & \verb|mxlMaxFlag=0| & Cap at total column depth & Debugging, very shallow domains \\
\addlinespace
1 & \verb|mxlMaxFlag=1| & Cap at distance to nearest boundary (surface or bottom) & Weakly stratified domains, coarse vertical resolution \\
\addlinespace
2 & \verb|mxlMaxFlag=2| & Two-way sweep (Blanke \& Delecluse 1993) & Realistic stratification, high vertical resolution, \textbf{recommended} (Appendix~\ref{app:eccov4_config}) \\
\bottomrule
\end{tabular}
\end{table}

\paragraph{Computational cost.} Method 2 requires two vertical passes (downward + upward) plus a final min operation, making it roughly twice as expensive as Method 1. However, the mixing length computation is a small fraction of total GGL90 cost (dominated by the implicit TKE solve), so this difference is typically negligible (< 1\% of total runtime).

\paragraph{Recommendation.} Use Method 2 unless you have a specific reason to do otherwise.

% ------------------------------------------------------------
\subsection{Surface and Bottom Constraints (\texttt{mxlSurfFlag})}
\label{sec:mxl_surface_constraints}
% ------------------------------------------------------------

In addition to the three main methods described above, the GGL90 scheme provides an optional surface constraint controlled by the runtime parameter \verb|mxlSurfFlag|.

\subsubsection{Purpose of \texttt{mxlSurfFlag}}

In some configurations, the Gaspar mixing length (Eq.~\ref{eq:mxl_gaspar}) may become very small near the surface due to weak TKE or strong stratification just below the surface. This can lead to unrealistically shallow mixed layers, particularly in coarse-resolution models where the first grid cell is thick (e.g., $\Delta z_1 > 10$ m).

Setting \verb|mxlSurfFlag = .TRUE.| forces the mixing length at the surface interface (between cells 1 and 2) to be at least the thickness of the first grid cell:

\begin{equation}
\ell(k=2) \geq \Delta z_1
\end{equation}

This ensures that the top two cells are always coupled by mixing, preventing the formation of an artificially thin "skin layer" at the surface.

\subsubsection{Implementation}

\begin{lstlisting}[language=Fortran,caption={(ggl90\_mixinglength.F:L147--160) Surface constraint}]
C-    Ensure mixing between first and second level
      IF (mxlSurfFlag) THEN
       DO j=jMin,jMax
        DO i=iMin,iMax
#ifdef ALLOW_SHELFICE
         IF ( useShelfIce ) THEN
          kSrf = MAX(1,kTopC(i,j,bi,bj))
          kTop = MIN(kSrf+1,Nr)
         ENDIF
#endif
         GGL90mixingLength(i,j,kTop)=drF(kSrf)*recip_coordFac
        ENDDO
       ENDDO
      ENDIF
\end{lstlisting}

\paragraph{Key details:}
\begin{itemize}
\item \verb|kSrf| is the surface level index (1 in z-coords, $N_r$ in p-coords).
\item \verb|kTop| is the first interior interface index ($k=2$ in z-coords, $k=N_r$ in p-coords).
\item \verb|drF(kSrf)| is the thickness of the surface cell.
\item The \verb|#ifdef ALLOW_SHELFICE| block adjusts \verb|kSrf| and \verb|kTop| for ice-shelf cavities, where the "surface" may be at a deeper level than $k=1$.
\end{itemize}

This constraint is applied \emph{before} Methods 0--2, so the subsequent geometric constraints can further reduce the mixing length if needed, but they cannot increase it below $\Delta z_1$ at $k=2$.

\subsubsection{When to use \texttt{mxlSurfFlag}}

\paragraph{Enable \texttt{mxlSurfFlag = .TRUE.} when:}
\begin{itemize}
\item The first grid cell is thick ($\Delta z_1 > 10$ m) and you observe unrealistically shallow mixed layers.
\item You are using a coarse vertical resolution in the upper ocean (e.g., typical global configurations, see Appendix~\ref{app:eccov4_config}).
\item You want to ensure that wind-driven mixing always penetrates at least one grid cell below the surface.
\end{itemize}

\paragraph{Disable \texttt{mxlSurfFlag = .FALSE.} when:}
\begin{itemize}
\item You have very high vertical resolution near the surface ($\Delta z_1 < 5$ m) and want to resolve thin diurnal warm layers or cool skin effects.
\item You are performing sensitivity studies where the mixing length should be determined purely by the Gaspar formula and boundary constraints, without artificial floors.
\end{itemize}

\paragraph{Example application.} Global ocean applications with coarse near-surface resolution (typically $\Delta z_1 = 10$ m) often use \verb|mxlSurfFlag = .TRUE.| to ensure robust surface mixing (see Appendix~\ref{app:eccov4_config} for a specific example).

\subsubsection{Bottom boundary condition}

At the bottom interface ($k = N_r$ in z-coords, $k = 1$ in p-coords), the mixing length is implicitly limited by the special treatment in the upward sweep (Method 2) or by the distance-to-bottom constraint (Method 1). No additional parameter analogous to \verb|mxlSurfFlag| exists for the bottom boundary.

However, the parameter \verb|GGL90mixingLengthMin| effectively sets a minimum mixing length everywhere, including at the bottom, to prevent numerical issues when $\ell \to 0$.

\subsubsection{Regularization: \texttt{GGL90\_REGULARIZE\_MIXINGLENGTH}}

An optional compile-time flag \verb|GGL90_REGULARIZE_MIXINGLENGTH| modifies the final mixing length calculation to prevent abrupt transitions. When enabled, the code replaces:

\begin{equation}
\ell_{\text{final}} = \max(\ell, \ell_{\text{min}})
\end{equation}

with:

\begin{equation}
\ell_{\text{final}} = \sqrt{ \ell^2 + \ell_{\text{min}}^2 }
\end{equation}

This provides a smoother transition as $\ell \to 0$, which can improve convergence in adjoint calculations (see Appendix~\ref{app:adjoint_considerations}). The code appears in lines 390--416 of \verb|ggl90_mixinglength.F|:

\begin{lstlisting}[language=Fortran,caption={(ggl90\_mixinglength.F:L402--415) Final mixing length with optional regularization}]
       DO k=2,Nr
        DO j=jMin,jMax
         DO i=iMin,iMax
#ifdef GGL90_REGULARIZE_MIXINGLENGTH
          MLtmp = SQRT( GGL90mixingLength(i,j,k)**2
     &         + GGL90mixingLengthMin**2 )
#else
          MLtmp = MAX(GGL90mixingLength(i,j,k),GGL90mixingLengthMin)
#endif
          GGL90mixingLength(i,j,k) = MLtmp
          rMixingLength(i,j,k) = 1. _d 0 / MLtmp
         ENDDO
        ENDDO
       ENDDO
\end{lstlisting}

\paragraph{When to use regularization:}
\begin{itemize}
\item Enable for adjoint/gradient-based optimization (see Appendix~\ref{app:adjoint_considerations}), where smooth derivatives are important.
\item Disable for forward-only simulations to maintain strict adherence to the published GGL90 formulation.
\end{itemize}

% ------------------------------------------------------------
\subsection{Langmuir Circulation Enhancement (Optional)}
\label{sec:mxl_langmuir}
% ------------------------------------------------------------

If the \verb|ALLOW_GGL90_LANGMUIR| compile-time option is enabled, the mixing length can be further modified to account for Langmuir circulation effects (wind-wave interactions). This is an advanced feature documented separately in Section~\ref{sec:langmuir} and is \emph{not} part of the standard GGL90 scheme.

The Langmuir enhancement computes a modified mixing length \verb|LCmixingLength| that is larger than the standard \verb|GGL90mixingLength| in the surface boundary layer, enhancing vertical mixing when wind and surface waves align. The details are beyond the scope of this section, but users should be aware that the mixing length values may differ from the standard GGL90 formulation when this option is active.

% ============================================================
\section{TKE Evolution and Budget}
\label{sec:tke_evolution}
% ============================================================

The turbulent kinetic energy (TKE) $e$ is the sole prognostic variable in the GGL90 scheme. Its evolution is governed by a budget equation that balances production (from shear and surface forcing), buoyancy effects (destruction by stable stratification, or enhancement by convection), dissipation (conversion to heat at small scales), and diffusive transport. This section documents the TKE equation as implemented in MITgcm, including the time-splitting between explicit and implicit terms.

% ------------------------------------------------------------
\subsection{Governing Equation}
\label{sec:tke_governing}
% ------------------------------------------------------------

The TKE evolution equation implemented in GGL90 is:

\begin{equation}
\frac{\partial e}{\partial t} = \underbrace{\kappa_m \, S^2}_{\text{shear production}}
- \underbrace{\kappa_h \, N^2}_{\text{buoyancy}}
- \underbrace{c_\varepsilon \frac{e^{3/2}}{\ell}}_{\text{dissipation}}
+ \underbrace{\frac{\partial}{\partial z} \left( \kappa_e \frac{\partial e}{\partial z} \right)}_{\text{vertical diffusion}}
+ \underbrace{F_{\text{ext}}}_{\text{external sources}}
\label{eq:tke_budget}
\end{equation}

where:
\begin{itemize}
\item $e$: TKE per unit mass (m$^2$ s$^{-2}$)
\item $\kappa_m$: Eddy viscosity (m$^2$ s$^{-1}$), computed as $\kappa_m = c_\mu \ell \sqrt{e}$ (Section~\ref{sec:eddy_coeffs})
\item $S^2$: Squared vertical shear of horizontal velocity, $S^2 = ({\partial u}/{\partial z})^2 + ({\partial v}/{\partial z})^2$ (s$^{-2}$)
\item $\kappa_h$: Eddy diffusivity for scalars (m$^2$ s$^{-1}$), related to $\kappa_m$ via the turbulent Prandtl number (Section~\ref{sec:eddy_coeffs})
\item $N^2$: Squared buoyancy frequency (s$^{-2}$), $N^2 = -(g/\rho_0) \, \partial\rho/\partial z$
\item $c_\varepsilon$: Dissipation constant (dimensionless), typically 0.7
\item $\ell$: Mixing length (m), computed by \verb|GGL90_MIXINGLENGTH| (Section~\ref{sec:mxl_method2})
\item $\kappa_e$: Diffusivity of TKE itself (m$^2$ s$^{-1}$), $\kappa_e = \alpha_e \kappa_m$ with $\alpha_e \approx 1$
\item $F_{\text{ext}}$: External TKE sources (optional), including IDEMIX internal wave dissipation (Section~\ref{sec:idemix}) and Langmuir circulation (Section~\ref{sec:langmuir})
\end{itemize}

\paragraph{Physical interpretation of each term:}
\begin{enumerate}
\item \textbf{Shear production ($\kappa_m S^2$):} Extraction of kinetic energy from the mean flow by turbulent Reynolds stresses. Always positive (source term).

\item \textbf{Buoyancy ($-\kappa_h N^2$):}
\begin{itemize}
\item If $N^2 > 0$ (stable stratification): Buoyancy work is negative (sink term), as turbulent eddies must do work against gravity to mix denser fluid upward.
\item If $N^2 < 0$ (convective instability): Buoyancy work is positive (source term), as potential energy is released by gravitational overturning.
\end{itemize}

\item \textbf{Dissipation ($c_\varepsilon e^{3/2}/\ell$):} Irreversible conversion of TKE to heat via molecular viscosity at the smallest turbulent scales. Always positive (sink term). Note that the dissipation rate scales as $e^{3/2}$ (Kolmogorov's $-5/3$ law) and is inversely proportional to the mixing length $\ell$.

\item \textbf{Vertical diffusion:} Redistribution of TKE in the vertical by turbulent transport. This is typically a down-gradient flux (from high TKE to low TKE regions).

\item \textbf{External sources ($F_{\text{ext}}$):} Additional TKE sources from parameterized processes not explicitly resolved in the shear production term. Examples include internal wave breaking (IDEMIX) and Langmuir circulation.
\end{enumerate}

\paragraph{Coordinate-invariant form.} The equation above is written in z-coordinates (depth increasing downward, $z$ positive upward). In pressure coordinates, vertical derivatives $\partial/\partial z$ are replaced by $\partial/\partial p$ with appropriate scaling factors (\verb|coordFac| in the code).

\paragraph{Boundary conditions.} Equation~\eqref{eq:tke_budget} is subject to:
\begin{itemize}
\item \textbf{Surface}: Dirichlet condition $e = B_0 u_*^2$ at $z=0$, where $u_*$ is the friction velocity and $B_0 \approx 3.75$ (Section~\ref{sec:bc_surface}).
\item \textbf{Bottom}: Either Dirichlet ($e = B_0 u_{*,\text{bottom}}^2$) or Neumann ($\partial e/\partial z = 0$), controlled by \verb|GGL90_dirichlet| (Section~\ref{sec:bc_bottom}).
\end{itemize}

% ------------------------------------------------------------
\subsection{Production Term}
\label{sec:tke_production}
% ------------------------------------------------------------

The shear production term $P = \kappa_m S^2$ represents the extraction of kinetic energy from the mean flow by turbulent Reynolds stresses.

\subsubsection{Vertical shear computation}

The squared vertical shear $S^2$ is computed in \verb|GGL90_CALC| immediately after the mixing length calculation. MITgcm provides two options for computing $S^2$, controlled by the runtime parameter \verb|calcMeanVertShear|:

\paragraph{Option 1: Average of four components (\texttt{calcMeanVertShear = .TRUE.})}

Compute the shear at each of the four surrounding velocity points (two U-points, two V-points), then average:

\begin{equation}
S^2(i,j,k) = \frac{1}{2} \left[
\left(\frac{\Delta u_1}{\Delta z}\right)^2 + \left(\frac{\Delta u_2}{\Delta z}\right)^2
+ \left(\frac{\Delta v_1}{\Delta z}\right)^2 + \left(\frac{\Delta v_2}{\Delta z}\right)^2
\right]
\label{eq:shear_option1}
\end{equation}

where $\Delta u_1 = u(i,j,k-1) - u(i,j,k)$, $\Delta u_2 = u(i+1,j,k-1) - u(i+1,j,k)$, and similarly for $v$.

\begin{lstlisting}[language=Fortran,caption={(ggl90\_calc.F:L468--482) Vertical shear: Option 1}]
       IF ( calcMeanVertShear ) THEN
C     by averaging (@ grid-cell center) the 4 vertical shear compon @ U,V pos.
        DO j=jMin,jMax
         DO i=iMin,iMax
          tempU  = ( uVel( i ,j,km1,bi,bj) - uVel( i ,j,k,bi,bj) )
          tempUp = ( uVel(i+1,j,km1,bi,bj) - uVel(i+1,j,k,bi,bj) )
          tempV  = ( vVel(i, j ,km1,bi,bj) - vVel(i, j ,k,bi,bj) )
          tempVp = ( vVel(i,j+1,km1,bi,bj) - vVel(i,j+1,k,bi,bj) )
          verticalShear(i,j) = (
     &                 ( tempU*tempU + tempUp*tempUp )
     &               + ( tempV*tempV + tempVp*tempVp )
     &                         )*halfRL*recip_drC(k)*recip_drC(k)
     &                          *coordFac*coordFac
         ENDDO
        ENDDO
\end{lstlisting}

\paragraph{Option 2: Shear of averaged flow (\texttt{calcMeanVertShear = .FALSE.})}

First average the velocity components to the cell center, then compute shear:

\begin{equation}
S^2(i,j,k) = \left(\frac{\bar{u}(k-1) - \bar{u}(k)}{\Delta z}\right)^2
+ \left(\frac{\bar{v}(k-1) - \bar{v}(k)}{\Delta z}\right)^2
\label{eq:shear_option2}
\end{equation}

where $\bar{u}(k) = \frac{1}{2}[u(i,j,k) + u(i+1,j,k)]$ and $\bar{v}(k) = \frac{1}{2}[v(i,j,k) + v(i,j+1,k)]$.

\begin{lstlisting}[language=Fortran,caption={(ggl90\_calc.F:L483--498) Vertical shear: Option 2}]
       ELSE
C     from the averaged flow at grid-cell center (2 compon x 2 pos.)
        DO j=jMin,jMax
         DO i=iMin,iMax
          tempU = ( ( uVel(i,j,km1,bi,bj) + uVel(i+1,j,km1,bi,bj) )
     &             -( uVel(i,j,k  ,bi,bj) + uVel(i+1,j,k  ,bi,bj) )
     &            )*halfRL*recip_drC(k)
     &             *coordFac
          tempV = ( ( vVel(i,j,km1,bi,bj) + vVel(i,j+1,km1,bi,bj) )
     &             -( vVel(i,j,k  ,bi,bj) + vVel(i,j+1,k  ,bi,bj) )
     &            )*halfRL*recip_drC(k)
     &             *coordFac
          verticalShear(i,j) = tempU*tempU + tempV*tempV
         ENDDO
        ENDDO
       ENDIF
\end{lstlisting}

\paragraph{Which option to use?} Option 1 is more accurate on the Arakawa C-grid, as it avoids spurious cancellation of checkerboard modes. Option 2 is slightly cheaper computationally. The difference is typically small in well-resolved simulations. Global applications typically use Option 1 (\verb|calcMeanVertShear = .TRUE.|, see Appendix~\ref{app:eccov4_config}).

\subsubsection{Production term calculation}

Once $S^2$ and $\kappa_m$ are known, the production term is straightforward:

\begin{lstlisting}[language=Fortran,caption={(ggl90\_calc.F:L605--610) Explicit TKE budget update: production term}]
         GGL90TKE(i,j,k,bi,bj) = GGL90TKE(i,j,k,bi,bj)
     &        + deltaTloc*(
     &        + KappaM(i,j)*verticalShear(i,j)
     &        - KappaH*Nsquare(i,j,k)
     &        - TKEdissipation
     &        )
\end{lstlisting}

The production term \verb|KappaM(i,j)*verticalShear(i,j)| is always non-negative, since both $\kappa_m$ and $S^2$ are non-negative.

\paragraph{Typical values.} In a wind-driven surface mixed layer with $\kappa_m \sim 0.01$ m$^2$ s$^{-1}$ and $S^2 \sim 10^{-4}$ s$^{-2}$, the production term is $P \sim 10^{-6}$ m$^2$ s$^{-3}$. This is comparable to the dissipation term (see below), indicating a balance between shear production and dissipation in quasi-steady turbulence.

% ------------------------------------------------------------
\subsection{Buoyancy Term}
\label{sec:tke_buoyancy}
% ------------------------------------------------------------

The buoyancy term $B = -\kappa_h N^2$ represents the work done by buoyancy forces on turbulent eddies.

\subsubsection{Sign convention and physical interpretation}

The sign of the buoyancy term depends on the stratification:

\begin{itemize}
\item \textbf{Stable stratification ($N^2 > 0$):} The term $-\kappa_h N^2 < 0$ is a sink of TKE. Turbulent eddies must work against buoyancy to lift heavy fluid (or depress light fluid), converting kinetic energy to potential energy. This suppresses turbulence.

\item \textbf{Neutral stratification ($N^2 \approx 0$):} The buoyancy term is negligible, and TKE evolution is controlled by shear production and dissipation.

\item \textbf{Unstable stratification ($N^2 < 0$):} The term $-\kappa_h N^2 > 0$ is a source of TKE. Convective overturning releases potential energy, which is converted to kinetic energy. This enhances turbulence.
\end{itemize}

\paragraph{Relation to the Richardson number.} The ratio of buoyancy to production defines the flux Richardson number:

\begin{equation}
Ri_f = \frac{-\kappa_h N^2}{\kappa_m S^2} = -\frac{N^2}{Pr_t \, S^2}
\end{equation}

where $Pr_t = \kappa_m / \kappa_h$ is the turbulent Prandtl number (Section~\ref{sec:tke_prandtl}). A critical value $Ri_f \approx 0.2$ separates turbulent ($Ri_f < 0.2$) from laminar ($Ri_f > 0.2$) regimes.

\subsubsection{Implementation}

The buoyancy term is computed alongside the production term in the explicit TKE update:

\begin{lstlisting}[language=Fortran,caption={(ggl90\_calc.F:L605--610) Explicit TKE budget update: buoyancy term}]
         GGL90TKE(i,j,k,bi,bj) = GGL90TKE(i,j,k,bi,bj)
     &        + deltaTloc*(
     &        + KappaM(i,j)*verticalShear(i,j)
     &        - KappaH*Nsquare(i,j,k)
     &        - TKEdissipation
     &        )
\end{lstlisting}

The term \verb|- KappaH*Nsquare(i,j,k)| directly implements $-\kappa_h N^2$.

\paragraph{Coordinate scaling.} Recall that \verb|Nsquare| is pre-multiplied by \verb|coordFac| (line 347--348 of \verb|ggl90_calc.F|) to ensure dimensional consistency in pressure coordinates. This is automatically accounted for in the buoyancy term.

% ------------------------------------------------------------
\subsection{Dissipation Term}
\label{sec:tke_dissipation}
% ------------------------------------------------------------

The dissipation term $\varepsilon = c_\varepsilon e^{3/2} / \ell$ represents the irreversible conversion of TKE to heat at the smallest scales of turbulence (the Kolmogorov scale).

\subsubsection{Physical basis}

In fully developed turbulence, the TKE dissipation rate $\varepsilon$ scales as:

\begin{equation}
\varepsilon \sim \frac{e^{3/2}}{\ell}
\label{eq:dissipation_scaling}
\end{equation}

This scaling arises from dimensional analysis: $\varepsilon$ has units of m$^2$ s$^{-3}$, so the only combination of TKE ($e$, units m$^2$ s$^{-2}$) and length scale ($\ell$, units m) that yields the correct dimensions is $e^{3/2}/\ell$.

The proportionality constant $c_\varepsilon$ is an empirical parameter, typically set to 0.7 based on atmospheric boundary layer observations and large-eddy simulations. In MITgcm, $c_\varepsilon$ is the runtime parameter \verb|GGL90ceps|.

\subsubsection{Implementation: Semi-implicit time integration}

The dissipation term is unique among the TKE budget terms in that it is treated \emph{semi-implicitly}: a fraction is computed explicitly, and the remainder is treated implicitly to ensure numerical stability.

\paragraph{Explicit part:}

\begin{lstlisting}[language=Fortran,caption={(ggl90\_calc.F:L600--610) Explicit dissipation term}]
C     dissipation term
         TKEdissipation = explDissFac*GGL90ceps
     &        *SQRTTKE(i,j,k)*rMixingLength(i,j,k)
     &        *GGL90TKE(i,j,k,bi,bj)
C     partial update with sum of explicit contributions
         GGL90TKE(i,j,k,bi,bj) = GGL90TKE(i,j,k,bi,bj)
     &        + deltaTloc*(
     &        + KappaM(i,j)*verticalShear(i,j)
     &        - KappaH*Nsquare(i,j,k)
     &        - TKEdissipation
     &        )
\end{lstlisting}

Here, \verb|explDissFac| is the explicit fraction of the dissipation term. The formula for \verb|TKEdissipation| is:

\begin{equation}
\varepsilon_{\text{explicit}} = f_{\text{expl}} \, c_\varepsilon \, \frac{\sqrt{e} \cdot e}{\ell} = f_{\text{expl}} \, c_\varepsilon \, \frac{e^{3/2}}{\ell}
\end{equation}

where $f_{\text{expl}}$ = \verb|explDissFac| and \verb|rMixingLength = 1/ell| is precomputed to avoid repeated divisions.

\paragraph{Implicit part:}

The implicit fraction $(1 - f_{\text{expl}})$ appears in the center diagonal of the tridiagonal matrix used for the implicit TKE solve:

\begin{lstlisting}[language=Fortran,caption={(ggl90\_calc.F:L746--749) Implicit dissipation in matrix diagonal}]
         b3d(i,j,k) = 1. _d 0 - c3d(i,j,k) - a3d(i,j,k)
     &        + implDissFac*deltaTloc*GGL90ceps*SQRTTKE(i,j,k)
     &        * rMixingLength(i,j,k)
     &        * maskC(i,j,k,bi,bj)*maskC(i,j,km1,bi,bj)
\end{lstlisting}

The term \verb|+ implDissFac*deltaTloc*GGL90ceps*SQRTTKE(i,j,k)*rMixingLength(i,j,k)| adds the implicit dissipation contribution to the diagonal of the matrix.

\subsubsection{Explicit/implicit splitting}

The fractions \verb|explDissFac| and \verb|implDissFac| are set in lines 226--246 of \verb|ggl90_calc.F|:

\begin{lstlisting}[language=Fortran,caption={(ggl90\_calc.F:L226--246) Explicit/implicit dissipation fractions}]
C     Explicit/implicit TKE dissipation
      explDissFac = 0. _d 0
      implDissFac = 1. _d 0
      IF ( GGL90dumpFreq.LT.dTtracerLev(1) ) THEN
C      Recover the old explicit-only-formulation by
C      setting explDissFac=1 and implDissFac=0
       explDissFac = 1. _d 0
       implDissFac = 0. _d 0
      ENDIF
C     The implicit part is not coded very efficiently in terms
C     of TAF dependencies (with the CG3D routine called by the
C     GGL90_DO_EXCH macro), so we default to DIVA_EXCH instead,
C     and still use explDissFac=0 and implDissFac=1
#ifdef ALLOW_AUTODIFF_TAMC
      IF ( useGGL90inAdMode ) THEN
       explDissFac = 0. _d 0
       implDissFac = 1. _d 0
      ENDIF
#endif /* ALLOW_AUTODIFF_TAMC */
\end{lstlisting}

\paragraph{Default behavior:}
\begin{itemize}
\item \textbf{Standard forward run}: \verb|explDissFac = 0|, \verb|implDissFac = 1| (fully implicit dissipation). This is unconditionally stable for any time step.
\item \textbf{Adjoint mode}: Same as forward (fully implicit).
\item \textbf{Legacy mode}: If \verb|GGL90dumpFreq < dTtracerLev(1)| (which is never true in practice), revert to the old fully explicit scheme (\verb|explDissFac = 1|, \verb|implDissFac = 0|). This is retained for backward compatibility but is \textbf{not recommended} due to potential instability.
\end{itemize}

\paragraph{Why semi-implicit?} The dissipation term is nonlinear in $e$ (proportional to $e^{3/2}$), which can cause CFL-like instabilities if treated purely explicitly with large time steps. Implicit treatment linearizes the term about the current time step ($\partial \varepsilon / \partial e \approx 3/2 \cdot \varepsilon / e$), ensuring stability without requiring prohibitively small time steps.

% ------------------------------------------------------------
\subsection{Vertical Diffusion of TKE}
\label{sec:tke_diffusion}
% ------------------------------------------------------------

The vertical diffusion term $\partial/\partial z (\kappa_e \, \partial e / \partial z)$ represents turbulent transport of TKE itself, spreading TKE from high-energy regions (e.g., near the surface) to low-energy regions (e.g., the thermocline).

\subsubsection{Diffusivity of TKE: $\kappa_e$}

The diffusivity of TKE is related to the eddy viscosity $\kappa_m$ via the constant $\alpha_e$:

\begin{equation}
\kappa_e = \alpha_e \, \kappa_m
\label{eq:kappa_e}
\end{equation}

In GGL90, $\alpha_e$ = \verb|GGL90alpha| is a runtime parameter, typically set to 1.0 (following Gaspar et al.\ 1990).

The code computes $\kappa_e$ in lines 593--598 of \verb|ggl90_calc.F|:

\begin{lstlisting}[language=Fortran,caption={(ggl90\_calc.F:L593--598) Diffusivity of TKE}]
       DO j=jMin,jMax
        DO i=iMin,iMax
C        diffusivity
         KappaH = KappaM(i,j)/TKEPrandtlNumber(i,j,k)
         KappaE(i,j,k) = GGL90alpha * KappaM(i,j)
     &        * maskC(i,j,k,bi,bj)*maskC(i,j,k-1,bi,bj)
\end{lstlisting}

Note that $\kappa_e$ is stored at cell interfaces (W-points), just like $\kappa_m$ and $\kappa_h$.

\subsubsection{Implicit vertical diffusion solver}

The vertical diffusion of TKE is treated \emph{fully implicitly} using a tridiagonal solve. This is necessary for numerical stability, as explicit diffusion would impose a restrictive CFL condition:

\begin{equation}
\Delta t \leq \frac{\Delta z^2}{2 \kappa_e}
\end{equation}

which would be prohibitively small in regions with large $\kappa_e$ (e.g., near the surface under strong wind forcing).

\paragraph{Tridiagonal matrix construction.} The implicit diffusion equation is discretized as:

\begin{equation}
\frac{e^{n+1}_k - e^*_k}{\Delta t} = \frac{1}{\Delta z_k} \left[
\kappa_{e,k-1/2} \frac{e^{n+1}_{k-1} - e^{n+1}_k}{\Delta z_{k-1/2}}
- \kappa_{e,k+1/2} \frac{e^{n+1}_k - e^{n+1}_{k+1}}{\Delta z_{k+1/2}}
\right]
\label{eq:tke_implicit_diffusion}
\end{equation}

where $e^*_k$ is the TKE after the explicit update (production, buoyancy, dissipation), and $e^{n+1}_k$ is the final TKE after implicit diffusion. Rearranging yields a tridiagonal system:

\begin{equation}
a_k e^{n+1}_{k-1} + b_k e^{n+1}_k + c_k e^{n+1}_{k+1} = e^*_k
\label{eq:tke_tridiagonal}
\end{equation}

The coefficients $a_k$, $b_k$, $c_k$ are computed in lines 668--760 of \verb|ggl90_calc.F|:

\begin{lstlisting}[language=Fortran,caption={(ggl90\_calc.F:L668--720) Tridiagonal matrix setup for implicit TKE diffusion}]
C     ==============================================
C     Implicit time step to update TKE for k=1,Nr;
C     TKE(Nr+1)=0 by default;
C     for pressure coordinates, this translates into
C     TKE(1)  = 0, TKE(Nr+1) is the surface value
C     ==============================================
C     set up matrix
C--   Lower diagonal
      DO j=jMin,jMax
       DO i=iMin,iMax
         a3d(i,j,1) = 0. _d 0
       ENDDO
      ENDDO
      DO k=2,Nr
       km1=MAX(2,k-1)
       DO j=jMin,jMax
        DO i=iMin,iMax
         IF ( usingPCoords) km1=MIN(Nr,MAX(kSurfC(i,j,bi,bj)+1,k-1))
C-    We keep recip_hFacC in the diffusive flux calculation,
C-    but no hFacC in TKE volume control
C-    No need for maskC(k-1) with recip_hFacC(k-1)
         a3d(i,j,k) = -deltaTloc
     &        *recip_drF(k-1)*recip_hFacC(i,j,k-1,bi,bj)
     &        *.5 _d 0*(KappaE(i,j, k )+KappaE(i,j,km1))
     &        *recip_drC(k)*maskC(i,j,k,bi,bj)*recip_hFacI(i,j,k)
     &        *coordFac*coordFac
        ENDDO
       ENDDO
      ENDDO
C--   Upper diagonal
      DO j=jMin,jMax
       DO i=iMin,iMax
         c3d(i,j,1)  = 0. _d 0
       ENDDO
      ENDDO
      DO k=2,Nr
       kp1=MIN(k+1,Nr)
       DO j=jMin,jMax
        DO i=iMin,iMax
         IF ( usingZCoords ) kp1=MAX(1,MIN(klowC(i,j,bi,bj),k+1))
C-    We keep recip_hFacC in the diffusive flux calculation,
C-    but no hFacC in TKE volume control
C-    No need for maskC(k) with recip_hFacC(k)
         c3d(i,j,k) = -deltaTloc
     &        *recip_drF( k ) * recip_hFacC(i,j,k,bi,bj)
     &        *.5 _d 0*(KappaE(i,j,k)+KappaE(i,j,kp1))
     &        *recip_drC(k)*maskC(i,j,k-1,bi,bj)*recip_hFacI(i,j,k)
     &        *coordFac*coordFac
        ENDDO
       ENDDO
      ENDDO
\end{lstlisting}

\paragraph{Key implementation details:}
\begin{itemize}
\item \verb|a3d(i,j,k)| is the lower diagonal coefficient (coupling to cell $k-1$).
\item \verb|c3d(i,j,k)| is the upper diagonal coefficient (coupling to cell $k+1$).
\item \verb|b3d(i,j,k)| is the center diagonal coefficient (cell $k$), computed as $b_k = 1 - a_k - c_k + \text{(implicit dissipation)}$.
\item The factor \verb|0.5*(KappaE(i,j,k)+KappaE(i,j,km1))| is the average of $\kappa_e$ at the two adjacent interfaces, which is appropriate for a cell-centered TKE variable with interface-located diffusivities.
\item \verb|coordFac*coordFac| accounts for coordinate transformation (z vs.\ p coords).
\end{itemize}

\subsubsection{Tridiagonal solve}

After the matrix is constructed, the system is solved by a call to the generic MITgcm tridiagonal solver \verb|SOLVE_TRIDIAGONAL|:

\begin{lstlisting}[language=Fortran,caption={(ggl90\_calc.F:L900--903) Solve tridiagonal system for TKE}]
      CALL SOLVE_TRIDIAGONAL( iMin, iMax, jMin, jMax,
     I                        a3d, b3d, c3d,
     U                        GGL90TKE(1-OLx,1-OLy,1,bi,bj),
     I                        errCode, bi, bj, myThid )
\end{lstlisting}

On entry, \verb|GGL90TKE| contains the right-hand side $e^*$ (TKE after explicit update). On exit, it contains the solution $e^{n+1}$ (TKE after implicit diffusion and dissipation).

% ------------------------------------------------------------
\subsection{Time Integration Scheme: Operator Splitting}
\label{sec:tke_time_integration}
% ------------------------------------------------------------

The GGL90 TKE equation is integrated using a \textbf{split-explicit-implicit} scheme:

\begin{enumerate}
\item \textbf{Explicit step}: Update TKE with production, buoyancy, explicit fraction of dissipation, and any external sources (IDEMIX, Langmuir, horizontal diffusion):
\begin{equation}
e^* = e^n + \Delta t \left( \kappa_m S^2 - \kappa_h N^2 - f_{\text{expl}} \frac{c_\varepsilon e^{3/2}}{\ell} + F_{\text{ext}} \right)
\end{equation}

\item \textbf{Implicit step}: Solve for $e^{n+1}$ including implicit dissipation and vertical diffusion:
\begin{equation}
e^{n+1} = \mathcal{L}^{-1} \left[ e^* - (1 - f_{\text{expl}}) \frac{c_\varepsilon (e^{n+1})^{3/2}}{\ell} \right]
\end{equation}
where $\mathcal{L}^{-1}$ represents the tridiagonal solve for vertical diffusion.
\end{enumerate}

This splitting is unconditionally stable for any time step $\Delta t$, provided that the explicit dissipation fraction $f_{\text{expl}}$ is small (default: $f_{\text{expl}} = 0$).

\subsubsection{Time step selection}

The TKE time step \verb|deltaTloc| is set equal to the tracer time step \verb|dTtracerLev(1)| (line 221 of \verb|ggl90_calc.F|):

\begin{lstlisting}[language=Fortran,caption={(ggl90\_calc.F:L221) TKE time step}]
      deltaTloc = dTtracerLev(1)
\end{lstlisting}

This ensures that TKE evolves synchronously with temperature and salinity, which is physically appropriate since turbulent mixing of TKE and scalars occurs on the same time scales.

\paragraph{Sub-cycling TKE (optional).} For very high Reynolds number flows or very coarse time stepping, it may be beneficial to sub-cycle the TKE equation (i.e., take multiple TKE steps per tracer step). However, this is not currently implemented in the standard GGL90 package and would require code modifications.

\subsubsection{Horizontal diffusion of TKE (optional)}

If the compile-time flag \verb|ALLOW_GGL90_HORIZDIFF| is enabled and \verb|GGL90diffTKEh > 0|, the code also includes horizontal diffusion of TKE:

\begin{equation}
\frac{\partial e}{\partial t} = \ldots + \nabla_h \cdot \left( \kappa_{h,\text{TKE}} \nabla_h e \right)
\end{equation}

This is computed explicitly (lines 381--434 of \verb|ggl90_calc.F|) and added to the TKE budget before the implicit solve (lines 638--648):

\begin{lstlisting}[language=Fortran,caption={(ggl90\_calc.F:L638--648) Horizontal diffusion of TKE}]
#ifdef ALLOW_GGL90_HORIZDIFF
       IF ( GGL90diffTKEh .GT. 0. _d 0 ) THEN
C--    Add horiz. diffusion tendency
        DO j=jMin,jMax
         DO i=iMin,iMax
          GGL90TKE(i,j,k,bi,bj) = GGL90TKE(i,j,k,bi,bj)
     &                          + gTKE(i,j)*deltaTloc
         ENDDO
        ENDDO
       ENDIF
#endif /* ALLOW_GGL90_HORIZDIFF */
\end{lstlisting}

\paragraph{When to use horizontal diffusion.} Horizontal TKE diffusion is rarely used in realistic ocean simulations, as it is not physically motivated (TKE should be generated locally by shear and buoyancy, not transported horizontally). However, it can be useful for:
\begin{itemize}
\item Smoothing spurious TKE gradients in coarse-resolution models.
\item Numerical stability in eddy-permitting or eddy-resolving configurations with very strong mesoscale variability.
\end{itemize}

The default is \verb|GGL90diffTKEh = 0| (no horizontal diffusion).

% ============================================================
\section{Eddy Coefficient Calculation}
\label{sec:eddy_coeffs}
% ============================================================

The GGL90 scheme diagnoses eddy viscosity $\kappa_m$ and eddy diffusivity $\kappa_h$ from the prognostic TKE field $e$ and the mixing length $\ell$. These coefficients are then used by the main model to compute vertical fluxes of momentum, temperature, and salinity.

The calculation proceeds in three steps:
\begin{enumerate}
\item Compute eddy viscosity $\kappa_m = c_\mu \ell \sqrt{e}$ (Section~\ref{sec:eddy_visc}).
\item Compute TKE Prandtl number $Pr_t = f(Ri)$ from the local Richardson number (Section~\ref{sec:tke_prandtl}).
\item Compute eddy diffusivity $\kappa_h = \kappa_m / Pr_t$ (Section~\ref{sec:eddy_diff}).
\item Apply background floors and maximum caps (Section~\ref{sec:eddy_floors}).
\end{enumerate}

These coefficients are computed at every interior interface (k=2 to Nr) and then transferred to the main model via \verb|GGL90_CALC_VISC| and \verb|GGL90_CALC_DIFF| (Section~\ref{sec:model_interface}).

% ------------------------------------------------------------
\subsection{Eddy Viscosity ($\kappa_m$)}
\label{sec:eddy_visc}
% ------------------------------------------------------------

\subsubsection{Theoretical basis}

The eddy viscosity in GGL90 is based on Prandtl's mixing length hypothesis combined with the TKE energy scale:

\begin{equation}
\kappa_m = c_\mu \, \ell \, \sqrt{e}
\label{eq:kappa_m}
\end{equation}

where:
\begin{itemize}
\item $c_\mu$ (\verb|GGL90ck|): Dimensionless constant, typically $c_\mu = 0.1$ (Gaspar et al.\ 1990).
\item $\ell$ (\verb|GGL90mixingLength|): Mixing length (m), computed by \verb|GGL90_MIXINGLENGTH| (Section~\ref{sec:mxl_method2}).
\item $\sqrt{e}$ (\verb|SQRTTKE|): Characteristic turbulent velocity scale (m s$^{-1}$).
\end{itemize}

\paragraph{Physical interpretation.} The product $\ell \sqrt{e}$ has dimensions of m$^2$ s$^{-1}$ (kinematic viscosity). It represents the rate at which turbulent eddies of size $\ell$ and velocity $\sqrt{e}$ transport momentum vertically. The constant $c_\mu$ is an empirical factor determined from observations and large-eddy simulations.

\paragraph{Comparison to K-profile parameterization (KPP).} In KPP, the eddy viscosity is prescribed as a shape function times a boundary layer depth scale. In GGL90, $\kappa_m$ emerges dynamically from the TKE and mixing length, which respond to local stratification and shear. This makes GGL90 more physically consistent, but also more computationally expensive due to the prognostic TKE equation.

\subsubsection{Implementation}

The eddy viscosity is computed in lines 436--465 of \verb|ggl90_calc.F|:

\begin{lstlisting}[language=Fortran,caption={(ggl90\_calc.F:L446--465) Eddy viscosity calculation}]
        DO j=jMin,jMax
         DO i=iMin,iMax
          KappaM(i,j) = GGL90ck*GGL90mixingLength(i,j,k)*SQRTTKE(i,j,k)
         ENDDO
        ENDDO
       DO j=jMin,jMax
        DO i=iMin,iMax
         GGL90visctmp(i,j,k) = MAX( KappaM(i,j),diffKrNrS(k)
     &                            * recip_coordFac*recip_coordFac )
     &        * maskC(i,j,k,bi,bj)*maskC(i,j,k-1,bi,bj)
C        note: storing GGL90visctmp like this, and using it later to compute
C              GGL9rdiffKr etc. is robust in case of smoothing (e.g. see OPA)
         KappaM(i,j) = MAX( KappaM(i,j),viscArNr(k)
     &                    * recip_coordFac*recip_coordFac )
     &                    * maskC(i,j,k,bi,bj)*maskC(i,j,k-1,bi,bj)
        ENDDO
       ENDDO
\end{lstlisting}

\paragraph{Key implementation details:}
\begin{itemize}
\item Line 448: $\kappa_m$ is computed from Eq.~\eqref{eq:kappa_m}.
\item Lines 456--458: \verb|GGL90visctmp| is a copy of $\kappa_m$ (before applying the viscosity floor), floored at the background scalar diffusivity \verb|diffKrNrS(k)|. This is used later to compute diffusivities for temperature and salinity, ensuring that scalar diffusion is never less than the background value.
\item Lines 461--463: The actual \verb|KappaM| used for momentum is floored at the background viscosity \verb|viscArNr(k)|, which is typically larger than \verb|diffKrNrS(k)| (see Section~\ref{sec:eddy_floors}).
\item \verb|maskC(i,j,k)*maskC(i,j,k-1)|: Ensures that $\kappa_m$ is zero at dry interfaces (e.g., where one of the adjacent cells is land or ice).
\item \verb|recip_coordFac*recip_coordFac|: Coordinate transformation factor for pressure coordinates. In z-coordinates, \verb|coordFac = 1|, so this factor is unity.
\end{itemize}

\subsubsection{Typical values}

For representative ocean conditions:
\begin{itemize}
\item \textbf{Surface mixed layer} (wind-driven turbulence): $\ell \sim 20$ m, $\sqrt{e} \sim 0.03$ m s$^{-1}$, $c_\mu = 0.1$ $\Rightarrow$ $\kappa_m \sim 0.06$ m$^2$ s$^{-1}$.
\item \textbf{Thermocline} (strong stratification): $\ell \sim 1$ m, $\sqrt{e} \sim 0.003$ m s$^{-1}$, $c_\mu = 0.1$ $\Rightarrow$ $\kappa_m \sim 3 \times 10^{-4}$ m$^2$ s$^{-1}$.
\item \textbf{Deep ocean} (weak background turbulence): $\ell \sim 5$ m, $\sqrt{e} \sim 0.001$ m s$^{-1}$, $c_\mu = 0.1$ $\Rightarrow$ $\kappa_m \sim 5 \times 10^{-4}$ m$^2$ s$^{-1}$.
\end{itemize}

These values span 2--3 orders of magnitude, illustrating the highly variable nature of ocean turbulence.

% ------------------------------------------------------------
\subsection{TKE Prandtl Number}
\label{sec:tke_prandtl}
% ------------------------------------------------------------

The turbulent Prandtl number $Pr_t$ relates the eddy viscosity $\kappa_m$ to the eddy diffusivity for scalars $\kappa_h$:

\begin{equation}
Pr_t = \frac{\kappa_m}{\kappa_h}
\label{eq:prandtl_def}
\end{equation}

In a molecularly-dominated flow, the Prandtl number is determined by fluid properties (e.g., $Pr \approx 7$ for seawater at 20°C). In turbulent flow, $Pr_t$ depends on the flow stability, quantified by the Richardson number.

\subsubsection{Formulation as a function of Richardson number}

GGL90 uses an empirical relationship between $Pr_t$ and the local gradient Richardson number $Ri$:

\begin{equation}
Ri = \frac{N^2}{S^2} = \frac{-(g/\rho_0) \, \partial\rho/\partial z}{(\partial u/\partial z)^2 + (\partial v/\partial z)^2}
\label{eq:richardson_number}
\end{equation}

The Prandtl number formulation is:

\begin{equation}
Pr_t = \begin{cases}
1 & \text{if } Ri < 0.2 \\
5 \, Ri & \text{if } Ri \geq 0.2
\end{cases}
\label{eq:prandtl_ri}
\end{equation}

with an upper cap $Pr_t \leq 10$ to prevent excessively small scalar diffusivities.

\paragraph{Physical interpretation:}
\begin{itemize}
\item \textbf{Strongly turbulent regime ($Ri < 0.2$):} Turbulence is vigorous enough to mix momentum and scalars equally efficiently, so $Pr_t = 1$ (i.e., $\kappa_m = \kappa_h$).
\item \textbf{Weakly turbulent / transitional regime ($Ri \geq 0.2$):} Stratification suppresses vertical motion, making scalar mixing less efficient than momentum mixing. The ratio $Pr_t$ increases linearly with $Ri$, reflecting the increasing dominance of buoyancy over shear.
\item \textbf{Cap at $Pr_t = 10$:} This prevents $\kappa_h$ from becoming unrealistically small in strongly stratified regions where $Ri$ may be very large (e.g., $Ri > 10$ in the deep thermocline). The cap ensures a minimum level of scalar diffusion.
\end{itemize}

The critical Richardson number $Ri_c = 0.2$ is a classical result from atmospheric boundary layer theory, corresponding to the threshold for turbulence suppression by stratification.

\subsubsection{Implementation}

The TKE Prandtl number is computed in lines 561--588 of \verb|ggl90_calc.F|:

\begin{lstlisting}[language=Fortran,caption={(ggl90\_calc.F:L577--585) TKE Prandtl number}]
        DO j=jMin,jMax
         DO i=iMin,iMax
          RiNumber = MAX(Nsquare(i,j,k),0. _d 0)
     &         /(verticalShear(i,j)+GGL90eps)
          prTemp = 1. _d 0
          IF ( RiNumber .GE. 0.2 _d 0 ) prTemp = 5. _d 0 * RiNumber
          TKEPrandtlNumber(i,j,k) = MIN(10. _d 0,prTemp)
         ENDDO
        ENDDO
\end{lstlisting}

\paragraph{Key implementation details:}
\begin{itemize}
\item Line 579: $Ri = \max(N^2, 0) / (S^2 + \varepsilon)$. The numerator is floored at zero to exclude convectively unstable regions ($N^2 < 0$), where the concept of a "Richardson number" is not well-defined. The denominator is floored at \verb|GGL90eps| to avoid division by zero when $S^2 \approx 0$.
\item Lines 581--582: Implement the piecewise formula \eqref{eq:prandtl_ri}.
\item Line 583: Cap at $Pr_t = 10$.
\end{itemize}

\paragraph{Behavior in limiting cases:}
\begin{itemize}
\item \textbf{Convective layer} ($N^2 < 0$): $Ri = 0$ $\Rightarrow$ $Pr_t = 1$.
\item \textbf{Neutral layer} ($N^2 \approx 0$, $S^2 > 0$): $Ri \approx 0$ $\Rightarrow$ $Pr_t = 1$.
\item \textbf{Strongly stratified layer} ($N^2 \gg S^2$): $Ri \gg 1$ $\Rightarrow$ $Pr_t = 10$ (capped).
\item \textbf{Zero shear} ($S^2 \approx 0$): $Ri \to \infty$ $\Rightarrow$ $Pr_t = 10$ (capped). This prevents $\kappa_h$ from vanishing in quiescent regions, ensuring a background level of scalar diffusion.
\end{itemize}

\subsubsection{Alternative formulation with IDEMIX}

If the IDEMIX internal wave module is enabled (Section~\ref{sec:idemix}), the Prandtl number formula is modified to account for the contribution of internal wave breaking to scalar mixing:

\begin{lstlisting}[language=Fortran,caption={(ggl90\_calc.F:L563--574) TKE Prandtl number with IDEMIX}]
#ifdef ALLOW_GGL90_IDEMIX
       IF ( useIDEMIX ) THEN
        DO j=jMin,jMax
         DO i=iMin,iMax
          RiNumber = MAX(Nsquare(i,j,k),0. _d 0)
     &         /(verticalShear(i,j)+GGL90eps)
          IDEMIX_RiNumber = MAX( KappaM(i,j)*Nsquare(i,j,k), 0. _d 0)/
     &         ( GGL90eps + IDEMIX_gTKE(i,j,k) )
          prTemp          = 6.6 _d 0 * MIN( RiNumber, IDEMIX_RiNumber )
          TKEPrandtlNumber(i,j,k) = MIN(10. _d 0,prTemp)
          TKEPrandtlNumber(i,j,k) = MAX( oneRL,TKEPrandtlNumber(i,j,k) )
         ENDDO
        ENDDO
       ELSE
#endif /* ALLOW_GGL90_IDEMIX */
\end{lstlisting}

The IDEMIX-modified formula computes two Richardson numbers:
\begin{itemize}
\item \verb|RiNumber|: Standard gradient Richardson number (as before).
\item \verb|IDEMIX_RiNumber|: An effective Richardson number based on the TKE production rate from internal wave dissipation.
\end{itemize}

The Prandtl number is then $Pr_t = 6.6 \times \min(Ri, Ri_{\text{IDEMIX}})$, capped at 10. This formulation reduces $Pr_t$ (i.e., enhances scalar mixing) when internal waves contribute significantly to TKE production, reflecting the enhanced vertical mixing by wave breaking.

% ------------------------------------------------------------
\subsection{Eddy Diffusivity ($\kappa_h$)}
\label{sec:eddy_diff}
% ------------------------------------------------------------

Once $\kappa_m$ and $Pr_t$ are known, the eddy diffusivity for scalars (temperature and salinity) is simply:

\begin{equation}
\kappa_h = \frac{\kappa_m}{Pr_t}
\label{eq:kappa_h}
\end{equation}

This is computed in lines 593--598 of \verb|ggl90_calc.F|:

\begin{lstlisting}[language=Fortran,caption={(ggl90\_calc.F:L593--598) Eddy diffusivity and TKE diffusivity}]
       DO j=jMin,jMax
        DO i=iMin,iMax
C        diffusivity
         KappaH = KappaM(i,j)/TKEPrandtlNumber(i,j,k)
         KappaE(i,j,k) = GGL90alpha * KappaM(i,j)
     &        * maskC(i,j,k,bi,bj)*maskC(i,j,k-1,bi,bj)
\end{lstlisting}

\paragraph{Note on variable naming.} The local variable \verb|KappaH| is computed but not stored; instead, the scalar diffusivities are computed later from \verb|GGL90visctmp| (the viscosity floored at the scalar background diffusivity) in the routines \verb|GGL90_CALC_DIFF| (Section~\ref{sec:model_interface}).

The second line computes $\kappa_e = \alpha_e \kappa_m$, the diffusivity of TKE itself (Section~\ref{sec:tke_diffusion}), where $\alpha_e$ = \verb|GGL90alpha| (typically 1.0).

\subsubsection{Typical values}

Continuing the examples from Section~\ref{sec:eddy_visc}:
\begin{itemize}
\item \textbf{Surface mixed layer} ($Ri < 0.2$): $\kappa_m \sim 0.06$ m$^2$ s$^{-1}$, $Pr_t = 1$ $\Rightarrow$ $\kappa_h \sim 0.06$ m$^2$ s$^{-1}$.
\item \textbf{Thermocline} ($Ri \sim 1$): $\kappa_m \sim 3 \times 10^{-4}$ m$^2$ s$^{-1}$, $Pr_t = 5$ $\Rightarrow$ $\kappa_h \sim 6 \times 10^{-5}$ m$^2$ s$^{-1}$.
\item \textbf{Deep ocean} ($Ri \sim 5$): $\kappa_m \sim 5 \times 10^{-4}$ m$^2$ s$^{-1}$, $Pr_t = 10$ (capped) $\Rightarrow$ $\kappa_h \sim 5 \times 10^{-5}$ m$^2$ s$^{-1}$.
\end{itemize}

The ratio $\kappa_m / \kappa_h$ ranges from 1 (neutral/convective) to 10 (strongly stratified), consistent with atmospheric and oceanic observations.

% ------------------------------------------------------------
\subsection{Background Floors and Maximum Caps}
\label{sec:eddy_floors}
% ------------------------------------------------------------

To ensure numerical stability and physical realism, the GGL90 scheme imposes minimum (background) values on the eddy coefficients, preventing them from vanishing in quiescent regions where TKE is very small.

\subsubsection{Background viscosity}

The eddy viscosity $\kappa_m$ is floored at the background vertical viscosity \verb|viscArNr(k)|:

\begin{equation}
\kappa_m^{\text{total}} = \max(\kappa_m^{\text{GGL90}}, \, \kappa_m^{\text{background}})
\label{eq:visc_floor}
\end{equation}

where $\kappa_m^{\text{background}}$ = \verb|viscArNr(k)| is a depth-dependent background viscosity specified in the main model's \verb|data| file. Typical values are $\kappa_m^{\text{background}} = 10^{-4}$ to $10^{-3}$ m$^2$ s$^{-1}$ in the upper ocean, increasing to $10^{-3}$ to $10^{-2}$ m$^2$ s$^{-1}$ in the deep ocean to represent mixing by internal waves, double diffusion, and other sub-grid-scale processes not explicitly captured by GGL90.

This is implemented in lines 461--463 of \verb|ggl90_calc.F| (shown in Section~\ref{sec:eddy_visc}).

\subsubsection{Background diffusivity}

Similarly, the eddy diffusivity for scalars is floored at \verb|diffKrNrS(k)|:

\begin{equation}
\kappa_h^{\text{total}} = \max(\kappa_h^{\text{GGL90}}, \, \kappa_h^{\text{background}})
\label{eq:diff_floor}
\end{equation}

where $\kappa_h^{\text{background}}$ = \verb|diffKrNrS(k)| is typically $10^{-5}$ to $10^{-4}$ m$^2$ s$^{-1}$ in the upper ocean, representing molecular diffusion and small-scale mixing processes.

Importantly, the background floor for scalars is applied to \verb|GGL90visctmp| (line 456--458 of \verb|ggl90_calc.F|), which is used later to compute the scalar diffusivities in \verb|GGL90_CALC_DIFF|. This ensures that the scalar diffusivity is never less than the background, even in strongly stratified regions where $Pr_t$ may be large.

\subsubsection{Coordinate transformation}

In pressure coordinates, the background viscosities and diffusivities must be scaled by \verb|recip_coordFac*recip_coordFac| to convert from z-coordinate units (m$^2$ s$^{-1}$) to pressure-coordinate units (Pa$^2$ s$^{-1}$ or equivalent). This is automatically applied in lines 457 and 462 of \verb|ggl90_calc.F|.

\subsubsection{Maximum cap (optional)}

MITgcm does not impose an explicit maximum cap on GGL90 eddy coefficients by default. However, users can add such a cap via the parameter \verb|viscArMax| in the main model's \verb|data| file, which applies globally to all viscosity schemes (not just GGL90). This is rarely necessary in practice, as the mixing length constraints (Section~\ref{sec:mxl_method2}) and TKE dissipation (Section~\ref{sec:tke_dissipation}) naturally limit the growth of $\kappa_m$ and $\kappa_h$.

\subsubsection{Masking at dry interfaces}

All eddy coefficients are multiplied by \verb|maskC(i,j,k)*maskC(i,j,k-1)| to ensure they are zero at dry interfaces (e.g., adjacent to land, ice, or outside the domain). This prevents spurious fluxes across solid boundaries.

% ============================================================
\section{Boundary Conditions}
\label{sec:boundary_conditions}
% ============================================================

The TKE equation requires boundary conditions at the ocean surface and bottom (and optionally at ice-shelf bases). GGL90 uses a Dirichlet (prescribed value) condition at the surface and either Dirichlet or Neumann (zero flux) at the bottom, controlled by the runtime parameter \verb|GGL90_dirichlet|.

% ------------------------------------------------------------
\subsection{Surface Boundary Condition}
\label{sec:bc_surface}
% ------------------------------------------------------------

\subsubsection{Physical basis}

At the ocean surface, TKE is generated by wind stress and surface wave breaking. The standard boundary condition relates the near-surface TKE to the friction velocity $u_*$:

\begin{equation}
e_{\text{surface}} = B_0 \, u_*^2
\label{eq:surface_bc}
\end{equation}

where:
\begin{itemize}
\item $u_*$: Friction velocity (m s$^{-1}$), defined by the surface wind stress $\tau$: $u_*^2 = |\tau| / \rho_0$.
\item $B_0$ (\verb|GGL90m2|): Dimensionless constant relating surface TKE to $u_*^2$. The default value is $B_0 = 3.75$, based on atmospheric boundary layer measurements and the Blanke \& Delecluse (1993) calibration for ocean applications.
\end{itemize}

\paragraph{Physical interpretation.} The factor $B_0 = 3.75$ indicates that the TKE at the surface is several times larger than the mean kinetic energy of the wind-driven current ($\sim u_*^2$). This excess energy comes from surface wave motion and wave breaking, which inject TKE directly into the upper ocean. The calibrated value of 3.75 balances the need for adequate surface layer mixing against observations of mixed layer depth and SST evolution.

\paragraph{Minimum surface TKE.} To prevent the surface TKE from vanishing under calm conditions (which would cause numerical issues), GGL90 imposes a minimum value:

\begin{equation}
e_{\text{surface}} = \max(e_{\text{surf,min}}, \, B_0 u_*^2)
\label{eq:surface_bc_min}
\end{equation}

where $e_{\text{surf,min}}$ = \verb|GGL90TKEsurfMin| is typically $10^{-4}$ to $10^{-6}$ m$^2$ s$^{-2}$ (corresponding to a turbulent velocity scale of 1--10 mm s$^{-1}$).

\subsubsection{Friction velocity calculation}

The friction velocity $u_*$ is computed from the surface wind stress components $\tau_x$, $\tau_y$:

\begin{equation}
u_* = \sqrt{\frac{|\tau|}{\rho_0}} = \sqrt{\frac{\sqrt{\tau_x^2 + \tau_y^2}}{\rho_0}}
\label{eq:ustar_definition}
\end{equation}

In MITgcm, the surface stresses are provided by the fields \verb|surfaceForcingU| and \verb|surfaceForcingV|, which have units of m$^2$ s$^{-2}$ (kinematic stress = stress / density). Therefore, the code computes:

\begin{equation}
u_*^2 = \sqrt{\tau_x^2 + \tau_y^2}
\label{eq:ustar_squared}
\end{equation}

(Note: MITgcm's \verb|surfaceForcingU/V| are already divided by $\rho_0$, so no additional normalization is needed.)

The implementation (lines 764--795 of \verb|ggl90_calc.F|) offers two options, controlled by \verb|calcMeanVertShear|:

\paragraph{Option 1: Average of four components (\texttt{calcMeanVertShear = .TRUE.})}

\begin{lstlisting}[language=Fortran,caption={(ggl90\_calc.F:L764--783) Surface friction velocity: Option 1}]
      IF ( calcMeanVertShear ) THEN
C     by averaging (@ grid-cell center) the 4 components @ U,V pos.
       DO j=jMin,jMax
        DO i=iMin,iMax
         tempU  = surfaceForcingU( i ,j,bi,bj)
         tempUp = surfaceForcingU(i+1,j,bi,bj)
         tempV  = surfaceForcingV(i, j ,bi,bj)
         tempVp = surfaceForcingV(i,j+1,bi,bj)
         uStarSquare(i,j) =
     &        ( tempU*tempU + tempUp*tempUp
     &        + tempV*tempV + tempVp*tempVp
     &        )*halfRL
        ENDDO
       ENDDO
\end{lstlisting}

This averages the squared stresses from the four surrounding U- and V-points, which is appropriate for the Arakawa C-grid staggering.

\paragraph{Option 2: Shear of averaged stress (\texttt{calcMeanVertShear = .FALSE.})}

\begin{lstlisting}[language=Fortran,caption={(ggl90\_calc.F:L784--795) Surface friction velocity: Option 2}]
      ELSE
       DO j=jMin,jMax
        DO i=iMin,iMax
C     estimate friction velocity uStar from surface forcing
         uStarSquare(i,j) =
     &     ( .5 _d 0*( surfaceForcingU(i,  j,  bi,bj)
     &               + surfaceForcingU(i+1,j,  bi,bj) ) )**2
     &   + ( .5 _d 0*( surfaceForcingV(i,  j,  bi,bj)
     &               + surfaceForcingV(i,  j+1,bi,bj) ) )**2
        ENDDO
       ENDDO
      ENDIF
\end{lstlisting}

This averages the stresses first, then computes the squared magnitude. The difference is typically small.

After computing \verb|uStarSquare|, the code takes its square root and applies coordinate scaling (lines 864--873):

\begin{lstlisting}[language=Fortran,caption={(ggl90\_calc.F:L864--873) Square root and coordinate transformation}]
      DO j=jMin,jMax
       DO i=iMin,iMax
#ifdef ALLOW_AUTODIFF
         IF ( uStarSquare(i,j) .GT. zeroRL )
     &        uStarSquare(i,j) = SQRT(uStarSquare(i,j))*recip_coordFac
#else
         uStarSquare(i,j) = SQRT(uStarSquare(i,j))*recip_coordFac
#endif
       ENDDO
      ENDDO
\end{lstlisting}

Despite the variable name, \verb|uStarSquare| after this step actually contains $u_*$ (not $u_*^2$), scaled by \verb|recip_coordFac| for pressure coordinates.

\subsubsection{Application of surface boundary condition}

The surface TKE is imposed as a Dirichlet boundary condition in the implicit TKE solve. The implementation differs slightly between z-coordinates and pressure coordinates.

\paragraph{Z-coordinates (surface at $k=1$):}

\begin{lstlisting}[language=Fortran,caption={(ggl90\_calc.F:L899--912) Surface BC in z-coordinates}]
       DO j=jMin,jMax
        DO i=iMin,iMax
#ifdef ALLOW_SHELFICE
         IF ( useShelfIce ) THEN
          kSrf = MAX(1,kTopC(i,j,bi,bj))
          kTop = MIN(kSrf+1,Nr)
         ENDIF
#endif
         GGL90TKE(i,j,kSrf,bi,bj) = maskC(i,j,kSrf,bi,bj)
     &        *MAX(GGL90TKEsurfMin,GGL90m2*uStarSquare(i,j))
         GGL90TKE(i,j,kTop,bi,bj) = GGL90TKE(i,j,kTop,bi,bj)
     &        - a3d(i,j,kTop)*GGL90TKE(i,j,kSrf,bi,bj)
         a3d(i,j,kTop) = 0. _d 0
        ENDDO
\end{lstlisting}

\paragraph{Key details:}
\begin{itemize}
\item Line 907--908: Set the surface cell TKE to $\max(e_{\text{surf,min}}, B_0 u_*^2)$.
\item Lines 909--911: Subtract the contribution of the surface TKE from the RHS of the implicit equation for the second cell (since the surface TKE is now known and moves to the RHS).
\item Line 911: Zero out the lower diagonal coefficient for the second cell, effectively decoupling it from the surface cell in the tridiagonal solve.
\end{itemize}

\paragraph{Pressure coordinates (surface at $k=N_r$):}

\begin{lstlisting}[language=Fortran,caption={(ggl90\_calc.F:L883--897) Surface BC in pressure coordinates}]
      IF ( usingPCoords ) THEN
       DO j=jMin,jMax
        DO i=iMin,iMax
         GGL90TKE(i,j,kSrf,bi,bj) = GGL90TKE(i,j,kSrf,bi,bj)
     &        - c3d(i,j,kSrf) * maskC(i,j,kSrf,bi,bj)
     &        *MAX(GGL90TKEsurfMin,GGL90m2*uStarSquare(i,j))
         c3d(i,j,kSrf) = 0. _d 0
        ENDDO
       ENDDO
\end{lstlisting}

The logic is similar, but the upper diagonal \verb|c3d| (instead of lower diagonal \verb|a3d|) is zeroed out, reflecting the reversed vertical indexing in pressure coordinates.

\subsubsection{Surface minimum TKE}

The parameter \verb|GGL90TKEsurfMin| (default $10^{-4}$ m$^2$ s$^{-2}$) ensures that the surface TKE never vanishes, even under zero wind forcing. This is essential for:
\begin{itemize}
\item \textbf{Numerical stability}: Prevents division by zero in the mixing length and eddy coefficient calculations.
\item \textbf{Physical realism}: Even under calm conditions, there is residual turbulence from diurnal heating cycles, residual wave motion, and surface roughness.
\end{itemize}

Typical values: $e_{\text{surf,min}} = 10^{-6}$ m$^2$ s$^{-2}$ (very quiescent) to $10^{-4}$ m$^2$ s$^{-2}$ (moderate background turbulence).

% ------------------------------------------------------------
\subsection{Bottom Boundary Condition}
\label{sec:bc_bottom}
% ------------------------------------------------------------

At the ocean bottom (or topography in the interior), GGL90 supports two types of boundary conditions, controlled by the runtime parameter \verb|GGL90_dirichlet|:

\begin{enumerate}
\item \textbf{Dirichlet (prescribed TKE)}: $e_{\text{bottom}} = B_0 u_{*,\text{bottom}}^2$, where $u_{*,\text{bottom}}$ is the bottom friction velocity. This is the default (\verb|GGL90_dirichlet = .TRUE.|).
\item \textbf{Neumann (zero flux)}: $\partial e / \partial z = 0$ at the bottom. This assumes no TKE injection from bottom drag (\verb|GGL90_dirichlet = .FALSE.|).
\end{enumerate}

\subsubsection{Dirichlet bottom boundary condition}

When \verb|GGL90_dirichlet = .TRUE.| (default), the bottom TKE is set to:

\begin{equation}
e_{\text{bottom}} = B_0 \, u_{*,\text{bottom}}^2
\label{eq:bottom_bc_dirichlet}
\end{equation}

where $u_{*,\text{bottom}}$ is the friction velocity associated with bottom drag. In MITgcm, this is computed from the bottom stress (diagnosed from the velocity field and a quadratic drag law) in a manner analogous to the surface friction velocity.

However, the standard GGL90 implementation in MITgcm \emph{does not explicitly impose a non-zero bottom TKE} based on bottom drag. Instead, the bottom boundary condition is effectively $e_{\text{bottom}} = 0$ (or \verb|GGL90TKEbottom|, if specified). This is because:
\begin{itemize}
\item Computing bottom friction velocity requires knowledge of the near-bottom velocity, which introduces a two-way coupling between TKE and momentum that complicates the solver.
\item In most ocean applications, wind-driven turbulence dominates in the surface mixed layer, and bottom-generated turbulence is much weaker (except in shallow coastal regions or over rough topography).
\end{itemize}

For applications where bottom turbulence is important, users can:
\begin{itemize}
\item Set a non-zero background diffusivity (\verb|diffKrNrS(Nr)|) in the bottom layer.
\item Use a separate bottom boundary layer parameterization (e.g., a logarithmic law-of-the-wall).
\item Modify the code to explicitly compute $u_{*,\text{bottom}}$ and impose Eq.~\eqref{eq:bottom_bc_dirichlet}.
\end{itemize}

\subsubsection{Neumann bottom boundary condition}

When \verb|GGL90_dirichlet = .FALSE.|, the code imposes a zero-flux condition:

\begin{equation}
\frac{\partial e}{\partial z}\bigg|_{\text{bottom}} = 0
\label{eq:bottom_bc_neumann}
\end{equation}

This is implemented by zeroing out the appropriate off-diagonal coefficient in the tridiagonal matrix:

\begin{lstlisting}[language=Fortran,caption={(ggl90\_calc.F:L722--739) Neumann bottom BC}]
      IF (.NOT.GGL90_dirichlet) THEN
C      Neumann bottom boundary condition for TKE: no flux from bottom
       IF ( usingPCoords ) THEN
        DO j=jMin,jMax
         DO i=iMin,iMax
          kBot = MIN(kSurfC(i,j,bi,bj)+1,Nr)
          a3d(i,j,kBot) = 0. _d 0
         ENDDO
        ENDDO
       ELSE
        DO j=jMin,jMax
         DO i=iMin,iMax
          kBot = MAX(kLowC(i,j,bi,bj),1)
          c3d(i,j,kBot) = 0. _d 0
         ENDDO
        ENDDO
       ENDIF
      ENDIF
\end{lstlisting}

\paragraph{Physical interpretation.} The Neumann condition assumes that TKE is not generated at the bottom, and any TKE present in the bottom layer is due to downward diffusion from above (or local shear production). This is appropriate for:
\begin{itemize}
\item Deep ocean basins with weak bottom currents.
\item Coarse-resolution models where bottom boundary layer processes are not resolved.
\end{itemize}

For shallow coastal regions or continental shelves with strong tidal currents, the Neumann condition is \emph{not} appropriate, and a Dirichlet condition (with bottom friction velocity) should be used instead.

\subsubsection{Choice of boundary condition}

\begin{table}[ht]
\centering
\caption{Comparison of GGL90 bottom boundary conditions}
\label{tab:bottom_bc}
\begin{tabular}{lp{5cm}p{5cm}}
\toprule
\textbf{Parameter} & \textbf{Dirichlet} & \textbf{Neumann} \\
\midrule
\verb|GGL90_dirichlet| & \verb|.TRUE.| (default) & \verb|.FALSE.| \\
\addlinespace
Bottom condition & $e_{\text{bottom}} = 0$ (or user-specified) & $\partial e/\partial z = 0$ \\
\addlinespace
Best for & Deep ocean, weak bottom currents, typical global applications (Appendix~\ref{app:eccov4_config}) & Shallow coastal regions (if bottom drag is negligible) \\
\addlinespace
Physical interpretation & No TKE generation at bottom & No TKE flux through bottom \\
\bottomrule
\end{tabular}
\end{table}

\paragraph{Recommendation.} Use the default \verb|GGL90_dirichlet = .TRUE.| unless you have a specific reason to change it. For applications requiring realistic bottom boundary layer mixing, consider implementing a bottom friction velocity calculation (not included in the standard package).

% ------------------------------------------------------------
\subsection{Under-Ice-Shelf Treatment (Optional)}
\label{sec:bc_iceshelf}
% ------------------------------------------------------------

When the \verb|ALLOW_SHELFICE| package is enabled and \verb|useShelfIce = .TRUE.|, GGL90 automatically adjusts the surface boundary condition to account for ice-shelf cavities, where the "surface" is the ice-shelf base rather than the open ocean surface.

\subsubsection{Ice-shelf drag}

In columns beneath ice shelves, the friction velocity $u_*$ is computed from the drag between the ocean and the ice-shelf base, rather than from wind stress. The code computes the sub-glacial stress using the same quadratic drag law as for the ocean bottom:

\begin{equation}
\tau_{\text{ice}} = C_D \, |\vec{u}_{\text{ice-ocean}}| \, \vec{u}_{\text{ice-ocean}}
\label{eq:iceshelf_stress}
\end{equation}

where $C_D$ is the ice-shelf drag coefficient (specified by \verb|SHELFICEDragLinear| or \verb|SHELFICEDragQuadratic|) and $\vec{u}_{\text{ice-ocean}}$ is the ocean velocity at the ice-shelf base.

The implementation (lines 796--860 of \verb|ggl90_calc.F|) is complex and involves iterating through vertical levels to accumulate the total stress:

\begin{lstlisting}[language=Fortran,caption={(ggl90\_calc.F:L849--858) Ice-shelf friction velocity}]
       DO j=jMin,jMax
        DO i=jMin,jMax
         IF ( kTopC(i,j,bi,bj) .GT. 0 ) THEN
          uStarSquare(i,j) =
     &         ( stressU(i,j)*stressU(i,j)+stressU(i+1,j)*stressU(i+1,j)
     &         + stressV(i,j)*stressV(i,j)+stressV(i,j+1)*stressV(i,j+1)
     &         )*halfRL
         ENDIF
        ENDDO
       ENDDO
\end{lstlisting}

\paragraph{Key details:}
\begin{itemize}
\item \verb|kTopC(i,j,bi,bj)|: Index of the top wet cell (i.e., the first cell below the ice shelf). If \verb|kTopC > 0|, the column is beneath an ice shelf.
\item \verb|stressU|, \verb|stressV|: Ice-shelf drag stresses, computed by \verb|SHELFICE_U_DRAG_COEFF| and \verb|SHELFICE_V_DRAG_COEFF|.
\item The computed \verb|uStarSquare| replaces the surface wind stress in the surface boundary condition.
\end{itemize}

\subsubsection{Surface mixing length adjustment}

When \verb|mxlSurfFlag = .TRUE.| (Section~\ref{sec:mxl_surface_constraints}), the surface mixing length constraint is applied at the ice-shelf base (level \verb|kTopC|) instead of at the ocean surface, ensuring that the first two wet cells beneath the ice shelf are coupled by mixing.

\subsubsection{When to use ice-shelf treatment}

Enable \verb|useShelfIce = .TRUE.| when:
\begin{itemize}
\item Simulating Antarctic or Greenland ice-shelf cavities.
\item Including Antarctic bottom water formation or ice-shelf basal melt processes.
\end{itemize}

The ice-shelf treatment is \emph{not} needed for sea ice (use the \verb|SEAICE| package instead, which has its own coupling to GGL90).

% ============================================================
\section{Model Interface Routines}
\label{sec:model_interface}
% ============================================================

The GGL90 package communicates with the main MITgcm model through a set of interface routines that transfer eddy coefficients from the GGL90 arrays to the main model's mixing arrays. These routines are called during the vertical mixing phase of the time-stepping sequence.

Three key interface routines handle this communication:
\begin{enumerate}
\item \verb|GGL90_CALC_VISC|: Transfers eddy viscosity to momentum equations (Section~\ref{sec:visc_transfer}).
\item \verb|GGL90_CALC_DIFF|: Transfers eddy diffusivity to tracer equations (Section~\ref{sec:diff_transfer}).
\item \verb|GGL90_EXCHANGES|: Updates halo regions for TKE when horizontal diffusion is active (Section~\ref{sec:exchanges}).
\end{enumerate}

These routines are designed to be minimal and efficient, adding only the GGL90 contribution to the existing background mixing coefficients already initialized by the main model.

% ------------------------------------------------------------
\subsection{\texttt{GGL90\_CALC\_VISC}: Viscosity Transfer}
\label{sec:visc_transfer}
% ------------------------------------------------------------

\subsubsection{Purpose}

The subroutine \verb|GGL90_CALC_VISC| transfers the GGL90 eddy viscosity to the main model's momentum solver. It is called once per vertical level during the momentum equation time-stepping.

\subsubsection{Call signature}

\begin{lstlisting}[language=Fortran,caption={(ggl90\_calc\_visc.F:L7--10) Subroutine interface}]
      SUBROUTINE GGL90_CALC_VISC(
     I        bi, bj, iMin, iMax, jMin, jMax, k,
     U        KappaRU, KappaRV,
     I        myThid )
\end{lstlisting}

\paragraph{Arguments:}
\begin{itemize}
\item \verb|bi, bj|: Current tile indices (for domain decomposition).
\item \verb|iMin, iMax, jMin, jMax|: Horizontal index range for the computation.
\item \verb|k|: Vertical level index (1 to Nr).
\item \verb|KappaRU, KappaRV|: Input/output arrays for vertical viscosity at U-points and V-points (m$^2$ s$^{-1}$).
\item \verb|myThid|: Thread ID number (for parallel execution).
\end{itemize}

\subsubsection{Implementation}

The routine is remarkably simple, consisting of just two loops that add the GGL90 viscosity increment to the existing background viscosity:

\begin{lstlisting}[language=Fortran,caption={(ggl90\_calc\_visc.F:L47--59) Add GGL90 viscosity to momentum arrays}]
      DO j=jMin,jMax
       DO i=iMin,iMax
        KappaRU(i,j,k) = KappaRU(i,j,k)
     &                 + ( GGL90viscArU(i,j,k,bi,bj) - viscArNr(k) )
       ENDDO
      ENDDO

      DO j=jMin,jMax
       DO i=iMin,iMax
        KappaRV(i,j,k) = KappaRV(i,j,k)
     &                 + ( GGL90viscArV(i,j,k,bi,bj) - viscArNr(k) )
       ENDDO
      ENDDO
\end{lstlisting}

\paragraph{Key details:}
\begin{itemize}
\item \verb|KappaRU|, \verb|KappaRV| on entry contain the background viscosity \verb|viscArNr(k)| (set by the main model).
\item The increment \verb|GGL90viscArU - viscArNr(k)| represents the \emph{additional} viscosity from GGL90, beyond the background.
\item This ensures that the total viscosity is $\kappa_{\text{total}} = \kappa_{\text{background}} + (\kappa_{\text{GGL90}} - \kappa_{\text{background}}) = \kappa_{\text{GGL90}}$ when GGL90 is active, but reverts to $\kappa_{\text{background}}$ in cells where GGL90 produces less mixing than the background.
\item The subtraction of \verb|viscArNr(k)| is necessary because \verb|GGL90viscArU| and \verb|GGL90viscArV| already include the background floor (see Section~\ref{sec:eddy_floors}).
\end{itemize}

\subsubsection{Location in time-stepping sequence}

\verb|GGL90_CALC_VISC| is called from the main momentum routine \verb|MOM_VECINV| (or \verb|MOM_FLUXFORM|, depending on the momentum advection scheme) at each vertical level, immediately before the vertical viscosity term is applied. This allows GGL90 to override the background viscosity on a level-by-level basis.

% ------------------------------------------------------------
\subsection{\texttt{GGL90\_CALC\_DIFF}: Diffusivity Transfer}
\label{sec:diff_transfer}
% ------------------------------------------------------------

\subsubsection{Purpose}

The subroutine \verb|GGL90_CALC_DIFF| transfers the GGL90 eddy diffusivity to the main model's tracer solver. Unlike \verb|GGL90_CALC_VISC|, which is called separately for U and V velocities, this routine handles diffusivity for all scalar tracers (temperature, salinity, and any passive tracers) with a single call.

\subsubsection{Call signature}

\begin{lstlisting}[language=Fortran,caption={(ggl90\_calc\_diff.F:L7--10) Subroutine interface}]
      SUBROUTINE GGL90_CALC_DIFF(
     I        bi, bj, iMin, iMax, jMin, jMax, kArg, kSize,
     U        KappaRx,
     I        myThid )
\end{lstlisting}

\paragraph{Arguments:}
\begin{itemize}
\item \verb|bi, bj|: Current tile indices.
\item \verb|iMin, iMax, jMin, jMax|: Horizontal index range.
\item \verb|kArg|: Vertical level index. If \verb|kArg = 0|, process all levels; if \verb|kArg > 0|, process only level \verb|kArg|.
\item \verb|kSize|: Third dimension of the \verb|KappaRx| array (usually \verb|Nr|, but may differ in some configurations).
\item \verb|KappaRx|: Input/output array for vertical diffusivity (m$^2$ s$^{-1}$).
\item \verb|myThid|: Thread ID number.
\end{itemize}

\subsubsection{Implementation}

The routine supports two calling modes: process all levels, or process a single level.

\paragraph{Mode 1: Process all levels (\texttt{kArg = 0})}

\begin{lstlisting}[language=Fortran,caption={(ggl90\_calc\_diff.F:L48--58) Add GGL90 diffusivity: all levels}]
      IF ( kArg .EQ. 0 ) THEN
C-    do all levels :
        DO k=1,MIN(Nr,kSize)
         DO j=jMin,jMax
          DO i=iMin,iMax
            KappaRx(i,j,k) = KappaRx(i,j,k)
     &                  +( GGL90diffKr(i,j,k,bi,bj)
     &                     - diffKrNrS(k) )
          ENDDO
         ENDDO
        ENDDO
\end{lstlisting}

\paragraph{Mode 2: Process single level (\texttt{kArg > 0})}

\begin{lstlisting}[language=Fortran,caption={(ggl90\_calc\_diff.F:L59--69) Add GGL90 diffusivity: single level}]
      ELSE
C-    do level k=kArg only :
         k = MIN(kArg,kSize)
         DO j=jMin,jMax
          DO i=iMin,iMax
            KappaRx(i,j,k) = KappaRx(i,j,k)
     &                  +( GGL90diffKr(i,j,kArg,bi,bj)
     &                     - diffKrNrS(kArg) )
          ENDDO
         ENDDO
      ENDIF
\end{lstlisting}

\paragraph{Key details:}
\begin{itemize}
\item Similar to \verb|GGL90_CALC_VISC|, the routine adds the GGL90 diffusivity increment \verb|GGL90diffKr - diffKrNrS(k)| to the existing background diffusivity.
\item \verb|diffKrNrS(k)| is the background scalar diffusivity (typically $10^{-5}$ to $10^{-4}$ m$^2$ s$^{-1}$ in the upper ocean).
\item The two-mode design allows the main model to call \verb|GGL90_CALC_DIFF| either once per time step (all levels at once) or once per level (level-by-level), depending on the tracer advection scheme.
\end{itemize}

% ------------------------------------------------------------
\subsection{\texttt{GGL90\_EXCHANGES}: Halo Updates}
\label{sec:exchanges}
% ------------------------------------------------------------

\subsubsection{Purpose}

The subroutine \verb|GGL90_EXCHANGES| updates the halo (overlap) regions of the TKE array \verb|GGL90TKE| to ensure consistency across tile boundaries in domain-decomposed simulations. This is necessary when horizontal diffusion of TKE is active (\verb|GGL90diffTKEh > 0|) or when the optional IDEMIX package is used with horizontal internal wave energy diffusion.

\subsubsection{Implementation}

\begin{lstlisting}[language=Fortran,caption={(ggl90\_exchanges.F:L31--40) Conditional halo exchanges}]
#ifdef ALLOW_GGL90_IDEMIX
      IF ( useIDEMIX .AND. IDEMIX_tau_h .GT. zeroRL ) THEN
       _EXCH_XYZ_RL( GGL90TKE, myThid )
       _EXCH_XYZ_RL( IDEMIX_E, myThid )
      ELSEIF ( GGL90diffTKEh .GT. zeroRL ) THEN
#else /* ALLOW_GGL90_IDEMIX */
      IF ( GGL90diffTKEh .GT. zeroRL ) THEN
#endif /* ALLOW_GGL90_IDEMIX */
       _EXCH_XYZ_RL( GGL90TKE, myThid )
      ENDIF
\end{lstlisting}

\paragraph{Key details:}
\begin{itemize}
\item \verb|_EXCH_XYZ_RL| is a MITgcm macro that expands to the appropriate halo exchange routine for 3D real*8 arrays.
\item Halo exchanges are \emph{only} performed when horizontal TKE diffusion or IDEMIX horizontal diffusion is enabled.
\item If neither feature is active, the routine is a no-op (does nothing).
\end{itemize}

% ============================================================
\section{Initialization and Restart}
\label{sec:initialization}
% ============================================================

GGL90 initialization proceeds in three stages: reading runtime parameters (Section~\ref{sec:readparms}), initializing time-invariant fields (Section~\ref{sec:init_fixed}), and initializing time-varying fields including restart capability (Section~\ref{sec:init_varia}).

% ------------------------------------------------------------
\subsection{\texttt{GGL90\_READPARMS}: Parameter Input}
\label{sec:readparms}
% ------------------------------------------------------------

The routine \verb|GGL90_READPARMS| reads runtime parameters from the file \verb|data.ggl90| in the run directory. This file contains three Fortran namelists:

\begin{itemize}
\item \verb|GGL90_PARM01|: Core GGL90 parameters (required)
\item \verb|GGL90_PARM02|: IDEMIX parameters (optional, only if \verb|ALLOW_GGL90_IDEMIX| is defined)
\item \verb|GGL90_PARM03|: Langmuir circulation parameters (optional, only if \verb|ALLOW_GGL90_LANGMUIR| is defined)
\end{itemize}

\paragraph{Default values} (from \verb|ggl90_readparms.F:L102--127}):
\begin{itemize}
\item \verb|GGL90ck = 0.1|: Eddy viscosity constant $c_\mu$
\item \verb|GGL90ceps = 0.7|: Dissipation constant $c_\varepsilon$
\item \verb|GGL90alpha = 1.0|: TKE diffusivity constant $\alpha_e$
\item \verb|GGL90m2 = 3.75|: Surface TKE constant $B_0$ (originally from Blanke \& Delecluse 1993)
\item \verb|GGL90TKEmin = 1e-11|: Minimum TKE (m$^2$ s$^{-2}$)
\item \verb|GGL90TKEsurfMin = 1e-4|: Minimum surface TKE (m$^2$ s$^{-2}$)
\item \verb|GGL90mixingLengthMin = 1e-8|: Minimum mixing length (m)
\item \verb|mxlMaxFlag = 0|: Mixing length method (0, 1, or 2)
\item \verb|mxlSurfFlag = .FALSE.|: Force surface mixing length to grid spacing
\item \verb|GGL90_dirichlet = .TRUE.|: Use Dirichlet bottom BC for TKE
\item \verb|calcMeanVertShear = .FALSE.|: Method for computing vertical shear
\end{itemize}

% ------------------------------------------------------------
\subsection{\texttt{GGL90\_INIT\_FIXED}: Time-Invariant Initialization}
\label{sec:init_fixed}
% ------------------------------------------------------------

The routine \verb|GGL90_INIT_FIXED| initializes fields that do not change during the simulation. This is a minimal routine that primarily registers diagnostics:

\begin{lstlisting}[language=Fortran,caption={(ggl90\_init\_fixed.F:L39--43) Diagnostics registration}]
#ifdef ALLOW_DIAGNOSTICS
      IF ( useDiagnostics ) THEN
        CALL GGL90_DIAGNOSTICS_INIT( myThid )
      ENDIF
#endif
\end{lstlisting}

If horizontal smoothing is enabled (\verb|ALLOW_GGL90_SMOOTH|), the routine also initializes a mask array for handling cubed-sphere corner points (lines 45--63).

% ------------------------------------------------------------
\subsection{\texttt{GGL90\_INIT\_VARIA}: Time-Varying Field Initialization}
\label{sec:init_varia}
% ------------------------------------------------------------

The routine \verb|GGL90_INIT_VARIA| initializes all GGL90 prognostic and diagnostic arrays at the start of a simulation. The key steps are:

\paragraph{1. Zero out eddy coefficient arrays:}
\begin{lstlisting}[language=Fortran,caption={(ggl90\_init\_varia.F:L42--56) Initialize arrays to zero}]
        DO k=1,Nr
         DO j=1-OLy,sNy+OLy
          DO i=1-OLx,sNx+OLx
           GGL90viscArU(i,j,k,bi,bj) = 0. _d 0
           GGL90viscArV(i,j,k,bi,bj) = 0. _d 0
           GGL90diffKr(i,j,k,bi,bj)  = 0. _d 0
           IF ( useIDEMIX) THEN
            GGL90TKE(i,j,k,bi,bj)=GGL90eps*maskC(i,j,k,bi,bj)
           ELSE
            GGL90TKE(i,j,k,bi,bj)=GGL90TKEmin*maskC(i,j,k,bi,bj)
           ENDIF
          ENDDO
         ENDDO
        ENDDO
\end{lstlisting}

Note that TKE is initialized to \verb|GGL90TKEmin| (or \verb|GGL90eps| if IDEMIX is active), not zero, to avoid division by zero in the first time step.

\paragraph{2. Read initial TKE from file (optional):}

If the parameter \verb|GGL90TKEFile| is specified and this is not a restart run:

\begin{lstlisting}[language=Fortran,caption={(ggl90\_init\_varia.F:L135--150) Read initial TKE from file}]
       IF ( GGL90TKEFile .NE. ' ' ) THEN
        CALL READ_FLD_XYZ_RL( GGL90TKEFile, ' ', GGL90TKE, 0, myThid )
        _EXCH_XYZ_RL(GGL90TKE,myThid)
        DO bj=myByLo(myThid),myByHi(myThid)
         DO bi=myBxLo(myThid),myBxHi(myThid)
          DO k=1,Nr
           DO j=1-OLy,sNy+OLy
            DO i=1-OLx,sNx+OLx
             GGL90TKE(i,j,k,bi,bj) = MAX(GGL90TKE(i,j,k,bi,bj),
     &            GGL90TKEmin)*maskC(i,j,k,bi,bj)
            ENDDO
           ENDDO
          ENDDO
         ENDDO
        ENDDO
       ENDIF
\end{lstlisting}

The file should contain a 3D field of TKE values (m$^2$ s$^{-2}$) in MITgcm's standard binary format.

\paragraph{3. Read restart pickup file:}

If this is a restart run (\verb|nIter0 != 0| or \verb|pickupSuff != ' '|), load the prognostic TKE field from the pickup file:

\begin{lstlisting}[language=Fortran,caption={(ggl90\_init\_varia.F:L131--133) Load from restart}]
      IF ( nIter0.NE.0 .OR. pickupSuff.NE.' ' ) THEN
       CALL GGL90_READ_PICKUP( nIter0, myThid )
      ELSE
\end{lstlisting}

% ------------------------------------------------------------
\subsection{Restart Files: \texttt{GGL90\_READ/WRITE\_PICKUP}}
\label{sec:restart}
% ------------------------------------------------------------

GGL90 uses MITgcm's standard pickup file mechanism to save and restore the prognostic TKE field during restart.

\paragraph{Pickup file contents:}
\begin{itemize}
\item \verb|GGL90TKE|: 3D TKE field (m$^2$ s$^{-2}$)
\item \verb|IDEMIX_E|: 3D internal wave energy field (m$^2$ s$^{-2}$), only if \verb|ALLOW_GGL90_IDEMIX| is active
\end{itemize}

\paragraph{File naming convention:}
\begin{itemize}
\item Standard restart: \verb|pickup_ggl90.ITERATION.data| and \verb|.meta|
\item Checkpoint: \verb|pickup_ggl90.ckptA.data| and \verb|.meta| (or \verb|.ckptB|)
\end{itemize}

\paragraph{Usage:} No user action is required. MITgcm automatically calls \verb|GGL90_WRITE_PICKUP| at the end of each run and \verb|GGL90_READ_PICKUP| at the start of a restart run.

% ============================================================
\section{Compile-Time Options}
\label{sec:cpp_options}
% ============================================================

GGL90 behavior is controlled by a combination of runtime parameters (in \verb|data.ggl90|) and compile-time CPP flags (in \verb|GGL90_OPTIONS.h|). This section documents the available compile-time options.

% ------------------------------------------------------------
\subsection{Core Options}
\label{sec:core_cpp}
% ------------------------------------------------------------

\paragraph{\texttt{ALLOW\_GGL90}}

Master flag to enable the GGL90 package. Must be defined in \verb|packages.conf| to compile GGL90 routines. If not defined, all GGL90 code is excluded from compilation.

\paragraph{\texttt{GGL90\_REGULARIZE\_MIXINGLENGTH}}

If defined, use a regularized formula for the minimum mixing length (Section~\ref{sec:mxl_surface_constraints}):
\begin{equation*}
\ell_{\text{final}} = \sqrt{\ell^2 + \ell_{\text{min}}^2}
\end{equation*}
instead of the standard $\max(\ell, \ell_{\text{min}})$. This provides smoother derivatives for adjoint calculations (see Appendix~\ref{app:adjoint_considerations}). Default: \emph{not defined}.

\paragraph{\texttt{GGL90\_MISSING\_HFAC\_BUG}}

Legacy flag related to partial cell thickness handling in IDEMIX. Retains a historical bug for backward compatibility with older configurations (see Appendix~\ref{app:adjoint_considerations}). Default: \emph{not defined}. New users should \emph{not} define this.

% ------------------------------------------------------------
\subsection{Optional Feature Flags}
\label{sec:optional_cpp}
% ------------------------------------------------------------

\paragraph{\texttt{ALLOW\_GGL90\_HORIZDIFF}}

Enable horizontal diffusion of TKE. When defined, the runtime parameter \verb|GGL90diffTKEh| controls the horizontal diffusivity (m$^2$ s$^{-1}$). Default: \emph{not defined}. This feature is rarely used in practice (see Section~\ref{sec:tke_time_integration}).

\paragraph{\texttt{ALLOW\_GGL90\_SMOOTH}}

Enable horizontal smoothing of eddy viscosity and diffusivity before transferring to the main model. Uses a 9-point stencil following the OPA/NEMO approach. Default: \emph{not defined}. Useful for coarse-resolution models to reduce grid-scale noise.

\paragraph{\texttt{ALLOW\_GGL90\_IDEMIX}}

Enable the IDEMIX internal wave energy parameterization (Section~\ref{sec:idemix}). Adds internal wave breaking as a TKE source. Default: \emph{not defined}. Requires additional parameters in \verb|GGL90_PARM02| namelist.

\paragraph{\texttt{ALLOW\_GGL90\_LANGMUIR}}

Enable Langmuir circulation parameterization. Enhances surface layer mixing under wind-wave forcing. Default: \emph{not defined}. Requires additional parameters in \verb|GGL90_PARM03| namelist.

% ============================================================
\section{Output and Diagnostics}
\label{sec:diagnostics}
% ============================================================

GGL90 provides a comprehensive set of diagnostic fields for monitoring mixing processes. All diagnostics are registered via MITgcm's standard diagnostics package.

% ------------------------------------------------------------
\subsection{Available Diagnostic Fields}
\label{sec:diag_fields}
% ------------------------------------------------------------

Table~\ref{tab:ggl90_diagnostics} lists the available GGL90 diagnostic fields.

\begin{table}[ht]
\centering
\caption{GGL90 diagnostic fields}
\label{tab:ggl90_diagnostics}
\small
\begin{tabular}{llp{5cm}l}
\toprule
\textbf{Name} & \textbf{Dims} & \textbf{Description} & \textbf{Units} \\
\midrule
\verb|GGL90TKE| & 3D & Turbulent kinetic energy & m$^2$ s$^{-2}$ \\
\verb|GGL90ArU| & 3D & Eddy viscosity at U-points & m$^2$ s$^{-1}$ \\
\verb|GGL90ArV| & 3D & Eddy viscosity at V-points & m$^2$ s$^{-1}$ \\
\verb|GGL90Kr | & 3D & Eddy diffusivity (scalars) & m$^2$ s$^{-1}$ \\
\verb|GGL90mixL| & 3D & Mixing length & m \\
\verb|GGL90Prandtl| & 3D & Turbulent Prandtl number & -- \\
\addlinespace
\verb|GGL90flx | & 2D & Surface TKE flux & m$^3$ s$^{-3}$ \\
\verb|GGL90tau | & 2D & Surface friction velocity & m s$^{-1}$ \\
\bottomrule
\end{tabular}
\end{table}

\paragraph{Usage example:} To output hourly snapshots of TKE and eddy diffusivity, add to \verb|data.diagnostics|:

\begin{verbatim}
 fields(1:2,1) = 'GGL90TKE','GGL90Kr ',
 filename(1) = 'ggl90_snap',
 frequency(1) = 3600.0,
\end{verbatim}

% ------------------------------------------------------------
\subsection{Diagnostics Registration}
\label{sec:diag_registration}
% ------------------------------------------------------------

Diagnostics are registered in \verb|GGL90_DIAGNOSTICS_INIT|, called from \verb|GGL90_INIT_FIXED|. The registration is automatic and requires no user action. To enable diagnostics output, set \verb|useDiagnostics = .TRUE.| in the main \verb|data| file and configure the desired fields in \verb|data.diagnostics|.

% ============================================================
\section{IDEMIX Extension (Optional)}
\label{sec:idemix}
% ============================================================

IDEMIX (Internal Wave Dissipation and MIXing) is an optional extension to GGL90 that parameterizes the generation, propagation, and dissipation of internal wave energy. When enabled (\verb|ALLOW_GGL90_IDEMIX| in \verb|GGL90_OPTIONS.h|), IDEMIX adds internal wave breaking as an additional source term in the TKE equation.

\subsection{Overview of Internal Wave Model}

IDEMIX represents the internal wave field as a single prognostic variable $E$ (internal wave energy per unit mass, m$^2$ s$^{-2}$), which evolves according to:

\begin{equation}
\frac{\partial E}{\partial t} = F_{\text{sources}} - \frac{E}{\tau_v} + \frac{\partial}{\partial z}\left(\kappa_E \frac{\partial E}{\partial z}\right) + \nabla_h \cdot (\kappa_{E,h} \nabla_h E)
\end{equation}

where:
\begin{itemize}
\item $F_{\text{sources}}$: Energy input from tidal forcing and near-inertial wind forcing
\item $\tau_v$: Vertical dissipation time scale (typically 2--10 days)
\item $\kappa_E$, $\kappa_{E,h}$: Vertical and horizontal diffusivities of internal wave energy
\end{itemize}

The dissipated internal wave energy is converted to TKE with an efficiency factor $\gamma_{\text{mix}}$ (typically 0.1--0.2).

\subsection{Coupling to GGL90}

IDEMIX couples to GGL90 through two mechanisms:

\paragraph{1. TKE source term:} The internal wave dissipation rate $\varepsilon_{\text{IW}} = E / \tau_v$ is added to the TKE equation (Section~\ref{sec:tke_evolution}):

\begin{equation}
\frac{\partial e}{\partial t} = \ldots + \gamma_{\text{mix}} \frac{E}{\tau_v}
\end{equation}

This enhances mixing in the ocean interior where internal wave breaking is significant (e.g., over rough topography, at the equator).

\paragraph{2. Modified Prandtl number:} The turbulent Prandtl number formula is modified to account for internal wave-driven mixing (Section~\ref{sec:tke_prandtl}), reducing $Pr_t$ and enhancing scalar diffusivity when $E$ is large.

\paragraph{When to use IDEMIX:} IDEMIX is beneficial for global ocean simulations where interior mixing from internal tides and near-inertial waves is important for water mass properties and overturning circulation. It requires additional forcing files (\verb|IDEMIX_tidal_file|, \verb|IDEMIX_wind_file|) and runtime parameters in the \verb|GGL90_PARM02| namelist. For regional or process studies focused on surface boundary layer mixing, IDEMIX can be omitted.

% ============================================================
\section{Package Validation: \texttt{GGL90\_CHECK}}
\label{sec:validation}
% ============================================================

The routine \verb|GGL90_CHECK| performs compile-time and runtime validation of GGL90 parameters. It is called during model initialization to catch common configuration errors before the simulation begins.

\paragraph{Checks performed:}
\begin{itemize}
\item All required runtime parameters have been set (not \verb|UNSET_RL|)
\item Physical constants are positive (e.g., \verb|GGL90ck > 0|, \verb|GGL90TKEmin > 0|)
\item Mixing length method \verb|mxlMaxFlag| is valid (0, 1, or 2)
\item Background mixing coefficients are physically reasonable
\item If IDEMIX is enabled, all IDEMIX parameters are set
\end{itemize}

If any check fails, the routine prints an error message and stops the model with \verb|STOP 'ABNORMAL END'|. This prevents silent failures due to misconfiguration.

\paragraph{User action:} If \verb|GGL90_CHECK| reports an error, review the parameters in \verb|data.ggl90| and correct the issue. Common problems include typos in parameter names (Fortran namelists are case-insensitive but require exact spelling) and missing parameters when optional features are enabled.

% ============================================================
\section{Frequently Encountered Issues}
\label{sec:issues}
% ============================================================

This section provides practical guidance for diagnosing and resolving common GGL90 problems.

\subsection{Parameter Tuning Guidance}

\paragraph{Example starting configuration:}

For a well-tested baseline configuration suitable for global ocean applications, see Appendix~\ref{app:eccov4_config}. This provides a complete \texttt{data.ggl90} file with physically motivated parameter choices.

\paragraph{Tuning for specific applications:}
\begin{itemize}
\item \textbf{Too much mixing:} Reduce \verb|GGL90ck| (e.g., 0.08) or increase \verb|GGL90ceps| (e.g., 0.9).
\item \textbf{Too little mixing:} Increase \verb|GGL90ck| (e.g., 0.12) or decrease \verb|GGL90ceps| (e.g., 0.5).
\item \textbf{Unrealistic surface mixed layer depth:} Adjust \verb|GGL90m2| (surface TKE constant). Lower values give shallower mixed layers.
\item \textbf{Numerical instability:} Ensure \verb|GGL90TKEmin > 1e-8| and \verb|GGL90mixingLengthMin > 1e-10|.
\end{itemize}

\subsection{Common Mistakes and Solutions}

\paragraph{1. Forgetting to enable GGL90 in data file:}

\textbf{Symptom:} Model runs but mixing is weak, diagnostics show zero TKE.

\textbf{Solution:} Set \verb|useGGL90 = .TRUE.| in the main \verb|data| file (not \verb|data.ggl90|).

\paragraph{2. Incompatible coordinate system:}

\textbf{Symptom:} Mixing length values are unrealistically small or large.

\textbf{Solution:} GGL90 handles z- and p-coordinates automatically via \verb|coordFac|. No user action needed, but verify that \verb|drF| and \verb|drC| are correctly set in \verb|data.grid|.

\paragraph{3. Missing pickup file on restart:}

\textbf{Symptom:} Model crashes with ``file not found: pickup_ggl90.XXX.data''.

\textbf{Solution:} Ensure the pickup file from the previous run is in the run directory. If starting from a checkpoint that predates GGL90 activation, set \verb|nIter0 = 0| to cold-start GGL90.

\paragraph{4. Conflicting mixing schemes:}

\textbf{Symptom:} Excessive or oscillatory mixing, conservation errors.

\textbf{Solution:} Disable other vertical mixing schemes (KPP, PP81, MY82) when using GGL90. Only one turbulence closure should be active at a time.

\subsection{Performance Considerations}

\paragraph{Computational cost:} GGL90 adds approximately 10--20\% to the total runtime of a typical ocean simulation, dominated by:
\begin{itemize}
\item TKE implicit solve (tridiagonal system, one solve per horizontal point per time step)
\item Mixing length computation (especially Method 2 with two vertical sweeps)
\item Eddy coefficient interpolation to velocity points
\end{itemize}

\paragraph{Optimization strategies:}
\begin{itemize}
\item Use Method 1 instead of Method 2 for mixing length if the extra accuracy is not needed (saves ~30\% of GGL90 cost).
\item Disable horizontal TKE diffusion (\verb|GGL90diffTKEh = 0|, the default) to avoid halo exchanges.
\item For adjoint calculations, see Appendix~\ref{app:adjoint_considerations} for optimization-specific settings and tuning recommendations.
\end{itemize}

% ============================================================
\section{Conclusions}
\label{sec:conclusions}
% ============================================================

This document provides a comprehensive technical reference for the GGL90 turbulence closure scheme as implemented in the MIT General Circulation Model (MITgcm). We have documented:

\begin{enumerate}
\item \textbf{Physical foundation:} The TKE-based turbulence closure theory of Gaspar et al.\ (1990), including the governing equations, mixing length formulation, and eddy coefficient parameterization (Sections~\ref{sec:background}--\ref{sec:eddy_coeffs}).

\item \textbf{Numerical implementation:} Complete line-by-line correspondence with MITgcm source code (\verb|pkg/ggl90/*.F|), covering initialization, time integration, boundary conditions, and model interfaces (Sections~\ref{sec:architecture}--\ref{sec:diagnostics}).

\item \textbf{Input/output contracts:} Comprehensive tables and data flow diagram documenting all exchanges between GGL90 and the main model (Appendix~\ref{app:io_mapping}).

\item \textbf{Practical guidance:} Parameter tuning recommendations, common issues, and performance considerations (Section~\ref{sec:issues}).
\end{enumerate}

\paragraph{Key design features of GGL90:}
\begin{itemize}
\item \textbf{Prognostic TKE:} Solves a transport equation for turbulent kinetic energy, allowing the turbulence field to respond dynamically to changing forcing and stratification.

\item \textbf{Mixing length from stability:} Computes the vertical scale of turbulent eddies from local TKE and buoyancy frequency, with geometric constraints preventing unphysical penetration across stable layers (Method 2, Blanke \& Delecluse 1993).

\item \textbf{Richardson number-dependent Prandtl number:} Reduces scalar diffusivity relative to momentum diffusivity in stratified regions, consistent with atmospheric and oceanic observations.

\item \textbf{Clean model interface:} GGL90 receives velocity, density gradient, and wind stress as inputs, and returns eddy coefficients as outputs, with no direct access to tracers or other model internals. This modular design facilitates testing, comparison with other schemes, and portability to other ocean models.
\end{itemize}

\paragraph{Recommendations for users:}
\begin{itemize}
\item Start with a well-validated configuration (e.g., Appendix~\ref{app:eccov4_config} for global applications) as a baseline.
\item Monitor TKE and eddy coefficient diagnostics (\verb|GGL90TKE|, \verb|GGL90Kr|) to verify that mixing behavior is physically reasonable.
\item Compare against observations (e.g., Argo mixed layer depth, mooring time series) when tuning parameters for regional applications.
\item Consult the original references (Gaspar et al.\ 1990; Blanke \& Delecluse 1993) for the underlying physics and theory.
\item For adjoint-based applications, see Appendix~\ref{app:adjoint_considerations} for gradient-specific considerations.
\end{itemize}

\paragraph{Future development:} Potential enhancements to GGL90 include:
\begin{itemize}
\item Bottom boundary layer parameterization with explicit bottom friction velocity
\item Langmuir turbulence with Stokes drift from a wave model (already available via \verb|ALLOW_GGL90_LANGMUIR|)
\item Anisotropic mixing length for submesoscale-resolving simulations
\item Integration with sea ice models for under-ice turbulence
\end{itemize}

The GGL90 package remains under active development within the MITgcm community, with ongoing refinements for global ocean state estimation and high-resolution process studies.

% ============================================================
\appendix
% ============================================================

% ============================================================
\section{ECCOv4 Configuration}
\label{app:eccov4_config}
% ============================================================

The ECCO (Estimating the Circulation and Climate of the Ocean) consortium uses non-default GGL90 settings optimized for global ocean state estimation with $\sim$1$^\circ$ resolution. This appendix consolidates all ECCOv4-specific configuration details, parameter choices, and physical rationale.

\subsection{Overview: ECCOv4 and GGL90}

ECCOv4 is a data-assimilating ocean state estimate that synthesizes satellite altimetry, Argo float profiles, hydrographic sections, and other observations to produce a dynamically consistent reconstruction of global ocean circulation from 1992 to present. The model component uses MITgcm with GGL90 vertical mixing to represent turbulent processes across a wide range of oceanic regimes, from polar convection to tropical thermoclines.

The ECCOv4 GGL90 configuration was determined through extensive sensitivity studies and adjoint-based parameter estimation. Key objectives included:
\begin{itemize}[leftmargin=*]
  \item Maintaining sharp thermoclines consistent with Argo observations
  \item Preventing unrealistic deep stratification in quiescent regions
  \item Providing adequate surface boundary layer mixing for accurate SST simulation
  \item Ensuring adjoint stability and convergence for gradient-based optimization
\end{itemize}

\subsection{ECCOv4 Parameter Choices}

Table~\ref{tab:eccov4_params} summarizes the key differences between ECCOv4 settings and the Gaspar et al.\ (1990) default values.

\begin{table}[h!]
\centering
\caption{ECCOv4 parameter choices compared to GGL90 defaults.}
\label{tab:eccov4_params}
\begin{tabular}{@{}llll@{}}
\toprule
\textbf{Parameter} & \textbf{Default} & \textbf{ECCOv4} & \textbf{Units} \\
\midrule
\texttt{GGL90alpha} & 1.0 & 30.0 & -- \\
\texttt{GGL90TKEmin} & $10^{-11}$ & $10^{-7}$ & m$^2$/s$^2$ \\
\texttt{GGL90TKEbottom} & $10^{-11}$ & $10^{-6}$ & m$^2$/s$^2$ \\
\texttt{mxlMaxFlag} & 0 & 2 & -- \\
\texttt{mxlSurfFlag} & \texttt{.FALSE.} & \texttt{.TRUE.} & logical \\
\texttt{GGL90ck} & 0.1 & 0.1 & -- \\
\texttt{GGL90ceps} & 0.7 & 0.7 & -- \\
\texttt{GGL90m2} & 3.75 & 3.75 & -- \\
\bottomrule
\end{tabular}
\end{table}

\subsection{Physical Rationale for ECCOv4 Settings}

\subsubsection{GGL90alpha = 30.0 (Turbulent Prandtl Number)}

The turbulent Prandtl number $\alpha$ relates eddy diffusivity to eddy viscosity:
\[
\kappa_h = \frac{\kappa_m}{\alpha}
\]

\textbf{Default ($\alpha = 1$):} In neutral boundary layers (atmospheric or oceanic mixed layers), momentum and heat are transported with similar efficiency, suggesting $\alpha \approx 1$. This is the Gaspar et al.\ (1990) choice.

\textbf{ECCOv4 ($\alpha = 30$):} In the stratified ocean interior, scalar diffusion is suppressed more strongly than momentum diffusion. A large $\alpha$ dramatically reduces $\kappa_h$ relative to $\kappa_m$, maintaining sharp thermoclines and haloclines consistent with Argo float observations. This prevents the excessive erosion of vertical gradients that occurs with $\alpha = 1$ in coarse-resolution models.

\textbf{Physical interpretation:} Stratified turbulence is anisotropic. Vertical scalar transport requires work against buoyancy forces, while horizontal momentum transport does not. The effective Prandtl number in the ocean interior can exceed 10--100 depending on stratification strength.

\subsubsection{GGL90TKEmin = $10^{-7}$ m$^2$/s$^2$ (Background TKE)}

\textbf{Default ($10^{-11}$):} Represents molecular-scale motions, essentially zero for practical purposes.

\textbf{ECCOv4 ($10^{-7}$):} Higher background TKE represents unresolved internal wave breaking and other sub-grid mixing processes. This is critical for:
\begin{itemize}[leftmargin=*]
  \item Preventing unrealistic stratification buildup in the deep ocean
  \item Maintaining diapycnal mixing rates ($\kappa \sim 10^{-5}$ m$^2$/s) consistent with microstructure observations
  \item Avoiding numerical issues in regions with near-zero shear and stratification
\end{itemize}

\textbf{Observational support:} Tracer release experiments (Ledwell et al., 1993) and fine-scale parameterizations (Polzin et al., 1995) suggest basin-averaged deep ocean diffusivities of $\sim 10^{-5}$ m$^2$/s, which GGL90 achieves with $\TKE_{\mathrm{min}} \sim 10^{-7}$ m$^2$/s$^2$.

\subsubsection{GGL90TKEbottom = $10^{-6}$ m$^2$/s$^2$ (Bottom TKE)}

\textbf{Default:} Equal to \texttt{GGL90TKEmin} (no bottom enhancement).

\textbf{ECCOv4 ($10^{-6}$):} One order of magnitude higher than interior background. Accounts for:
\begin{itemize}[leftmargin=*]
  \item Rough topography and abyssal hills (unresolved at 1$^\circ$ resolution)
  \item Bottom drag from geostrophic currents interacting with seafloor
  \item Internal wave generation at topography
\end{itemize}

\textbf{Note:} Bottom TKE is applied as a Dirichlet boundary condition in the lowest wet cell. It ensures adequate near-bottom mixing in regions with weak resolved currents.

\subsubsection{mxlMaxFlag = 2 (Two-Way Sweep)}

\textbf{Default (Method 0):} Mixing length limited by distance to nearest boundary.

\textbf{ECCOv4 (Method 2):} Two-way sweep algorithm from Blanke \& Delecluse (1993). Produces the smoothest vertical profiles and is most compatible with adjoint-based optimization. See main text Section~\ref{sec:mixing_length} for detailed algorithm description.

\textbf{Advantages:}
\begin{itemize}[leftmargin=*]
  \item Prevents unphysical penetration of mixing across stable layers
  \item Produces realistic mixed layer depth and thermocline structure
  \item More differentiable (fewer discontinuities) than Method 0, improving adjoint convergence
\end{itemize}

\subsubsection{mxlSurfFlag = .TRUE. (Surface Mixing Length Enforcement)}

When \texttt{.TRUE.}, forces the mixing length between the surface and first model level to equal the layer thickness:
\[
\ell(k=1) = \Delta z_1
\]

\textbf{ECCOv4 rationale:} With coarse near-surface resolution (typically $\Delta z_1 = 10$ m), the Gaspar estimate can underestimate surface mixing length in weakly stratified conditions. Forcing $\ell = \Delta z_1$ ensures adequate surface boundary layer mixing for accurate SST simulation.

\textbf{Physical justification:} In the actively mixing surface layer, the mixing length is geometrically constrained by the layer thickness. This flag prevents unphysical $\ell < \Delta z_1$ in the surface cell.

\subsection{Minimum Alpha for Oscillation-Free Solutions}
\label{sec:alpha_min_analysis}

\subsubsection{Motivation: Why Alpha Matters for Numerical Accuracy}

While GGL90 uses fully implicit time-stepping for TKE diffusion (ensuring unconditional numerical stability), the choice of $\alpha$ (turbulent Prandtl number) has a critical impact on the \textbf{accuracy} and \textbf{smoothness} of the solution. Small values of $\alpha$ can produce cell-to-cell oscillations in TKE and mixing coefficients, even though the simulation remains stable.

The physical mechanism is:
\begin{itemize}
\item Small $\alpha$ $\Rightarrow$ small $\kappa_h = \kappa_m / \alpha$ $\Rightarrow$ weak TKE diffusion
\item Weak TKE diffusion $\Rightarrow$ sharp, localized TKE gradients
\item Sharp gradients on finite-difference grid $\Rightarrow$ under-resolved features and oscillations
\end{itemize}

This is a \textbf{resolution issue}, not a stability issue. The implicit treatment of diffusion eliminates CFL-like constraints, but cannot create information at scales finer than the grid spacing.

\subsubsection{Oscillation Diagnosis Results}

Systematic testing of the arctic convection scenario (strong cooling, ice formation, deep convection) across $\alpha = 1, 3, 5, 7.5, 10, 15, 20, 30, 50$ reveals:

\begin{table}[ht]
\centering
\caption{TKE oscillation metrics vs.\ $\alpha$ for arctic convection scenario}
\label{tab:alpha_oscillations}
\small
\begin{tabular}{ccccc}
\toprule
$\alpha$ & \textbf{TKE Roughness} & \textbf{Normalized} & \textbf{Oscillation} & \textbf{Status} \\
 & $\sum |\partial^2 \TKE/\partial z^2|$ & \textbf{Roughness} & \textbf{Count} &  \\
\midrule
1.0  & $1.22 \times 10^{-2}$ & 1.961 & 11 & \textcolor{red}{OSCILLATORY} \\
3.0  & $4.53 \times 10^{-3}$ & 0.959 & 9  & \textcolor{red}{OSCILLATORY} \\
5.0  & $4.61 \times 10^{-4}$ & 0.128 & 3  & \textcolor{orange}{THRESHOLD} \\
7.5  & $4.12 \times 10^{-4}$ & 0.112 & 3  & \textcolor{green!60!black}{SMOOTH} \\
10.0 & $3.82 \times 10^{-4}$ & 0.103 & 3  & \textcolor{green!60!black}{SMOOTH} \\
15.0 & $3.40 \times 10^{-4}$ & 0.090 & 3  & \textcolor{green!60!black}{SMOOTH} \\
20.0 & $3.09 \times 10^{-4}$ & 0.081 & 1  & \textcolor{green!60!black}{SMOOTH} \\
30.0 & $2.61 \times 10^{-4}$ & 0.068 & 1  & \textcolor{green!60!black}{SMOOTH} \\
50.0 & $2.15 \times 10^{-4}$ & 0.055 & 1  & \textcolor{green!60!black}{SMOOTH} \\
\bottomrule
\end{tabular}
\end{table}

\textbf{Key finding:} TKE roughness drops by factor of 27 from $\alpha=1$ to $\alpha=5$. The threshold for acceptable solution quality is \textbf{$\alpha \geq 5$} (roughness $< 10\%$ of maximum).

\subsubsection{Why ECCOv4 Uses $\alpha = 30$}

ECCOv4's choice of $\alpha = 30$ is well-justified:

\begin{itemize}
\item \textbf{Safety factor:} $30 / 5 = 6\times$ above the oscillation threshold
\item \textbf{Physical realism:} Observational studies suggest $\alpha \sim 20$--$50$ in stratified ocean interior
\item \textbf{Adjoint stability:} Smoother mixing fields reduce gradient noise in adjoint model
\item \textbf{Thermocline preservation:} Reduced $\kappa_h$ prevents excessive erosion of seasonal thermocline
\end{itemize}

\subsubsection{Practical Recommendations}

For new GGL90 applications:

\begin{enumerate}
\item \textbf{Minimum:} Use $\alpha \geq 5$ to avoid numerical oscillations
\item \textbf{Standard:} Use $\alpha = 10$--$20$ for most regional/process studies
\item \textbf{Global/adjoint:} Use $\alpha = 30$ (ECCOv4 value) for state estimation
\item \textbf{High stratification:} Consider $\alpha = 50$ for strongly stratified regimes
\end{enumerate}

\textbf{Grid dependence:} The minimum $\alpha$ scales with grid Reynolds number. Coarser grids (larger $\Delta z$) may tolerate smaller $\alpha$, but testing is recommended.

\textbf{Testing protocol:} When adapting GGL90 to a new configuration, plot TKE and $\kappa_m$ vertical profiles. If you observe cell-to-cell oscillations, increase $\alpha$ until profiles are smooth.

\subsection{Complete ECCOv4 data.ggl90 File}

\begin{lstlisting}[caption={Complete ECCOv4 Release 4 \texttt{data.ggl90} parameter file.},language=bash,basicstyle=\ttfamily\small]
&GGL90_PARM01
# Core physics parameters (ECCOv4 Release 4)
  GGL90ck              = 0.1,
  GGL90ceps            = 0.7,
  GGL90alpha           = 30.0,      # Reduced diffusivity for sharp thermoclines
  GGL90m2              = 3.75,

# TKE thresholds (ECCOv4 values)
  GGL90TKEmin          = 1.0e-7,    # Background internal wave mixing
  GGL90TKEsurfMin      = 1.0e-4,    # Default surface minimum
  GGL90TKEbottom       = 1.0e-6,    # Enhanced bottom mixing

# Mixing length (Method 2 for smoothness and adjoint compatibility)
  GGL90mixingLengthMin = 1.0e-8,
  mxlMaxFlag           = 2,          # Two-way sweep (Blanke & Delecluse)
  mxlSurfFlag          = .TRUE.,     # Force surface mixing length

# Caps (defaults are adequate)
  GGL90viscMax         = 1.0e2,
  GGL90diffMax         = 1.0e2,

# Boundary conditions
  GGL90_dirichlet      = .TRUE.,

# Shear computation (Method 1 for C-grid accuracy)
  calcMeanVertShear    = .TRUE.,

# I/O
  GGL90dumpFreq        = 86400.0,    # Daily output (adjust as needed)
/
\end{lstlisting}

\subsection{Performance Characteristics}

ECCOv4's GGL90 configuration has been extensively validated in global ocean simulations:

\begin{itemize}[leftmargin=*]
  \item \textbf{Mixed layer depth:} Agrees with Argo-derived climatologies to within 10--20\% globally
  \item \textbf{Thermocline sharpness:} Maintains realistic vertical temperature gradients in subtropical gyres
  \item \textbf{Deep ocean stratification:} Prevents spurious stratification buildup over multi-decadal integrations
  \item \textbf{Adjoint convergence:} Gradient checks pass to machine precision; optimization converges reliably
  \item \textbf{Computational cost:} GGL90 adds $\sim$15\% to total runtime on typical HPC systems (vs.\ constant background diffusivity)
\end{itemize}

\subsection{When to Use ECCOv4 Settings vs.\ Defaults}

\textbf{Use ECCOv4 settings if:}
\begin{itemize}[leftmargin=*]
  \item Running global ocean simulations at $\sim$1$^\circ$ resolution
  \item Performing data assimilation or adjoint-based optimization
  \item Comparing against Argo or hydrographic observations
  \item Prioritizing realistic thermocline structure over mixed layer detail
\end{itemize}

\textbf{Consider alternative settings if:}
\begin{itemize}[leftmargin=*]
  \item Running high-resolution ($<$ 10 km) process studies where $\alpha = 30$ may be too extreme
  \item Simulating specific boundary layer regimes (e.g., polar convection) where $\alpha \approx 1$ is more appropriate
  \item Coupling to sea ice or atmosphere models with different mixing assumptions
  \item Performing sensitivity studies where you want to isolate effects of individual parameters
\end{itemize}

\textbf{Recommended starting point:} Use ECCOv4 defaults as a baseline, then adjust based on comparison with observations or higher-resolution reference simulations.

% ============================================================
\section{Adjoint Model Considerations}
\label{app:adjoint_considerations}
% ============================================================

The GGL90 package has been designed with adjoint-based optimization in mind, making it suitable for gradient-based data assimilation, parameter estimation, and sensitivity studies. This appendix documents adjoint-specific features, configuration options, and best practices.

\subsection{Overview: GGL90 in Adjoint Mode}

MITgcm's automatic differentiation (AD) tool generates adjoint code from the forward model source. The adjoint model propagates sensitivities backward in time, computing gradients of a cost function (e.g., misfit to observations) with respect to control parameters (e.g., initial conditions, surface forcing, mixing parameters).

GGL90's prognostic TKE equation makes it well-suited for adjoint applications compared to diagnostic schemes like KPP:

\begin{itemize}[leftmargin=*]
  \item \textbf{Smooth dependencies:} TKE evolution is governed by smooth differential operators (advection, diffusion, source/sink terms), which produce well-behaved gradients.
  \item \textbf{Fewer branches:} Unlike KPP's many conditional branches (boundary layer depth search, profile shape functions), GGL90 has fewer \texttt{IF} statements that cause gradient discontinuities.
  \item \textbf{Simpler parameter space:} Core physics controlled by 3--4 key parameters ($c_k$, $c_\varepsilon$, $\alpha$, $m_2$) rather than KPP's $\sim$15 profile coefficients.
\end{itemize}

However, certain aspects of GGL90 require careful handling in adjoint mode, as documented below.

\subsection{Adjoint-Specific Parameters}

\subsubsection{adMxlMaxFlag: Alternative Mixing Length Method for Adjoint}

\textbf{Location:} \texttt{data.ggl90}, parameter \texttt{adMxlMaxFlag}

\textbf{Purpose:} Allows using a simpler mixing length method in adjoint mode than in forward mode.

\textbf{Default:} \texttt{UNSET\_I} (uses same method as forward, controlled by \texttt{mxlMaxFlag})

\textbf{Usage:} If \texttt{adMxlMaxFlag} is set to a valid value (0, 1, or 2), the adjoint model uses this method instead of \texttt{mxlMaxFlag}. Typical use case:
\begin{itemize}[leftmargin=*]
  \item Forward: \texttt{mxlMaxFlag = 2} (two-way sweep, most accurate)
  \item Adjoint: \texttt{adMxlMaxFlag = 0} (distance to boundaries, simpler, fewer sweep operations)
\end{itemize}

\textbf{Rationale:} Method 2's two-way sweep involves iterative vertical coupling that can amplify gradient noise in the adjoint. Method 0 is simpler and more local, improving adjoint stability in some cases.

\textbf{Recommendation:} Start with \texttt{adMxlMaxFlag = mxlMaxFlag} (default). Only change if gradient checks reveal issues or optimization convergence is poor.

\textbf{Code reference:} \texttt{ggl90\_readparms.F:945--949}, \texttt{ggl90\_mixinglength.F} (method selection logic).

\subsection{Adjoint-Specific CPP Flags}

Two compile-time flags in \texttt{GGL90\_OPTIONS.h} affect adjoint behavior:

\subsubsection{GGL90\_REGULARIZE\_MIXINGLENGTH}

\textbf{Purpose:} Provides a smoother transition as mixing length approaches its minimum value, improving gradient smoothness for adjoint calculations.

\textbf{Default:} Not defined (standard $\max(\ell, \ell_{\mathrm{min}})$ formulation)

\textbf{Effect when defined:} Replaces the sharp $\max$ function with a smooth regularization:
\[
\ell_{\mathrm{reg}} = \sqrt{\ell^2 + \ell_{\mathrm{min}}^2}
\]
instead of:
\[
\ell_{\mathrm{std}} = \max(\ell, \ell_{\mathrm{min}})
\]

\textbf{Trade-offs:}
\begin{itemize}[leftmargin=*]
  \item \textbf{Pro:} Eliminates gradient discontinuity at $\ell = \ell_{\mathrm{min}}$, improving adjoint convergence
  \item \textbf{Con:} Slightly increases mixing length in the $\ell \approx \ell_{\mathrm{min}}$ regime (typically deep ocean), affecting forward solution by $\sim$1--2\%
  \item \textbf{Pro:} Still enforces $\ell \geq \ell_{\mathrm{min}}$ asymptotically
\end{itemize}

\textbf{Recommendation:} Enable for adjoint-based optimization if gradient checks show noise near $\ell_{\mathrm{min}}$. Disable for forward-only simulations to avoid unnecessary overhead.

\textbf{Code reference:} \texttt{ggl90\_mixinglength.F:390--416}

\subsubsection{GGL90\_IDEMIX\_DIAG\_MISTAKE\_FIXED}

\textbf{Purpose:} Fixes a legacy bug in IDEMIX partial cell thickness handling. Relevant only when \texttt{useIDEMIX = .TRUE.}.

\textbf{Default:} Not defined (retains historical behavior for backward compatibility)

\textbf{Background:} Older ECCOv4 configurations used an incorrect thickness factor in IDEMIX diagnostics. Defining this flag corrects the mistake but changes output, breaking bit-for-bit reproducibility with older runs.

\textbf{Recommendation:}
\begin{itemize}[leftmargin=*]
  \item \textbf{New users:} Always define this flag to get the correct physics
  \item \textbf{Legacy ECCOv4 users:} Leave undefined to maintain consistency with previous state estimates
  \item \textbf{Adjoint impact:} Minimal (diagnostic-level change), but include for completeness if using IDEMIX
\end{itemize}

\textbf{Code reference:} \texttt{ggl90\_idemix.F}, \texttt{GGL90\_OPTIONS.h:line~35}

\subsection{GGL90 vs.\ KPP: Adjoint Complexity Comparison}

Table~\ref{tab:ggl90_kpp_adjoint} compares the adjoint characteristics of GGL90 and KPP.

\begin{table}[h!]
\centering
\caption{Adjoint complexity comparison: GGL90 vs.\ KPP.}
\label{tab:ggl90_kpp_adjoint}
\begin{tabular}{@{}lll@{}}
\toprule
\textbf{Aspect} & \textbf{GGL90} & \textbf{KPP} \\
\midrule
Primary equation type & PDE (TKE transport) & Diagnostic (BL depth search) \\
Conditional branches & Few ($\sim$5--10) & Many ($\sim$30--50) \\
Gradient smoothness & High (smooth PDEs) & Moderate (discontinuous at $h_{\mathrm{bl}}$) \\
Parameter count & 3--4 core params & $\sim$15 profile coefficients \\
Adjoint memory & Moderate (TKE field) & Low (no prognostic state) \\
Adjoint runtime & $\sim$1.5$\times$ forward & $\sim$1.3$\times$ forward \\
Gradient check pass rate & High ($>$95\%) & Moderate (70--90\%) \\
Optimization convergence & Fast (few iterations) & Slower (many iterations) \\
\bottomrule
\end{tabular}
\end{table}

\textbf{Key takeaway:} GGL90's prognostic formulation trades slightly higher memory and runtime for significantly better gradient quality and faster convergence in optimization problems.

\subsection{Smoothing Options for Adjoint Stability}

GGL90 provides optional horizontal smoothing to reduce grid-scale noise that can grow in adjoint mode.

\subsubsection{ALLOW\_GGL90\_SMOOTH (Compile-Time Flag)}

\textbf{Purpose:} Enables horizontal averaging of TKE and eddy coefficients after each timestep.

\textbf{Method:} OPA-style 5-point stencil:
\[
\TKE_{\mathrm{smooth}}(i,j) = \frac{1}{4}\left[\TKE(i-1,j) + \TKE(i+1,j) + \TKE(i,j-1) + \TKE(i,j+1)\right] + \epsilon \TKE(i,j)
\]
where $\epsilon$ is a relaxation parameter (typically 0.1--0.3).

\textbf{Use cases:}
\begin{itemize}[leftmargin=*]
  \item High-resolution simulations ($<$ 10 km) where grid-scale TKE noise appears
  \item Adjoint runs where gradient noise amplifies through the smoothing operator
\end{itemize}

\textbf{Trade-offs:}
\begin{itemize}[leftmargin=*]
  \item \textbf{Pro:} Reduces grid-scale variability, improving solution quality
  \item \textbf{Con:} Adds halo exchanges (communication cost in parallel runs)
  \item \textbf{Con:} Introduces implicit horizontal coupling (affects adjoint sensitivities)
\end{itemize}

\textbf{Recommendation:} Only enable if grid-scale noise is observed. Not used in standard ECCOv4 configuration (1$^\circ$ resolution is smooth enough).

\textbf{Code reference:} \texttt{ggl90\_calc.F:1050--1100}, \texttt{GGL90\_OPTIONS.h}

\subsection{Tuning Recommendations for Adjoint Runs}

\begin{enumerate}[leftmargin=*]
  \item \textbf{Start with forward-validated parameters:} Tune GGL90 in forward mode first to match observations. Only then enable adjoint optimization.

  \item \textbf{Gradient checking:} Run \texttt{testreport -adm} or manual finite-difference gradient checks for key cost function components. Expect agreement to 6--10 significant digits.

  \item \textbf{Parameter bounds:} Constrain optimization parameters to physical ranges:
  \begin{itemize}
    \item $0.05 < c_k < 0.2$
    \item $0.3 < c_\varepsilon < 1.5$
    \item $1 < \alpha < 100$
    \item $1 < m_2 < 10$
  \end{itemize}

  \item \textbf{Cost function design:} Avoid cost terms that depend on sharp thresholds (e.g., "mixed layer depth exceeds 100 m"). Use smooth penalties (e.g., $L^2$ norm of temperature misfit).

  \item \textbf{Mixing length method:} Start with Method 2 (\texttt{mxlMaxFlag = 2}). Only switch to Method 0 in adjoint (\texttt{adMxlMaxFlag = 0}) if gradient checks fail.

  \item \textbf{Regularization:} Consider enabling \texttt{GGL90\_REGULARIZE\_MIXINGLENGTH} if optimization stalls due to gradient noise near $\ell_{\mathrm{min}}$.

  \item \textbf{Checkpoint frequency:} GGL90's TKE field must be checkpointed for adjoint mode. Balance checkpoint frequency against memory constraints (typical: every 10--50 timesteps).
\end{enumerate}

\subsection{Known Issues and Workarounds}

\subsubsection{Issue 1: Gradient Noise in Deep Convection Regions}

\textbf{Symptom:} Adjoint gradients are noisy (large oscillations) in regions of active deep convection (e.g., Labrador Sea, Weddell Sea).

\textbf{Cause:} Rapid TKE growth during convective events creates strong nonlinearities. Adjoint sensitivities amplify through the TKE equation's implicit coupling.

\textbf{Workarounds:}
\begin{itemize}[leftmargin=*]
  \item Increase \texttt{GGL90TKEmin} to $10^{-6}$ (reduces dynamic range of TKE)
  \item Enable \texttt{GGL90\_REGULARIZE\_MIXINGLENGTH}
  \item Use \texttt{adMxlMaxFlag = 0} in convection-prone regions
  \item Exclude convective regions from cost function during initial optimization passes
\end{itemize}

\subsubsection{Issue 2: Adjoint Memory Requirements}

\textbf{Symptom:} Out-of-memory errors during adjoint compilation or runtime.

\textbf{Cause:} GGL90's TKE field (3D array) must be stored at every checkpoint for adjoint recomputation.

\textbf{Workarounds:}
\begin{itemize}[leftmargin=*]
  \item Reduce checkpoint frequency (store fewer snapshots, accept longer adjoint recomputation)
  \item Use revolving checkpoints (Griewank algorithm, available in MITgcm)
  \item Run on high-memory nodes or distributed-memory systems
\end{itemize}

\subsubsection{Issue 3: Divergence in Long Adjoint Integrations}

\textbf{Symptom:} Adjoint gradients grow exponentially or produce NaN values after $\sim$1000 timesteps.

\textbf{Cause:} Adjoint instability due to numerical issues in forward model (e.g., CFL violations, sharp gradients).

\textbf{Workarounds:}
\begin{itemize}[leftmargin=*]
  \item Verify forward model is stable and well-resolved (no CFL violations, adequate vertical resolution)
  \item Reduce timestep in regions with rapid TKE changes
  \item Apply vertical smoothing to TKE field (increase \texttt{GGL90mixingLengthMin} to regularize)
\end{itemize}

\subsection{References for Adjoint Modeling}

For further background on adjoint methods and their application to ocean models:

\begin{itemize}[leftmargin=*]
  \item Marotzke, J., R.\ Giering, K.\ Q.\ Zhang, D.\ Stammer, C.\ Hill, and T.\ Lee (1999). \textit{Construction of the adjoint MIT ocean general circulation model and application to Atlantic heat transport sensitivity}. Journal of Geophysical Research, 104(C12), 29,529--29,547.

  \item Heimbach, P., C.\ Hill, and R.\ Giering (2005). \textit{An efficient exact adjoint of the parallel MIT General Circulation Model, generated via automatic differentiation}. Future Generation Computer Systems, 21(8), 1356--1371.

  \item Forget, G., J.-M.\ Campin, P.\ Heimbach, C.\ N.\ Hill, R.\ M.\ Ponte, and C.\ Wunsch (2015). \textit{ECCO version 4: An integrated framework for non-linear inverse modeling and global ocean state estimation}. Geoscientific Model Development, 8, 3071--3104.
\end{itemize}

% ============================================================
\section{Glossary of Frequently Used Symbols}
\label{app:glossary}
% ============================================================

This glossary defines the most frequently used mathematical symbols and variable names in the GGL90 documentation and MITgcm source code.

\subsection*{Physical Variables}

\begin{longtable}{p{2.5cm}p{4cm}p{6cm}l}
\toprule
\textbf{Symbol} & \textbf{Code Variable} & \textbf{Description} & \textbf{Units} \\
\midrule
\endhead
$e$ & \verb|GGL90TKE| & Turbulent kinetic energy per unit mass & m$^2$ s$^{-2}$ \\
\addlinespace
$\ell$ & \verb|GGL90mixingLength| & Mixing length (vertical scale of turbulent eddies) & m \\
\addlinespace
$\kappa_m$ & \verb|GGL90viscArU|, \verb|GGL90viscArV| & Eddy viscosity for momentum & m$^2$ s$^{-1}$ \\
\addlinespace
$\kappa_h$ & \verb|GGL90diffKr| & Eddy diffusivity for scalars (T, S) & m$^2$ s$^{-1}$ \\
\addlinespace
$\kappa_e$ & \verb|KappaE| & Eddy diffusivity for TKE itself & m$^2$ s$^{-1}$ \\
\addlinespace
$N^2$ & \verb|Nsquare| & Squared buoyancy frequency (stratification), $N^2 = -(g/\rho_0) \partial\rho/\partial z$ & s$^{-2}$ \\
\addlinespace
$S^2$ & \verb|verticalShear| & Squared vertical shear of horizontal velocity, $S^2 = (\partial u/\partial z)^2 + (\partial v/\partial z)^2$ & s$^{-2}$ \\
\addlinespace
$Ri$ & \verb|RiNumber| & Richardson number, $Ri = N^2 / S^2$ & -- \\
\addlinespace
$Pr_t$ & \verb|TKEPrandtlNumber| & Turbulent Prandtl number, $Pr_t = \kappa_m / \kappa_h$ & -- \\
\addlinespace
$u_*$ & \verb|uStarSquare| (actually $u_*$ after sqrt) & Friction velocity at surface or bottom & m s$^{-1}$ \\
\addlinespace
$\rho$ & (not stored in GGL90) & In-situ density & kg m$^{-3}$ \\
\addlinespace
$\partial\rho/\partial z$ & \verb|sigmaR| & Vertical density gradient & kg m$^{-4}$ \\
\bottomrule
\end{longtable}

\subsection*{Dimensionless Constants}

\begin{longtable}{p{2.5cm}p{4cm}p{6cm}l}
\toprule
\textbf{Symbol} & \textbf{Code Parameter} & \textbf{Description} & \textbf{Default} \\
\midrule
\endhead
$c_\mu$ & \verb|GGL90ck| & Eddy viscosity constant & 0.1 \\
\addlinespace
$c_\varepsilon$ & \verb|GGL90ceps| & TKE dissipation constant & 0.7 \\
\addlinespace
$\alpha_e$ & \verb|GGL90alpha| & TKE diffusivity constant & 1.0 \\
\addlinespace
$B_0$ & \verb|GGL90m2| & Surface TKE boundary condition constant & 3.75 \\
\addlinespace
$e_{\text{min}}$ & \verb|GGL90TKEmin| & Minimum TKE (floor value) & 10$^{-11}$ m$^2$ s$^{-2}$ \\
\addlinespace
$\ell_{\text{min}}$ & \verb|GGL90mixingLengthMin| & Minimum mixing length (floor value) & 10$^{-8}$ m \\
\addlinespace
$\sqrt{2}$ & \verb|SQRTTWO| & Constant in Gaspar mixing length formula & 1.41421356 \\
\bottomrule
\end{longtable}

\subsection*{Grid and Coordinate Variables}

\begin{longtable}{p{2.5cm}p{4cm}p{6cm}l}
\toprule
\textbf{Symbol} & \textbf{Code Variable} & \textbf{Description} & \textbf{Units} \\
\midrule
\endhead
$k$ & \verb|k| & Vertical level index (1 to Nr) & -- \\
\addlinespace
$N_r$ & \verb|Nr| & Number of vertical levels & -- \\
\addlinespace
$\Delta z_k$ & \verb|drF(k)| & Thickness of cell $k$ at interface & m \\
\addlinespace
-- & \verb|drC(k)| & Distance between cell centers $k$ and $k+1$ & m \\
\addlinespace
$z$ & \verb|rF(k)| & Vertical position of interface $k$ (negative down in ocean) & m \\
\addlinespace
-- & \verb|coordFac| & Coordinate scaling factor (1 in z-coords, $g\rho_0$ in p-coords) & -- \\
\addlinespace
-- & \verb|maskC|, \verb|maskU|, \verb|maskV| & Land/ocean masks (0 = land, 1 = ocean) & -- \\
\bottomrule
\end{longtable}

\subsection*{Acronyms and Abbreviations}

\begin{itemize}
\item \textbf{GGL90}: Gaspar-Grégoris-Lefevre 1990 turbulence closure scheme
\item \textbf{TKE}: Turbulent Kinetic Energy
\item \textbf{KPP}: K-Profile Parameterization (alternative mixing scheme)
\item \textbf{IDEMIX}: Internal wave Dissipation and MIXing parameterization
\item \textbf{MITgcm}: Massachusetts Institute of Technology General Circulation Model
\item \textbf{ECCOv4}: Estimating the Circulation and Climate of the Ocean, version 4
\item \textbf{CPP}: C PreProcessor (compile-time flags)
\item \textbf{EOS}: Equation Of State
\item \textbf{BC}: Boundary Condition
\item \textbf{RHS}: Right Hand Side (of equation)
\item \textbf{LHS}: Left Hand Side (of equation)
\end{itemize}

% ============================================================
\section{Complete Input/Output Mapping}
\label{app:io_mapping}
% ============================================================

This appendix provides a comprehensive reference for all data exchanges between the GGL90 package and the main MITgcm model. Understanding these interfaces is essential for:
\begin{itemize}
\item Debugging GGL90 behavior in coupled simulations
\item Implementing GGL90 ports to other ocean models
\item Modifying GGL90 to use alternative input fields (e.g., from observations)
\item Diagnosing conservation properties and boundary condition handling
\end{itemize}

The GGL90 package operates as a ``mixing engine'' that receives velocity, density, and surface forcing fields from the main model, computes eddy coefficients and updates TKE, and returns the eddy coefficients to the main model for use in momentum and tracer equations. Figure~\ref{fig:ggl90_dataflow} summarizes this bidirectional exchange.

% ------------------------------------------------------------
\subsection{External Inputs from Main Model}
\label{app:external_inputs}
% ------------------------------------------------------------

Table~\ref{tab:external_inputs} lists all fields that GGL90 reads from the main model. These are \emph{inputs} from GGL90's perspective, but are computed elsewhere in MITgcm.

\begin{table}[ht]
\centering
\caption{External inputs from main MITgcm model to GGL90 package}
\label{tab:external_inputs}
\small
\begin{tabular}{p{3cm}p{3cm}p{5.5cm}lp{2.5cm}}
\toprule
\textbf{Variable Name} & \textbf{Source Module} & \textbf{Description} & \textbf{Units} & \textbf{Update Frequency} \\
\midrule
\verb|sigmaR| & \verb|CALC_PHI_HYD| (via \verb|FIND_RHO|) & Vertical gradient of in-situ density or potential density reference level, used to compute buoyancy frequency $N^2$ & kg m$^{-4}$ & Every time step \\
\addlinespace
\verb|uVel|, \verb|vVel| & Momentum solver & Horizontal velocity components on C-grid (U-points, V-points). Used to compute vertical shear $S^2 = (\partial u/\partial z)^2 + (\partial v/\partial z)^2$ & m s$^{-1}$ & Every time step \\
\addlinespace
\verb|surfaceForcingU|, \verb|surfaceForcingV| & External forcing / surface boundary layer & Surface wind stress (kinematic, i.e., divided by $\rho_0$). Used to compute friction velocity $u_*$ and surface TKE boundary condition & m$^2$ s$^{-2}$ & Every time step \\
\addlinespace
\verb|viscArNr(k)| & \verb|PARAMS.h| & Background vertical viscosity (depth-dependent array). GGL90 eddy viscosity is floored at this value & m$^2$ s$^{-1}$ & Initialization only \\
\addlinespace
\verb|diffKrNrS(k)| & \verb|PARAMS.h| & Background vertical diffusivity for scalars (depth-dependent array). GGL90 eddy diffusivity is floored at this value & m$^2$ s$^{-1}$ & Initialization only \\
\addlinespace
\verb|drF(k)|, \verb|drC(k)| & \verb|GRID.h| & Vertical grid spacing arrays: \verb|drF| = cell thickness at interfaces, \verb|drC| = distance between cell centers & m & Initialization only \\
\addlinespace
\verb|maskC|, \verb|maskU|, \verb|maskV| & \verb|GRID.h| & Land/ocean masks (0 = land, 1 = ocean) at tracer points (C), U-points, V-points & -- & Initialization only \\
\addlinespace
\verb|Ro_surf|, \verb|R_low| & \verb|GRID.h| & Vertical position of surface and bottom (z-coords: depth in m; p-coords: pressure in Pa). Used for mixing length boundary constraints & m or Pa & Initialization only \\
\addlinespace
\verb|theta|, \verb|salt| & Tracer solver & Potential temperature and salinity. \textbf{Not directly used by GGL90} except for optional diagnostics. Stratification is passed via \verb|sigmaR|, not recomputed from T/S & °C, psu & Every time step \\
\bottomrule
\end{tabular}
\end{table}

\paragraph{Key clarifications:}
\begin{itemize}
\item \textbf{Density gradient vs.\ T/S:} GGL90 does \emph{not} compute the equation of state internally. It receives the pre-computed density gradient \verb|sigmaR| from the main model's \verb|FIND_RHO| routine. This ensures consistency with the model's pressure gradient calculation.

\item \textbf{Surface forcing sign convention:} \verb|surfaceForcingU/V| are kinematic stresses (stress divided by $\rho_0$), with units of m$^2$ s$^{-2}$. Positive values indicate stress in the positive U/V direction (eastward/northward in geographic coordinates).

\item \textbf{Background mixing:} The arrays \verb|viscArNr| and \verb|diffKrNrS| are depth-dependent (varying with level index $k$) to allow specification of enhanced mixing in specific depth ranges (e.g., deep ocean background mixing from internal wave breaking).

\item \textbf{Coordinate-independent:} GGL90 works in both z-coordinates and pressure coordinates. The \verb|coordFac| scaling factor (computed internally from \verb|usingPCoords|) handles the conversion.
\end{itemize}

% ------------------------------------------------------------
\subsection{Package Outputs to Main Model}
\label{app:package_outputs}
% ------------------------------------------------------------

Table~\ref{tab:package_outputs} lists all fields that GGL90 writes for use by the main model. These are \emph{outputs} from GGL90's perspective.

\begin{table}[ht]
\centering
\caption{GGL90 package outputs to main MITgcm model}
\label{tab:package_outputs}
\small
\begin{tabular}{p{3cm}p{3.5cm}p{5cm}lp{2.5cm}}
\toprule
\textbf{Variable Name} & \textbf{Destination Module} & \textbf{Description} & \textbf{Units} & \textbf{Update Frequency} \\
\midrule
\verb|GGL90viscArU|, \verb|GGL90viscArV| & Momentum solver (via \verb|GGL90_CALC_VISC|) & Eddy viscosity at U-points and V-points, including background floor. Used in vertical viscosity term $\partial/\partial z(\kappa_m \partial u/\partial z)$ & m$^2$ s$^{-1}$ & Every time step \\
\addlinespace
\verb|GGL90diffKr| & Tracer solver (via \verb|GGL90_CALC_DIFF|) & Eddy diffusivity for all scalars (T, S, passive tracers), including background floor. Used in vertical diffusion term $\partial/\partial z(\kappa_h \partial \theta/\partial z)$ & m$^2$ s$^{-1}$ & Every time step \\
\bottomrule
\end{tabular}
\end{table}

\paragraph{Key clarifications:}
\begin{itemize}
\item \textbf{Single diffusivity for all tracers:} GGL90 computes one diffusivity field \verb|GGL90diffKr| that is applied to temperature, salinity, and all passive tracers. This assumes the turbulent Prandtl number is the same for all scalars.

\item \textbf{Already includes background:} The output arrays \verb|GGL90viscArU|, \verb|GGL90viscArV|, and \verb|GGL90diffKr| already include the background mixing coefficients (\verb|viscArNr|, \verb|diffKrNrS|). The interface routines \verb|GGL90_CALC_VISC| and \verb|GGL90_CALC_DIFF| subtract the background to avoid double-counting (see Section~\ref{sec:model_interface}).

\item \textbf{Update timing:} GGL90 is called once per time step, after the tracer and momentum equations have been advected but before the implicit vertical mixing solve. The updated eddy coefficients are used immediately in the same time step.
\end{itemize}

% ------------------------------------------------------------
\subsection{Internal Prognostic State}
\label{app:prognostic_state}
% ------------------------------------------------------------

Table~\ref{tab:prognostic_state} lists the internal prognostic variables that GGL90 carries between time steps. These are stored in the GGL90 common block and written to/read from pickup files during restart.

\begin{table}[ht]
\centering
\caption{GGL90 internal prognostic state (carried between time steps)}
\label{tab:prognostic_state}
\small
\begin{tabular}{p{3cm}p{4cm}p{5cm}l}
\toprule
\textbf{Variable Name} & \textbf{Storage Location} & \textbf{Description} & \textbf{Units} \\
\midrule
\verb|GGL90TKE| & \verb|GGL90.h| common block, \verb|pickup_ggl90.*.data| file & Turbulent kinetic energy per unit mass. \textbf{Only prognostic variable} in standard GGL90 (without IDEMIX) & m$^2$ s$^{-2}$ \\
\addlinespace
\verb|IDEMIX_E| & \verb|GGL90.h| common block (optional), \verb|pickup_ggl90.*.data| file & Internal wave energy per unit mass. Prognostic variable when \verb|ALLOW_GGL90_IDEMIX| is enabled & m$^2$ s$^{-2}$ \\
\bottomrule
\end{tabular}
\end{table}

\paragraph{Key clarifications:}
\begin{itemize}
\item \textbf{TKE is the only standard prognostic variable:} Unlike some other turbulence closures (e.g., $k$-$\varepsilon$ or $k$-$\omega$), GGL90 solves for TKE only. The mixing length $\ell$ is diagnosed algebraically from TKE and stratification (Section~\ref{sec:mxl_method2}), not evolved prognostically.

\item \textbf{Eddy coefficients are diagnostic:} \verb|GGL90viscArU|, \verb|GGL90viscArV|, and \verb|GGL90diffKr| are recomputed from TKE at every time step. They are \emph{not} prognostic variables and are \emph{not} written to pickup files.

\item \textbf{Restart behavior:} At the start of a simulation (\verb|nIter0 = 0|), TKE is initialized to \verb|GGL90TKEmin| or read from \verb|GGL90TKEFile| if specified. During a restart (\verb|nIter0 > 0|), TKE is read from the pickup file, preserving the exact state from the previous run.
\end{itemize}

% ------------------------------------------------------------
\subsection{Data Flow Diagram}
\label{app:dataflow}
% ------------------------------------------------------------

Figure~\ref{fig:ggl90_dataflow} illustrates the complete data flow between the main MITgcm model and the GGL90 package within a single time step.

\begin{figure}[ht]
\centering
\begin{tikzpicture}[
  node distance=1.5cm and 2.5cm,
  block/.style={rectangle, draw, fill=blue!10, text width=4.5cm, align=center, minimum height=1.2cm, rounded corners},
  ggl90block/.style={rectangle, draw, fill=green!15, text width=4.5cm, align=center, minimum height=1.2cm, rounded corners, thick},
  arrow/.style={->, >=stealth, thick},
  datalabel/.style={font=\small, align=left}
]

% Main model blocks
\node[block] (mom) {Main Model:\\Momentum Equations};
\node[block, below=of mom] (tracer) {Main Model:\\Tracer Equations};
\node[block, below=of tracer] (eos) {Main Model:\\Equation of State};

% GGL90 block
\node[ggl90block, right=of tracer] (ggl90) {\textbf{GGL90\_CALC}\\Compute TKE, $\ell$, $\kappa_m$, $\kappa_h$};

% Interface routines
\node[block, above right=0.8cm and 1cm of ggl90] (visc) {GGL90\_CALC\_VISC};
\node[block, below right=0.8cm and 1cm of ggl90] (diff) {GGL90\_CALC\_DIFF};

% Arrows: Inputs to GGL90
\draw[arrow] (mom.east) -- node[above, datalabel] {$u$, $v$\\(velocity)} ++(1.5cm,0) |- (ggl90.west);
\draw[arrow] (eos.east) -- node[above, datalabel] {$\partial\rho/\partial z$\\(density gradient)} ++(1.5cm,0) |- (ggl90.west);

% Surface forcing input (from above)
\draw[arrow] ($(mom.north) + (1.5cm,0)$) -- node[right, datalabel] {$\tau_x$, $\tau_y$\\(wind stress)} ($(mom.north) + (1.5cm,0) |- (ggl90.north)$) -- (ggl90.north);

% Arrows: Outputs from GGL90
\draw[arrow] (ggl90.east) -- ++(1.5cm,0) |- (visc.west);
\draw[arrow] (visc.west) -- ++(-1.5cm,0) node[above right, datalabel] {$\kappa_m^U$, $\kappa_m^V$\\(eddy viscosity)} |- (mom.east);

\draw[arrow] (ggl90.east) -- ++(1.5cm,0) |- (diff.west);
\draw[arrow] (diff.west) -- ++(-1.5cm,0) node[below right, datalabel] {$\kappa_h$\\(eddy diffusivity)} |- (tracer.east);

% TKE prognostic loop (feedback)
\draw[arrow, dashed, red] (ggl90.south) -- ++(0,-0.8cm) -| ++(-6cm,0) node[below, datalabel, pos=0.5] {$e^{n+1}$ (updated TKE)\\carried to next time step} |- ($(ggl90.west) - (0,0.5cm)$);

% Legend
\node[below=2.5cm of ggl90, anchor=north] (legend) {
  \begin{minipage}{10cm}
  \textbf{Legend:}\\
  \textcolor{blue!60}{\rule{1cm}{0.4pt}} Solid arrows: data flow within single time step\\
  \textcolor{red}{\rule{1cm}{0.4pt}} Dashed arrow: prognostic variable carried between time steps
  \end{minipage}
};

\end{tikzpicture}
\caption{Data flow diagram for GGL90 package within a single MITgcm time step. The main model provides velocity, density gradient, and surface forcing to GGL90. GGL90 computes TKE, mixing length, and eddy coefficients, then returns the eddy coefficients to the momentum and tracer solvers via interface routines. TKE is the only prognostic variable, carried between time steps (dashed red arrow).}
\label{fig:ggl90_dataflow}
\end{figure}

\paragraph{Interpretation of Figure~\ref{fig:ggl90_dataflow}:}

\begin{enumerate}
\item \textbf{Input phase (left side):} At the beginning of the vertical mixing phase, the main model has already computed:
\begin{itemize}
\item Updated velocity fields $u^{n+1/2}$, $v^{n+1/2}$ from momentum advection
\item Density gradient $\partial\rho/\partial z$ from the equation of state applied to $\theta^{n+1/2}$, $S^{n+1/2}$
\item Surface wind stress $\tau_x$, $\tau_y$ from atmospheric forcing
\end{itemize}
These fields are passed to \verb|GGL90_CALC| as inputs.

\item \textbf{GGL90 computation (center):} \verb|GGL90_CALC| performs the full TKE equation solve:
\begin{itemize}
\item Compute stratification $N^2$ from $\partial\rho/\partial z$
\item Compute vertical shear $S^2$ from $(u, v)$
\item Compute mixing length $\ell$ from TKE and $N^2$ (Sections~\ref{sec:mxl_gaspar}--\ref{sec:mxl_method2})
\item Update TKE: $e^{n+1} = e^n + \Delta t \, (\text{production} - \text{buoyancy} - \text{dissipation} + \text{diffusion})$ (Section~\ref{sec:tke_evolution})
\item Diagnose eddy coefficients: $\kappa_m = c_\mu \ell \sqrt{e^{n+1}}$, $\kappa_h = \kappa_m / Pr_t$ (Section~\ref{sec:eddy_coeffs})
\end{itemize}

\item \textbf{Output phase (right side):} The eddy coefficients are transferred to the main model:
\begin{itemize}
\item \verb|GGL90_CALC_VISC| adds $\kappa_m^U$, $\kappa_m^V$ to the momentum solver's viscosity arrays
\item \verb|GGL90_CALC_DIFF| adds $\kappa_h$ to the tracer solver's diffusivity arrays
\item The main model then applies implicit vertical mixing using these coefficients
\end{itemize}

\item \textbf{Prognostic feedback (dashed red loop):} The updated TKE $e^{n+1}$ is stored in the GGL90 common block and used as the starting point for the next time step. This is the \emph{only} variable that couples time steps; all other fields are recomputed from scratch.
\end{enumerate}

\paragraph{Frequency of data exchange:}

\begin{itemize}
\item \textbf{Every time step:} Velocity, density gradient, wind stress $\to$ GGL90; eddy coefficients $\to$ main model.
\item \textbf{Initialization only:} Grid parameters (\verb|drF|, \verb|drC|, masks), background mixing coefficients (\verb|viscArNr|, \verb|diffKrNrS|).
\item \textbf{Restart (pickup) files:} TKE field written at checkpoint intervals, read at restart.
\end{itemize}

This clean separation of concerns—where GGL90 is a self-contained ``black box'' that receives state variables and returns mixing coefficients—makes it straightforward to replace GGL90 with alternative turbulence closures (e.g., KPP, Mellor-Yamada) without modifying the main model code.

% ============================================================
\end{document}
